% sec:geo-timelike — Timelike Geodesics
\section{Timelike Geodesics}
\label{sec:geo-timelike}

Massive particles trace out timelike geodesics --- the trajectories
that determine accretion disc structure, binary star motion, and the
orbital mechanics of infalling matter.  Where does the inner edge
of an accretion disc sit, and why?  The answer lies in the stability
of circular orbits, which reduces to an analysis of the effective
potential for massive particles.  The Hamiltonian formulation
developed in \cref{sec:geo-formulations} applies with a single
change: the constraint level shifts from zero to $-\tfrac{1}{2}$,
and the arbitrary affine parameter of null geodesics is replaced by
the particle's proper time.

\paragraph{The timelike constraint.}

A timelike geodesic has a tangent vector with negative norm:
$\metric\,\dot{x}^\mu\dot{x}^\nu = -1$, where the overdot
denotes $\d/\d\tau$ and the normalisation follows from the
$(-,+,+,+)$ signature convention.  In Hamiltonian language, this
becomes
%
\begin{equation}
  \mathcal{H} = \tfrac{1}{2}\,\inversemetric\,p_\mu\,p_\nu
  = -\tfrac{1}{2} \,.
  \label{eq:timelike-constraint}
\end{equation}
%
The constraint preservation proof of \cref{eq:H-conservation}
applies verbatim: $\d\mathcal{H}/\d\lambda = 0$ for any autonomous
Hamiltonian, so $\mathcal{H} = -\tfrac{1}{2}$ is maintained
throughout the trajectory.  As a numerical diagnostic, the departure
of $\mathcal{H}$ from $-\tfrac{1}{2}$ measures accumulated
integration error, exactly as the departure from zero does for null
rays (\cref{sec:geo-null}).

Unlike null geodesics, where the affine parameter has a rescaling
freedom (\cref{eq:affine-rescaling}), the normalisation
$\metric\,\dot{x}^\mu\dot{x}^\nu = -1$ fixes the future-directed
parametrisation uniquely (up to an additive constant): the affine
parameter is proper time $\tau$, the time measured by a clock
co-moving with the particle.  Every conserved quantity built from
$p_\mu$ is therefore a \emph{specific} quantity --- energy per unit
rest mass, angular momentum per unit rest mass --- as noted in
\cref{sec:schw-orbits}.
\physnote{The rest frame of a massive particle is the frame in
which the particle is momentarily at rest.  No such frame exists
for a photon, which is why null geodesics have no canonical
parametrisation.}

\paragraph{The effective potential from the Hamiltonian.}

The conserved quantities from \cref{sec:geo-constants} reduce the
eight-dimensional phase space to a single radial equation.  To see
how, expand the constraint \cref{eq:timelike-constraint}
in the equatorial Schwarzschild spacetime.  With $f = 1 - 2M/r$,
the inverse metric is diagonal: $g^{tt} = -1/f$,
$g^{rr} = f$, $g^{\varphi\varphi} = 1/r^2$.  Substituting
$p_t = -E$ and $p_\varphi = L$ from
\cref{eq:geo-energy,eq:geo-angmom} gives
%
\begin{equation}
  \tfrac{1}{2}\bigl[-E^2/f + f\,p_r^2 + L^2/r^2\bigr]
  = -\tfrac{1}{2} \,.
  \label{eq:timelike-H-expanded}
\end{equation}
%
The position equation \cref{eq:geo-xdot} relates $p_r$ to the
coordinate velocity: $\dot{r} = g^{rr}\,p_r = f\,p_r$, so
$p_r = \dot{r}/f$ and hence $f\,p_r^2 = \dot{r}^2/f$.
Substituting and rearranging yields
%
\begin{equation}
  \dot{r}^2 = E^2
  - \underbrace{\left(1 - \frac{2M}{r}\right)
    \!\left(1 + \frac{L^2}{r^2}\right)}_{V_{\text{eff}}(r)} \,,
  \label{eq:timelike-radial}
\end{equation}
%
recovering the effective potential \cref{eq:veff-timelike} derived
in \cref{sec:schw-orbits} by expanding the metric norm
constraint directly.  The Hamiltonian route makes the reduction
transparent: each conserved quantity eliminates one momentum
component ($E$ removes $p_t$, $L$ removes $p_\varphi$, equatorial
symmetry removes $p_\theta$), and the constraint
$\mathcal{H} = -\tfrac{1}{2}$ eliminates $p_r$ in favour of
$\dot{r}$, leaving a single first-order ODE for $r(\tau)$.
\xrefnote{For the Kerr geometry, the Carter constant
$\mathcal{Q}$ (\cref{eq:carter-constant}) provides a fourth
conservation law, reducing the system to two separated first-order
equations --- one for $r(\tau)$ and one for $\theta(\tau)$ ---
each involving only its own coordinate and the four constants
$E$, $L$, $\mathcal{Q}$, $\mathcal{H}$.  The specific form of
the Kerr radial equation is \cref{eq:kerr-radial-R}.}

\paragraph{Circular orbits and stability.}

The shape of $V_{\text{eff}}(r)$ determines the orbit type.
Circular orbits sit at extrema of the effective potential, where two
conditions hold simultaneously:
%
\begin{equation}
  V_{\text{eff}}(r_c) = E^2 \,, \qquad
  V'_{\text{eff}}(r_c) = 0 \,.
  \label{eq:geo-circular-conditions}
\end{equation}
%
The first condition ensures $\dot{r} = 0$ --- the particle stays at
constant radius.  The second ensures the equilibrium is maintained:
differentiating the radial equation \cref{eq:timelike-radial} with
respect to $r$ gives $2\ddot{r} = -V'_{\text{eff}}$, so
$V'_{\text{eff}} = 0$ is equivalent to zero radial acceleration.
Stability depends on the curvature of the potential at $r_c$:
%
\begin{itemize}
\item $V''_{\text{eff}}(r_c) > 0$ (local minimum): the orbit is
  \emph{stable}.  Small perturbations cause radial oscillations
  about $r_c$, and the particle traces a precessing ellipse.
\item $V''_{\text{eff}}(r_c) < 0$ (local maximum): the orbit is
  \emph{unstable}.  Any perturbation sends the particle spiralling
  inward or escaping outward.
\item $V''_{\text{eff}}(r_c) = 0$ (with the circular-orbit
  conditions still holding): the minimum and maximum have merged
  --- this is the \emph{innermost stable circular orbit} (ISCO),
  the boundary between the two regimes.
\end{itemize}
%
The three simultaneous conditions $V_{\text{eff}} = E^2$,
$V'_{\text{eff}} = 0$, $V''_{\text{eff}} = 0$ are precisely those
solved in \cref{sec:schw-orbits} to obtain
$r_{\text{ISCO}} = 6M$ for Schwarzschild and in
\cref{sec:kerr-isco} to obtain the spin-dependent ISCO formula
\cref{eq:kerr-isco} for Kerr.

More generally, three classes of timelike orbit emerge from the
effective potential landscape:
%
\begin{enumerate}
\item \textbf{Bound orbits}: $E^2$ lies below the local
  maximum of $V_{\text{eff}}$, and the particle oscillates between
  two turning points $r_{\text{min}}$ and $r_{\text{max}}$.  In
  the equatorial plane, the trajectory traces a rosette pattern ---
  an ellipse whose perihelion advances with each orbit.  The
  relativistic perihelion advance, famously $43''$ per century for
  Mercury, is a direct consequence of the $-2ML^2/r^3$ correction
  in $V_{\text{eff}}$ (\cref{eq:veff-timelike}); for orbits near
  a black hole the precession per orbit is much larger.
\item \textbf{Marginal orbits}: $E^2$ equals the local
  maximum of $V_{\text{eff}}$, and the particle asymptotically
  approaches the unstable circular orbit, winding an infinite
  number of times without reaching it.
\item \textbf{Plunge orbits}: $E^2$ exceeds the peak of
  $V_{\text{eff}}$, or the angular momentum vanishes (radial
  infall), leaving no centrifugal barrier.  The particle falls
  through the horizon without a turning point.
\end{enumerate}

\begin{figure}[htbp]
  \centering
  %% Creator: Matplotlib, PGF backend
%%
%% To include the figure in your LaTeX document, write
%%   \input{<filename>.pgf}
%%
%% Make sure the required packages are loaded in your preamble
%%   \usepackage{pgf}
%%
%% Also ensure that all the required font packages are loaded; for instance,
%% the lmodern package is sometimes necessary when using math font.
%%   \usepackage{lmodern}
%%
%% Figures using additional raster images can only be included by \input if
%% they are in the same directory as the main LaTeX file. For loading figures
%% from other directories you can use the `import` package
%%   \usepackage{import}
%%
%% and then include the figures with
%%   \import{<path to file>}{<filename>.pgf}
%%
%% Matplotlib used the following preamble
%%   \def\mathdefault#1{#1}
%%   \everymath=\expandafter{\the\everymath\displaystyle}
%%   \IfFileExists{scrextend.sty}{
%%     \usepackage[fontsize=10.000000pt]{scrextend}
%%   }{
%%     \renewcommand{\normalsize}{\fontsize{10.000000}{12.000000}\selectfont}
%%     \normalsize
%%   }
%%   \usepackage{fontspec}
%%   \usepackage{unicode-math}
%%   \setmainfont{STIX Two Text}[Numbers=OldStyle, Ligatures=TeX]
%%   \setmathfont{STIX Two Math}
%%   \setmonofont{Source Code Pro}[Scale=MatchLowercase]
%%   \ifdefined\pdftexversion\else  % non-pdftex case.
%%     \usepackage{fontspec}
%%     \setmainfont{DejaVuSerif.ttf}[Path=\detokenize{/Users/gvn/.pyenv/versions/3.12.11/lib/python3.12/site-packages/matplotlib/mpl-data/fonts/ttf/}]
%%     \setsansfont{DejaVuSans.ttf}[Path=\detokenize{/Users/gvn/.pyenv/versions/3.12.11/lib/python3.12/site-packages/matplotlib/mpl-data/fonts/ttf/}]
%%     \setmonofont{DejaVuSansMono.ttf}[Path=\detokenize{/Users/gvn/.pyenv/versions/3.12.11/lib/python3.12/site-packages/matplotlib/mpl-data/fonts/ttf/}]
%%   \fi
%%   \makeatletter\@ifpackageloaded{underscore}{}{\usepackage[strings]{underscore}}\makeatother
%%
\begingroup%
\makeatletter%
\begin{pgfpicture}%
\pgfpathrectangle{\pgfpointorigin}{\pgfqpoint{3.844235in}{3.724603in}}%
\pgfusepath{use as bounding box, clip}%
\begin{pgfscope}%
\pgfsetbuttcap%
\pgfsetmiterjoin%
\definecolor{currentfill}{rgb}{1.000000,1.000000,1.000000}%
\pgfsetfillcolor{currentfill}%
\pgfsetlinewidth{0.000000pt}%
\definecolor{currentstroke}{rgb}{1.000000,1.000000,1.000000}%
\pgfsetstrokecolor{currentstroke}%
\pgfsetdash{}{0pt}%
\pgfpathmoveto{\pgfqpoint{0.000000in}{-0.000000in}}%
\pgfpathlineto{\pgfqpoint{3.844235in}{-0.000000in}}%
\pgfpathlineto{\pgfqpoint{3.844235in}{3.724603in}}%
\pgfpathlineto{\pgfqpoint{0.000000in}{3.724603in}}%
\pgfpathlineto{\pgfqpoint{0.000000in}{-0.000000in}}%
\pgfpathclose%
\pgfusepath{fill}%
\end{pgfscope}%
\begin{pgfscope}%
\pgfsetbuttcap%
\pgfsetmiterjoin%
\definecolor{currentfill}{rgb}{1.000000,1.000000,1.000000}%
\pgfsetfillcolor{currentfill}%
\pgfsetlinewidth{0.000000pt}%
\definecolor{currentstroke}{rgb}{0.000000,0.000000,0.000000}%
\pgfsetstrokecolor{currentstroke}%
\pgfsetstrokeopacity{0.000000}%
\pgfsetdash{}{0pt}%
\pgfpathmoveto{\pgfqpoint{0.472301in}{0.352669in}}%
\pgfpathlineto{\pgfqpoint{3.824235in}{0.352669in}}%
\pgfpathlineto{\pgfqpoint{3.824235in}{3.704603in}}%
\pgfpathlineto{\pgfqpoint{0.472301in}{3.704603in}}%
\pgfpathlineto{\pgfqpoint{0.472301in}{0.352669in}}%
\pgfpathclose%
\pgfusepath{fill}%
\end{pgfscope}%
\begin{pgfscope}%
\pgfpathrectangle{\pgfqpoint{0.472301in}{0.352669in}}{\pgfqpoint{3.351934in}{3.351934in}}%
\pgfusepath{clip}%
\pgfsetbuttcap%
\pgfsetmiterjoin%
\definecolor{currentfill}{rgb}{0.000000,0.000000,0.000000}%
\pgfsetfillcolor{currentfill}%
\pgfsetfillopacity{0.120000}%
\pgfsetlinewidth{1.003750pt}%
\definecolor{currentstroke}{rgb}{0.000000,0.000000,0.000000}%
\pgfsetstrokecolor{currentstroke}%
\pgfsetstrokeopacity{0.120000}%
\pgfsetdash{}{0pt}%
\pgfpathmoveto{\pgfqpoint{2.300628in}{2.028636in}}%
\pgfpathlineto{\pgfqpoint{2.300595in}{2.031837in}}%
\pgfpathlineto{\pgfqpoint{2.300494in}{2.035038in}}%
\pgfpathlineto{\pgfqpoint{2.300326in}{2.038235in}}%
\pgfpathlineto{\pgfqpoint{2.300091in}{2.041428in}}%
\pgfpathlineto{\pgfqpoint{2.299788in}{2.044615in}}%
\pgfpathlineto{\pgfqpoint{2.299419in}{2.047795in}}%
\pgfpathlineto{\pgfqpoint{2.298983in}{2.050967in}}%
\pgfpathlineto{\pgfqpoint{2.298481in}{2.054129in}}%
\pgfpathlineto{\pgfqpoint{2.297912in}{2.057280in}}%
\pgfpathlineto{\pgfqpoint{2.297277in}{2.060418in}}%
\pgfpathlineto{\pgfqpoint{2.296576in}{2.063542in}}%
\pgfpathlineto{\pgfqpoint{2.295810in}{2.066651in}}%
\pgfpathlineto{\pgfqpoint{2.294978in}{2.069742in}}%
\pgfpathlineto{\pgfqpoint{2.294082in}{2.072816in}}%
\pgfpathlineto{\pgfqpoint{2.293122in}{2.075870in}}%
\pgfpathlineto{\pgfqpoint{2.292097in}{2.078904in}}%
\pgfpathlineto{\pgfqpoint{2.291009in}{2.081915in}}%
\pgfpathlineto{\pgfqpoint{2.289858in}{2.084902in}}%
\pgfpathlineto{\pgfqpoint{2.288645in}{2.087865in}}%
\pgfpathlineto{\pgfqpoint{2.287369in}{2.090802in}}%
\pgfpathlineto{\pgfqpoint{2.286032in}{2.093711in}}%
\pgfpathlineto{\pgfqpoint{2.284634in}{2.096591in}}%
\pgfpathlineto{\pgfqpoint{2.283176in}{2.099441in}}%
\pgfpathlineto{\pgfqpoint{2.281659in}{2.102261in}}%
\pgfpathlineto{\pgfqpoint{2.280082in}{2.105047in}}%
\pgfpathlineto{\pgfqpoint{2.278448in}{2.107800in}}%
\pgfpathlineto{\pgfqpoint{2.276756in}{2.110518in}}%
\pgfpathlineto{\pgfqpoint{2.275007in}{2.113200in}}%
\pgfpathlineto{\pgfqpoint{2.273202in}{2.115844in}}%
\pgfpathlineto{\pgfqpoint{2.271342in}{2.118450in}}%
\pgfpathlineto{\pgfqpoint{2.269427in}{2.121016in}}%
\pgfpathlineto{\pgfqpoint{2.267459in}{2.123542in}}%
\pgfpathlineto{\pgfqpoint{2.265439in}{2.126025in}}%
\pgfpathlineto{\pgfqpoint{2.263367in}{2.128466in}}%
\pgfpathlineto{\pgfqpoint{2.261244in}{2.130862in}}%
\pgfpathlineto{\pgfqpoint{2.259071in}{2.133214in}}%
\pgfpathlineto{\pgfqpoint{2.256849in}{2.135519in}}%
\pgfpathlineto{\pgfqpoint{2.254579in}{2.137777in}}%
\pgfpathlineto{\pgfqpoint{2.252262in}{2.139987in}}%
\pgfpathlineto{\pgfqpoint{2.249899in}{2.142147in}}%
\pgfpathlineto{\pgfqpoint{2.247492in}{2.144258in}}%
\pgfpathlineto{\pgfqpoint{2.245040in}{2.146317in}}%
\pgfpathlineto{\pgfqpoint{2.242546in}{2.148325in}}%
\pgfpathlineto{\pgfqpoint{2.240010in}{2.150279in}}%
\pgfpathlineto{\pgfqpoint{2.237434in}{2.152180in}}%
\pgfpathlineto{\pgfqpoint{2.234818in}{2.154026in}}%
\pgfpathlineto{\pgfqpoint{2.232165in}{2.155817in}}%
\pgfpathlineto{\pgfqpoint{2.229474in}{2.157552in}}%
\pgfpathlineto{\pgfqpoint{2.226747in}{2.159230in}}%
\pgfpathlineto{\pgfqpoint{2.223986in}{2.160850in}}%
\pgfpathlineto{\pgfqpoint{2.221191in}{2.162412in}}%
\pgfpathlineto{\pgfqpoint{2.218364in}{2.163915in}}%
\pgfpathlineto{\pgfqpoint{2.215506in}{2.165358in}}%
\pgfpathlineto{\pgfqpoint{2.212618in}{2.166740in}}%
\pgfpathlineto{\pgfqpoint{2.209702in}{2.168062in}}%
\pgfpathlineto{\pgfqpoint{2.206759in}{2.169322in}}%
\pgfpathlineto{\pgfqpoint{2.203789in}{2.170520in}}%
\pgfpathlineto{\pgfqpoint{2.200796in}{2.171655in}}%
\pgfpathlineto{\pgfqpoint{2.197779in}{2.172728in}}%
\pgfpathlineto{\pgfqpoint{2.194740in}{2.173736in}}%
\pgfpathlineto{\pgfqpoint{2.191681in}{2.174681in}}%
\pgfpathlineto{\pgfqpoint{2.188603in}{2.175561in}}%
\pgfpathlineto{\pgfqpoint{2.185507in}{2.176376in}}%
\pgfpathlineto{\pgfqpoint{2.182394in}{2.177126in}}%
\pgfpathlineto{\pgfqpoint{2.179266in}{2.177810in}}%
\pgfpathlineto{\pgfqpoint{2.176125in}{2.178428in}}%
\pgfpathlineto{\pgfqpoint{2.172971in}{2.178981in}}%
\pgfpathlineto{\pgfqpoint{2.169807in}{2.179466in}}%
\pgfpathlineto{\pgfqpoint{2.166633in}{2.179886in}}%
\pgfpathlineto{\pgfqpoint{2.163451in}{2.180238in}}%
\pgfpathlineto{\pgfqpoint{2.160262in}{2.180524in}}%
\pgfpathlineto{\pgfqpoint{2.157068in}{2.180742in}}%
\pgfpathlineto{\pgfqpoint{2.153870in}{2.180894in}}%
\pgfpathlineto{\pgfqpoint{2.150669in}{2.180978in}}%
\pgfpathlineto{\pgfqpoint{2.147467in}{2.180995in}}%
\pgfpathlineto{\pgfqpoint{2.144266in}{2.180944in}}%
\pgfpathlineto{\pgfqpoint{2.141067in}{2.180826in}}%
\pgfpathlineto{\pgfqpoint{2.137870in}{2.180641in}}%
\pgfpathlineto{\pgfqpoint{2.134679in}{2.180389in}}%
\pgfpathlineto{\pgfqpoint{2.131493in}{2.180070in}}%
\pgfpathlineto{\pgfqpoint{2.128315in}{2.179684in}}%
\pgfpathlineto{\pgfqpoint{2.125145in}{2.179232in}}%
\pgfpathlineto{\pgfqpoint{2.121986in}{2.178713in}}%
\pgfpathlineto{\pgfqpoint{2.118838in}{2.178127in}}%
\pgfpathlineto{\pgfqpoint{2.115703in}{2.177476in}}%
\pgfpathlineto{\pgfqpoint{2.112583in}{2.176759in}}%
\pgfpathlineto{\pgfqpoint{2.109479in}{2.175976in}}%
\pgfpathlineto{\pgfqpoint{2.106391in}{2.175129in}}%
\pgfpathlineto{\pgfqpoint{2.103322in}{2.174216in}}%
\pgfpathlineto{\pgfqpoint{2.100273in}{2.173240in}}%
\pgfpathlineto{\pgfqpoint{2.097245in}{2.172199in}}%
\pgfpathlineto{\pgfqpoint{2.094240in}{2.171096in}}%
\pgfpathlineto{\pgfqpoint{2.091258in}{2.169929in}}%
\pgfpathlineto{\pgfqpoint{2.088302in}{2.168700in}}%
\pgfpathlineto{\pgfqpoint{2.085372in}{2.167409in}}%
\pgfpathlineto{\pgfqpoint{2.082470in}{2.166057in}}%
\pgfpathlineto{\pgfqpoint{2.079597in}{2.164644in}}%
\pgfpathlineto{\pgfqpoint{2.076754in}{2.163171in}}%
\pgfpathlineto{\pgfqpoint{2.073943in}{2.161638in}}%
\pgfpathlineto{\pgfqpoint{2.071165in}{2.160047in}}%
\pgfpathlineto{\pgfqpoint{2.068421in}{2.158398in}}%
\pgfpathlineto{\pgfqpoint{2.065712in}{2.156692in}}%
\pgfpathlineto{\pgfqpoint{2.063039in}{2.154929in}}%
\pgfpathlineto{\pgfqpoint{2.060404in}{2.153110in}}%
\pgfpathlineto{\pgfqpoint{2.057808in}{2.151236in}}%
\pgfpathlineto{\pgfqpoint{2.055252in}{2.149308in}}%
\pgfpathlineto{\pgfqpoint{2.052737in}{2.147327in}}%
\pgfpathlineto{\pgfqpoint{2.050264in}{2.145294in}}%
\pgfpathlineto{\pgfqpoint{2.047834in}{2.143209in}}%
\pgfpathlineto{\pgfqpoint{2.045449in}{2.141073in}}%
\pgfpathlineto{\pgfqpoint{2.043109in}{2.138888in}}%
\pgfpathlineto{\pgfqpoint{2.040816in}{2.136654in}}%
\pgfpathlineto{\pgfqpoint{2.038570in}{2.134372in}}%
\pgfpathlineto{\pgfqpoint{2.036372in}{2.132044in}}%
\pgfpathlineto{\pgfqpoint{2.034224in}{2.129670in}}%
\pgfpathlineto{\pgfqpoint{2.032126in}{2.127251in}}%
\pgfpathlineto{\pgfqpoint{2.030080in}{2.124789in}}%
\pgfpathlineto{\pgfqpoint{2.028086in}{2.122284in}}%
\pgfpathlineto{\pgfqpoint{2.026144in}{2.119738in}}%
\pgfpathlineto{\pgfqpoint{2.024257in}{2.117152in}}%
\pgfpathlineto{\pgfqpoint{2.022424in}{2.114527in}}%
\pgfpathlineto{\pgfqpoint{2.020647in}{2.111864in}}%
\pgfpathlineto{\pgfqpoint{2.018927in}{2.109164in}}%
\pgfpathlineto{\pgfqpoint{2.017263in}{2.106428in}}%
\pgfpathlineto{\pgfqpoint{2.015658in}{2.103658in}}%
\pgfpathlineto{\pgfqpoint{2.014111in}{2.100855in}}%
\pgfpathlineto{\pgfqpoint{2.012623in}{2.098020in}}%
\pgfpathlineto{\pgfqpoint{2.011195in}{2.095155in}}%
\pgfpathlineto{\pgfqpoint{2.009827in}{2.092260in}}%
\pgfpathlineto{\pgfqpoint{2.008521in}{2.089337in}}%
\pgfpathlineto{\pgfqpoint{2.007276in}{2.086387in}}%
\pgfpathlineto{\pgfqpoint{2.006094in}{2.083411in}}%
\pgfpathlineto{\pgfqpoint{2.004974in}{2.080412in}}%
\pgfpathlineto{\pgfqpoint{2.003918in}{2.077390in}}%
\pgfpathlineto{\pgfqpoint{2.002925in}{2.074346in}}%
\pgfpathlineto{\pgfqpoint{2.001997in}{2.071282in}}%
\pgfpathlineto{\pgfqpoint{2.001133in}{2.068199in}}%
\pgfpathlineto{\pgfqpoint{2.000334in}{2.065098in}}%
\pgfpathlineto{\pgfqpoint{1.999601in}{2.061982in}}%
\pgfpathlineto{\pgfqpoint{1.998933in}{2.058851in}}%
\pgfpathlineto{\pgfqpoint{1.998331in}{2.055706in}}%
\pgfpathlineto{\pgfqpoint{1.997795in}{2.052549in}}%
\pgfpathlineto{\pgfqpoint{1.997326in}{2.049382in}}%
\pgfpathlineto{\pgfqpoint{1.996924in}{2.046206in}}%
\pgfpathlineto{\pgfqpoint{1.996588in}{2.043022in}}%
\pgfpathlineto{\pgfqpoint{1.996319in}{2.039832in}}%
\pgfpathlineto{\pgfqpoint{1.996117in}{2.036637in}}%
\pgfpathlineto{\pgfqpoint{1.995983in}{2.033438in}}%
\pgfpathlineto{\pgfqpoint{1.995916in}{2.030237in}}%
\pgfpathlineto{\pgfqpoint{1.995916in}{2.027035in}}%
\pgfpathlineto{\pgfqpoint{1.995983in}{2.023834in}}%
\pgfpathlineto{\pgfqpoint{1.996117in}{2.020635in}}%
\pgfpathlineto{\pgfqpoint{1.996319in}{2.017440in}}%
\pgfpathlineto{\pgfqpoint{1.996588in}{2.014250in}}%
\pgfpathlineto{\pgfqpoint{1.996924in}{2.011066in}}%
\pgfpathlineto{\pgfqpoint{1.997326in}{2.007890in}}%
\pgfpathlineto{\pgfqpoint{1.997795in}{2.004722in}}%
\pgfpathlineto{\pgfqpoint{1.998331in}{2.001566in}}%
\pgfpathlineto{\pgfqpoint{1.998933in}{1.998421in}}%
\pgfpathlineto{\pgfqpoint{1.999601in}{1.995290in}}%
\pgfpathlineto{\pgfqpoint{2.000334in}{1.992174in}}%
\pgfpathlineto{\pgfqpoint{2.001133in}{1.989073in}}%
\pgfpathlineto{\pgfqpoint{2.001997in}{1.985990in}}%
\pgfpathlineto{\pgfqpoint{2.002925in}{1.982926in}}%
\pgfpathlineto{\pgfqpoint{2.003918in}{1.979882in}}%
\pgfpathlineto{\pgfqpoint{2.004974in}{1.976860in}}%
\pgfpathlineto{\pgfqpoint{2.006094in}{1.973860in}}%
\pgfpathlineto{\pgfqpoint{2.007276in}{1.970885in}}%
\pgfpathlineto{\pgfqpoint{2.008521in}{1.967935in}}%
\pgfpathlineto{\pgfqpoint{2.009827in}{1.965012in}}%
\pgfpathlineto{\pgfqpoint{2.011195in}{1.962117in}}%
\pgfpathlineto{\pgfqpoint{2.012623in}{1.959252in}}%
\pgfpathlineto{\pgfqpoint{2.014111in}{1.956417in}}%
\pgfpathlineto{\pgfqpoint{2.015658in}{1.953614in}}%
\pgfpathlineto{\pgfqpoint{2.017263in}{1.950844in}}%
\pgfpathlineto{\pgfqpoint{2.018927in}{1.948108in}}%
\pgfpathlineto{\pgfqpoint{2.020647in}{1.945408in}}%
\pgfpathlineto{\pgfqpoint{2.022424in}{1.942745in}}%
\pgfpathlineto{\pgfqpoint{2.024257in}{1.940120in}}%
\pgfpathlineto{\pgfqpoint{2.026144in}{1.937534in}}%
\pgfpathlineto{\pgfqpoint{2.028086in}{1.934988in}}%
\pgfpathlineto{\pgfqpoint{2.030080in}{1.932483in}}%
\pgfpathlineto{\pgfqpoint{2.032126in}{1.930021in}}%
\pgfpathlineto{\pgfqpoint{2.034224in}{1.927602in}}%
\pgfpathlineto{\pgfqpoint{2.036372in}{1.925228in}}%
\pgfpathlineto{\pgfqpoint{2.038570in}{1.922900in}}%
\pgfpathlineto{\pgfqpoint{2.040816in}{1.920618in}}%
\pgfpathlineto{\pgfqpoint{2.043109in}{1.918384in}}%
\pgfpathlineto{\pgfqpoint{2.045449in}{1.916199in}}%
\pgfpathlineto{\pgfqpoint{2.047834in}{1.914063in}}%
\pgfpathlineto{\pgfqpoint{2.050264in}{1.911978in}}%
\pgfpathlineto{\pgfqpoint{2.052737in}{1.909945in}}%
\pgfpathlineto{\pgfqpoint{2.055252in}{1.907963in}}%
\pgfpathlineto{\pgfqpoint{2.057808in}{1.906036in}}%
\pgfpathlineto{\pgfqpoint{2.060404in}{1.904162in}}%
\pgfpathlineto{\pgfqpoint{2.063039in}{1.902343in}}%
\pgfpathlineto{\pgfqpoint{2.065712in}{1.900580in}}%
\pgfpathlineto{\pgfqpoint{2.068421in}{1.898874in}}%
\pgfpathlineto{\pgfqpoint{2.071165in}{1.897225in}}%
\pgfpathlineto{\pgfqpoint{2.073943in}{1.895634in}}%
\pgfpathlineto{\pgfqpoint{2.076754in}{1.894101in}}%
\pgfpathlineto{\pgfqpoint{2.079597in}{1.892628in}}%
\pgfpathlineto{\pgfqpoint{2.082470in}{1.891215in}}%
\pgfpathlineto{\pgfqpoint{2.085372in}{1.889863in}}%
\pgfpathlineto{\pgfqpoint{2.088302in}{1.888572in}}%
\pgfpathlineto{\pgfqpoint{2.091258in}{1.887343in}}%
\pgfpathlineto{\pgfqpoint{2.094240in}{1.886176in}}%
\pgfpathlineto{\pgfqpoint{2.097245in}{1.885072in}}%
\pgfpathlineto{\pgfqpoint{2.100273in}{1.884032in}}%
\pgfpathlineto{\pgfqpoint{2.103322in}{1.883056in}}%
\pgfpathlineto{\pgfqpoint{2.106391in}{1.882143in}}%
\pgfpathlineto{\pgfqpoint{2.109479in}{1.881296in}}%
\pgfpathlineto{\pgfqpoint{2.112583in}{1.880513in}}%
\pgfpathlineto{\pgfqpoint{2.115703in}{1.879796in}}%
\pgfpathlineto{\pgfqpoint{2.118838in}{1.879145in}}%
\pgfpathlineto{\pgfqpoint{2.121986in}{1.878559in}}%
\pgfpathlineto{\pgfqpoint{2.125145in}{1.878040in}}%
\pgfpathlineto{\pgfqpoint{2.128315in}{1.877588in}}%
\pgfpathlineto{\pgfqpoint{2.131493in}{1.877202in}}%
\pgfpathlineto{\pgfqpoint{2.134679in}{1.876883in}}%
\pgfpathlineto{\pgfqpoint{2.137870in}{1.876631in}}%
\pgfpathlineto{\pgfqpoint{2.141067in}{1.876446in}}%
\pgfpathlineto{\pgfqpoint{2.144266in}{1.876328in}}%
\pgfpathlineto{\pgfqpoint{2.147467in}{1.876277in}}%
\pgfpathlineto{\pgfqpoint{2.150669in}{1.876294in}}%
\pgfpathlineto{\pgfqpoint{2.153870in}{1.876378in}}%
\pgfpathlineto{\pgfqpoint{2.157068in}{1.876530in}}%
\pgfpathlineto{\pgfqpoint{2.160262in}{1.876748in}}%
\pgfpathlineto{\pgfqpoint{2.163451in}{1.877034in}}%
\pgfpathlineto{\pgfqpoint{2.166633in}{1.877386in}}%
\pgfpathlineto{\pgfqpoint{2.169807in}{1.877805in}}%
\pgfpathlineto{\pgfqpoint{2.172971in}{1.878291in}}%
\pgfpathlineto{\pgfqpoint{2.176125in}{1.878844in}}%
\pgfpathlineto{\pgfqpoint{2.179266in}{1.879462in}}%
\pgfpathlineto{\pgfqpoint{2.182394in}{1.880146in}}%
\pgfpathlineto{\pgfqpoint{2.185507in}{1.880896in}}%
\pgfpathlineto{\pgfqpoint{2.188603in}{1.881711in}}%
\pgfpathlineto{\pgfqpoint{2.191681in}{1.882591in}}%
\pgfpathlineto{\pgfqpoint{2.194740in}{1.883536in}}%
\pgfpathlineto{\pgfqpoint{2.197779in}{1.884544in}}%
\pgfpathlineto{\pgfqpoint{2.200796in}{1.885616in}}%
\pgfpathlineto{\pgfqpoint{2.203789in}{1.886752in}}%
\pgfpathlineto{\pgfqpoint{2.206759in}{1.887950in}}%
\pgfpathlineto{\pgfqpoint{2.209702in}{1.889210in}}%
\pgfpathlineto{\pgfqpoint{2.212618in}{1.890532in}}%
\pgfpathlineto{\pgfqpoint{2.215506in}{1.891914in}}%
\pgfpathlineto{\pgfqpoint{2.218364in}{1.893357in}}%
\pgfpathlineto{\pgfqpoint{2.221191in}{1.894860in}}%
\pgfpathlineto{\pgfqpoint{2.223986in}{1.896422in}}%
\pgfpathlineto{\pgfqpoint{2.226747in}{1.898042in}}%
\pgfpathlineto{\pgfqpoint{2.229474in}{1.899720in}}%
\pgfpathlineto{\pgfqpoint{2.232165in}{1.901455in}}%
\pgfpathlineto{\pgfqpoint{2.234818in}{1.903246in}}%
\pgfpathlineto{\pgfqpoint{2.237434in}{1.905092in}}%
\pgfpathlineto{\pgfqpoint{2.240010in}{1.906993in}}%
\pgfpathlineto{\pgfqpoint{2.242546in}{1.908947in}}%
\pgfpathlineto{\pgfqpoint{2.245040in}{1.910955in}}%
\pgfpathlineto{\pgfqpoint{2.247492in}{1.913014in}}%
\pgfpathlineto{\pgfqpoint{2.249899in}{1.915125in}}%
\pgfpathlineto{\pgfqpoint{2.252262in}{1.917285in}}%
\pgfpathlineto{\pgfqpoint{2.254579in}{1.919495in}}%
\pgfpathlineto{\pgfqpoint{2.256849in}{1.921753in}}%
\pgfpathlineto{\pgfqpoint{2.259071in}{1.924058in}}%
\pgfpathlineto{\pgfqpoint{2.261244in}{1.926410in}}%
\pgfpathlineto{\pgfqpoint{2.263367in}{1.928806in}}%
\pgfpathlineto{\pgfqpoint{2.265439in}{1.931247in}}%
\pgfpathlineto{\pgfqpoint{2.267459in}{1.933730in}}%
\pgfpathlineto{\pgfqpoint{2.269427in}{1.936256in}}%
\pgfpathlineto{\pgfqpoint{2.271342in}{1.938822in}}%
\pgfpathlineto{\pgfqpoint{2.273202in}{1.941428in}}%
\pgfpathlineto{\pgfqpoint{2.275007in}{1.944072in}}%
\pgfpathlineto{\pgfqpoint{2.276756in}{1.946754in}}%
\pgfpathlineto{\pgfqpoint{2.278448in}{1.949472in}}%
\pgfpathlineto{\pgfqpoint{2.280082in}{1.952225in}}%
\pgfpathlineto{\pgfqpoint{2.281659in}{1.955011in}}%
\pgfpathlineto{\pgfqpoint{2.283176in}{1.957830in}}%
\pgfpathlineto{\pgfqpoint{2.284634in}{1.960681in}}%
\pgfpathlineto{\pgfqpoint{2.286032in}{1.963561in}}%
\pgfpathlineto{\pgfqpoint{2.287369in}{1.966470in}}%
\pgfpathlineto{\pgfqpoint{2.288645in}{1.969407in}}%
\pgfpathlineto{\pgfqpoint{2.289858in}{1.972370in}}%
\pgfpathlineto{\pgfqpoint{2.291009in}{1.975357in}}%
\pgfpathlineto{\pgfqpoint{2.292097in}{1.978368in}}%
\pgfpathlineto{\pgfqpoint{2.293122in}{1.981402in}}%
\pgfpathlineto{\pgfqpoint{2.294082in}{1.984456in}}%
\pgfpathlineto{\pgfqpoint{2.294978in}{1.987530in}}%
\pgfpathlineto{\pgfqpoint{2.295810in}{1.990621in}}%
\pgfpathlineto{\pgfqpoint{2.296576in}{1.993730in}}%
\pgfpathlineto{\pgfqpoint{2.297277in}{1.996854in}}%
\pgfpathlineto{\pgfqpoint{2.297912in}{1.999992in}}%
\pgfpathlineto{\pgfqpoint{2.298481in}{2.003143in}}%
\pgfpathlineto{\pgfqpoint{2.298983in}{2.006305in}}%
\pgfpathlineto{\pgfqpoint{2.299419in}{2.009477in}}%
\pgfpathlineto{\pgfqpoint{2.299788in}{2.012657in}}%
\pgfpathlineto{\pgfqpoint{2.300091in}{2.015844in}}%
\pgfpathlineto{\pgfqpoint{2.300326in}{2.019037in}}%
\pgfpathlineto{\pgfqpoint{2.300494in}{2.022234in}}%
\pgfpathlineto{\pgfqpoint{2.300595in}{2.025435in}}%
\pgfpathlineto{\pgfqpoint{2.300628in}{2.028636in}}%
\pgfpathlineto{\pgfqpoint{2.300628in}{2.028636in}}%
\pgfpathclose%
\pgfusepath{stroke,fill}%
\end{pgfscope}%
\begin{pgfscope}%
\pgfsetbuttcap%
\pgfsetroundjoin%
\definecolor{currentfill}{rgb}{0.000000,0.000000,0.000000}%
\pgfsetfillcolor{currentfill}%
\pgfsetlinewidth{0.501875pt}%
\definecolor{currentstroke}{rgb}{0.000000,0.000000,0.000000}%
\pgfsetstrokecolor{currentstroke}%
\pgfsetdash{}{0pt}%
\pgfsys@defobject{currentmarker}{\pgfqpoint{0.000000in}{0.000000in}}{\pgfqpoint{0.000000in}{0.041667in}}{%
\pgfpathmoveto{\pgfqpoint{0.000000in}{0.000000in}}%
\pgfpathlineto{\pgfqpoint{0.000000in}{0.041667in}}%
\pgfusepath{stroke,fill}%
}%
\begin{pgfscope}%
\pgfsys@transformshift{0.624661in}{0.352669in}%
\pgfsys@useobject{currentmarker}{}%
\end{pgfscope}%
\end{pgfscope}%
\begin{pgfscope}%
\definecolor{textcolor}{rgb}{0.000000,0.000000,0.000000}%
\pgfsetstrokecolor{textcolor}%
\pgfsetfillcolor{textcolor}%
\pgftext[x=0.624661in,y=0.304058in,,top]{\color{textcolor}{\sffamily\fontsize{8.000000}{9.600000}\selectfont\catcode`\^=\active\def^{\ifmmode\sp\else\^{}\fi}\catcode`\%=\active\def%{\%}\ensuremath{-}20}}%
\end{pgfscope}%
\begin{pgfscope}%
\pgfsetbuttcap%
\pgfsetroundjoin%
\definecolor{currentfill}{rgb}{0.000000,0.000000,0.000000}%
\pgfsetfillcolor{currentfill}%
\pgfsetlinewidth{0.501875pt}%
\definecolor{currentstroke}{rgb}{0.000000,0.000000,0.000000}%
\pgfsetstrokecolor{currentstroke}%
\pgfsetdash{}{0pt}%
\pgfsys@defobject{currentmarker}{\pgfqpoint{0.000000in}{0.000000in}}{\pgfqpoint{0.000000in}{0.041667in}}{%
\pgfpathmoveto{\pgfqpoint{0.000000in}{0.000000in}}%
\pgfpathlineto{\pgfqpoint{0.000000in}{0.041667in}}%
\pgfusepath{stroke,fill}%
}%
\begin{pgfscope}%
\pgfsys@transformshift{1.005563in}{0.352669in}%
\pgfsys@useobject{currentmarker}{}%
\end{pgfscope}%
\end{pgfscope}%
\begin{pgfscope}%
\definecolor{textcolor}{rgb}{0.000000,0.000000,0.000000}%
\pgfsetstrokecolor{textcolor}%
\pgfsetfillcolor{textcolor}%
\pgftext[x=1.005563in,y=0.304058in,,top]{\color{textcolor}{\sffamily\fontsize{8.000000}{9.600000}\selectfont\catcode`\^=\active\def^{\ifmmode\sp\else\^{}\fi}\catcode`\%=\active\def%{\%}\ensuremath{-}15}}%
\end{pgfscope}%
\begin{pgfscope}%
\pgfsetbuttcap%
\pgfsetroundjoin%
\definecolor{currentfill}{rgb}{0.000000,0.000000,0.000000}%
\pgfsetfillcolor{currentfill}%
\pgfsetlinewidth{0.501875pt}%
\definecolor{currentstroke}{rgb}{0.000000,0.000000,0.000000}%
\pgfsetstrokecolor{currentstroke}%
\pgfsetdash{}{0pt}%
\pgfsys@defobject{currentmarker}{\pgfqpoint{0.000000in}{0.000000in}}{\pgfqpoint{0.000000in}{0.041667in}}{%
\pgfpathmoveto{\pgfqpoint{0.000000in}{0.000000in}}%
\pgfpathlineto{\pgfqpoint{0.000000in}{0.041667in}}%
\pgfusepath{stroke,fill}%
}%
\begin{pgfscope}%
\pgfsys@transformshift{1.386464in}{0.352669in}%
\pgfsys@useobject{currentmarker}{}%
\end{pgfscope}%
\end{pgfscope}%
\begin{pgfscope}%
\definecolor{textcolor}{rgb}{0.000000,0.000000,0.000000}%
\pgfsetstrokecolor{textcolor}%
\pgfsetfillcolor{textcolor}%
\pgftext[x=1.386464in,y=0.304058in,,top]{\color{textcolor}{\sffamily\fontsize{8.000000}{9.600000}\selectfont\catcode`\^=\active\def^{\ifmmode\sp\else\^{}\fi}\catcode`\%=\active\def%{\%}\ensuremath{-}10}}%
\end{pgfscope}%
\begin{pgfscope}%
\pgfsetbuttcap%
\pgfsetroundjoin%
\definecolor{currentfill}{rgb}{0.000000,0.000000,0.000000}%
\pgfsetfillcolor{currentfill}%
\pgfsetlinewidth{0.501875pt}%
\definecolor{currentstroke}{rgb}{0.000000,0.000000,0.000000}%
\pgfsetstrokecolor{currentstroke}%
\pgfsetdash{}{0pt}%
\pgfsys@defobject{currentmarker}{\pgfqpoint{0.000000in}{0.000000in}}{\pgfqpoint{0.000000in}{0.041667in}}{%
\pgfpathmoveto{\pgfqpoint{0.000000in}{0.000000in}}%
\pgfpathlineto{\pgfqpoint{0.000000in}{0.041667in}}%
\pgfusepath{stroke,fill}%
}%
\begin{pgfscope}%
\pgfsys@transformshift{1.767366in}{0.352669in}%
\pgfsys@useobject{currentmarker}{}%
\end{pgfscope}%
\end{pgfscope}%
\begin{pgfscope}%
\definecolor{textcolor}{rgb}{0.000000,0.000000,0.000000}%
\pgfsetstrokecolor{textcolor}%
\pgfsetfillcolor{textcolor}%
\pgftext[x=1.767366in,y=0.304058in,,top]{\color{textcolor}{\sffamily\fontsize{8.000000}{9.600000}\selectfont\catcode`\^=\active\def^{\ifmmode\sp\else\^{}\fi}\catcode`\%=\active\def%{\%}\ensuremath{-}5}}%
\end{pgfscope}%
\begin{pgfscope}%
\pgfsetbuttcap%
\pgfsetroundjoin%
\definecolor{currentfill}{rgb}{0.000000,0.000000,0.000000}%
\pgfsetfillcolor{currentfill}%
\pgfsetlinewidth{0.501875pt}%
\definecolor{currentstroke}{rgb}{0.000000,0.000000,0.000000}%
\pgfsetstrokecolor{currentstroke}%
\pgfsetdash{}{0pt}%
\pgfsys@defobject{currentmarker}{\pgfqpoint{0.000000in}{0.000000in}}{\pgfqpoint{0.000000in}{0.041667in}}{%
\pgfpathmoveto{\pgfqpoint{0.000000in}{0.000000in}}%
\pgfpathlineto{\pgfqpoint{0.000000in}{0.041667in}}%
\pgfusepath{stroke,fill}%
}%
\begin{pgfscope}%
\pgfsys@transformshift{2.148268in}{0.352669in}%
\pgfsys@useobject{currentmarker}{}%
\end{pgfscope}%
\end{pgfscope}%
\begin{pgfscope}%
\definecolor{textcolor}{rgb}{0.000000,0.000000,0.000000}%
\pgfsetstrokecolor{textcolor}%
\pgfsetfillcolor{textcolor}%
\pgftext[x=2.148268in,y=0.304058in,,top]{\color{textcolor}{\sffamily\fontsize{8.000000}{9.600000}\selectfont\catcode`\^=\active\def^{\ifmmode\sp\else\^{}\fi}\catcode`\%=\active\def%{\%}0}}%
\end{pgfscope}%
\begin{pgfscope}%
\pgfsetbuttcap%
\pgfsetroundjoin%
\definecolor{currentfill}{rgb}{0.000000,0.000000,0.000000}%
\pgfsetfillcolor{currentfill}%
\pgfsetlinewidth{0.501875pt}%
\definecolor{currentstroke}{rgb}{0.000000,0.000000,0.000000}%
\pgfsetstrokecolor{currentstroke}%
\pgfsetdash{}{0pt}%
\pgfsys@defobject{currentmarker}{\pgfqpoint{0.000000in}{0.000000in}}{\pgfqpoint{0.000000in}{0.041667in}}{%
\pgfpathmoveto{\pgfqpoint{0.000000in}{0.000000in}}%
\pgfpathlineto{\pgfqpoint{0.000000in}{0.041667in}}%
\pgfusepath{stroke,fill}%
}%
\begin{pgfscope}%
\pgfsys@transformshift{2.529169in}{0.352669in}%
\pgfsys@useobject{currentmarker}{}%
\end{pgfscope}%
\end{pgfscope}%
\begin{pgfscope}%
\definecolor{textcolor}{rgb}{0.000000,0.000000,0.000000}%
\pgfsetstrokecolor{textcolor}%
\pgfsetfillcolor{textcolor}%
\pgftext[x=2.529169in,y=0.304058in,,top]{\color{textcolor}{\sffamily\fontsize{8.000000}{9.600000}\selectfont\catcode`\^=\active\def^{\ifmmode\sp\else\^{}\fi}\catcode`\%=\active\def%{\%}5}}%
\end{pgfscope}%
\begin{pgfscope}%
\pgfsetbuttcap%
\pgfsetroundjoin%
\definecolor{currentfill}{rgb}{0.000000,0.000000,0.000000}%
\pgfsetfillcolor{currentfill}%
\pgfsetlinewidth{0.501875pt}%
\definecolor{currentstroke}{rgb}{0.000000,0.000000,0.000000}%
\pgfsetstrokecolor{currentstroke}%
\pgfsetdash{}{0pt}%
\pgfsys@defobject{currentmarker}{\pgfqpoint{0.000000in}{0.000000in}}{\pgfqpoint{0.000000in}{0.041667in}}{%
\pgfpathmoveto{\pgfqpoint{0.000000in}{0.000000in}}%
\pgfpathlineto{\pgfqpoint{0.000000in}{0.041667in}}%
\pgfusepath{stroke,fill}%
}%
\begin{pgfscope}%
\pgfsys@transformshift{2.910071in}{0.352669in}%
\pgfsys@useobject{currentmarker}{}%
\end{pgfscope}%
\end{pgfscope}%
\begin{pgfscope}%
\definecolor{textcolor}{rgb}{0.000000,0.000000,0.000000}%
\pgfsetstrokecolor{textcolor}%
\pgfsetfillcolor{textcolor}%
\pgftext[x=2.910071in,y=0.304058in,,top]{\color{textcolor}{\sffamily\fontsize{8.000000}{9.600000}\selectfont\catcode`\^=\active\def^{\ifmmode\sp\else\^{}\fi}\catcode`\%=\active\def%{\%}10}}%
\end{pgfscope}%
\begin{pgfscope}%
\pgfsetbuttcap%
\pgfsetroundjoin%
\definecolor{currentfill}{rgb}{0.000000,0.000000,0.000000}%
\pgfsetfillcolor{currentfill}%
\pgfsetlinewidth{0.501875pt}%
\definecolor{currentstroke}{rgb}{0.000000,0.000000,0.000000}%
\pgfsetstrokecolor{currentstroke}%
\pgfsetdash{}{0pt}%
\pgfsys@defobject{currentmarker}{\pgfqpoint{0.000000in}{0.000000in}}{\pgfqpoint{0.000000in}{0.041667in}}{%
\pgfpathmoveto{\pgfqpoint{0.000000in}{0.000000in}}%
\pgfpathlineto{\pgfqpoint{0.000000in}{0.041667in}}%
\pgfusepath{stroke,fill}%
}%
\begin{pgfscope}%
\pgfsys@transformshift{3.290973in}{0.352669in}%
\pgfsys@useobject{currentmarker}{}%
\end{pgfscope}%
\end{pgfscope}%
\begin{pgfscope}%
\definecolor{textcolor}{rgb}{0.000000,0.000000,0.000000}%
\pgfsetstrokecolor{textcolor}%
\pgfsetfillcolor{textcolor}%
\pgftext[x=3.290973in,y=0.304058in,,top]{\color{textcolor}{\sffamily\fontsize{8.000000}{9.600000}\selectfont\catcode`\^=\active\def^{\ifmmode\sp\else\^{}\fi}\catcode`\%=\active\def%{\%}15}}%
\end{pgfscope}%
\begin{pgfscope}%
\pgfsetbuttcap%
\pgfsetroundjoin%
\definecolor{currentfill}{rgb}{0.000000,0.000000,0.000000}%
\pgfsetfillcolor{currentfill}%
\pgfsetlinewidth{0.501875pt}%
\definecolor{currentstroke}{rgb}{0.000000,0.000000,0.000000}%
\pgfsetstrokecolor{currentstroke}%
\pgfsetdash{}{0pt}%
\pgfsys@defobject{currentmarker}{\pgfqpoint{0.000000in}{0.000000in}}{\pgfqpoint{0.000000in}{0.041667in}}{%
\pgfpathmoveto{\pgfqpoint{0.000000in}{0.000000in}}%
\pgfpathlineto{\pgfqpoint{0.000000in}{0.041667in}}%
\pgfusepath{stroke,fill}%
}%
\begin{pgfscope}%
\pgfsys@transformshift{3.671874in}{0.352669in}%
\pgfsys@useobject{currentmarker}{}%
\end{pgfscope}%
\end{pgfscope}%
\begin{pgfscope}%
\definecolor{textcolor}{rgb}{0.000000,0.000000,0.000000}%
\pgfsetstrokecolor{textcolor}%
\pgfsetfillcolor{textcolor}%
\pgftext[x=3.671874in,y=0.304058in,,top]{\color{textcolor}{\sffamily\fontsize{8.000000}{9.600000}\selectfont\catcode`\^=\active\def^{\ifmmode\sp\else\^{}\fi}\catcode`\%=\active\def%{\%}20}}%
\end{pgfscope}%
\begin{pgfscope}%
\pgfsetbuttcap%
\pgfsetroundjoin%
\definecolor{currentfill}{rgb}{0.000000,0.000000,0.000000}%
\pgfsetfillcolor{currentfill}%
\pgfsetlinewidth{0.301125pt}%
\definecolor{currentstroke}{rgb}{0.000000,0.000000,0.000000}%
\pgfsetstrokecolor{currentstroke}%
\pgfsetdash{}{0pt}%
\pgfsys@defobject{currentmarker}{\pgfqpoint{0.000000in}{0.000000in}}{\pgfqpoint{0.000000in}{0.027778in}}{%
\pgfpathmoveto{\pgfqpoint{0.000000in}{0.000000in}}%
\pgfpathlineto{\pgfqpoint{0.000000in}{0.027778in}}%
\pgfusepath{stroke,fill}%
}%
\begin{pgfscope}%
\pgfsys@transformshift{0.472301in}{0.352669in}%
\pgfsys@useobject{currentmarker}{}%
\end{pgfscope}%
\end{pgfscope}%
\begin{pgfscope}%
\pgfsetbuttcap%
\pgfsetroundjoin%
\definecolor{currentfill}{rgb}{0.000000,0.000000,0.000000}%
\pgfsetfillcolor{currentfill}%
\pgfsetlinewidth{0.301125pt}%
\definecolor{currentstroke}{rgb}{0.000000,0.000000,0.000000}%
\pgfsetstrokecolor{currentstroke}%
\pgfsetdash{}{0pt}%
\pgfsys@defobject{currentmarker}{\pgfqpoint{0.000000in}{0.000000in}}{\pgfqpoint{0.000000in}{0.027778in}}{%
\pgfpathmoveto{\pgfqpoint{0.000000in}{0.000000in}}%
\pgfpathlineto{\pgfqpoint{0.000000in}{0.027778in}}%
\pgfusepath{stroke,fill}%
}%
\begin{pgfscope}%
\pgfsys@transformshift{0.548481in}{0.352669in}%
\pgfsys@useobject{currentmarker}{}%
\end{pgfscope}%
\end{pgfscope}%
\begin{pgfscope}%
\pgfsetbuttcap%
\pgfsetroundjoin%
\definecolor{currentfill}{rgb}{0.000000,0.000000,0.000000}%
\pgfsetfillcolor{currentfill}%
\pgfsetlinewidth{0.301125pt}%
\definecolor{currentstroke}{rgb}{0.000000,0.000000,0.000000}%
\pgfsetstrokecolor{currentstroke}%
\pgfsetdash{}{0pt}%
\pgfsys@defobject{currentmarker}{\pgfqpoint{0.000000in}{0.000000in}}{\pgfqpoint{0.000000in}{0.027778in}}{%
\pgfpathmoveto{\pgfqpoint{0.000000in}{0.000000in}}%
\pgfpathlineto{\pgfqpoint{0.000000in}{0.027778in}}%
\pgfusepath{stroke,fill}%
}%
\begin{pgfscope}%
\pgfsys@transformshift{0.700842in}{0.352669in}%
\pgfsys@useobject{currentmarker}{}%
\end{pgfscope}%
\end{pgfscope}%
\begin{pgfscope}%
\pgfsetbuttcap%
\pgfsetroundjoin%
\definecolor{currentfill}{rgb}{0.000000,0.000000,0.000000}%
\pgfsetfillcolor{currentfill}%
\pgfsetlinewidth{0.301125pt}%
\definecolor{currentstroke}{rgb}{0.000000,0.000000,0.000000}%
\pgfsetstrokecolor{currentstroke}%
\pgfsetdash{}{0pt}%
\pgfsys@defobject{currentmarker}{\pgfqpoint{0.000000in}{0.000000in}}{\pgfqpoint{0.000000in}{0.027778in}}{%
\pgfpathmoveto{\pgfqpoint{0.000000in}{0.000000in}}%
\pgfpathlineto{\pgfqpoint{0.000000in}{0.027778in}}%
\pgfusepath{stroke,fill}%
}%
\begin{pgfscope}%
\pgfsys@transformshift{0.777022in}{0.352669in}%
\pgfsys@useobject{currentmarker}{}%
\end{pgfscope}%
\end{pgfscope}%
\begin{pgfscope}%
\pgfsetbuttcap%
\pgfsetroundjoin%
\definecolor{currentfill}{rgb}{0.000000,0.000000,0.000000}%
\pgfsetfillcolor{currentfill}%
\pgfsetlinewidth{0.301125pt}%
\definecolor{currentstroke}{rgb}{0.000000,0.000000,0.000000}%
\pgfsetstrokecolor{currentstroke}%
\pgfsetdash{}{0pt}%
\pgfsys@defobject{currentmarker}{\pgfqpoint{0.000000in}{0.000000in}}{\pgfqpoint{0.000000in}{0.027778in}}{%
\pgfpathmoveto{\pgfqpoint{0.000000in}{0.000000in}}%
\pgfpathlineto{\pgfqpoint{0.000000in}{0.027778in}}%
\pgfusepath{stroke,fill}%
}%
\begin{pgfscope}%
\pgfsys@transformshift{0.853202in}{0.352669in}%
\pgfsys@useobject{currentmarker}{}%
\end{pgfscope}%
\end{pgfscope}%
\begin{pgfscope}%
\pgfsetbuttcap%
\pgfsetroundjoin%
\definecolor{currentfill}{rgb}{0.000000,0.000000,0.000000}%
\pgfsetfillcolor{currentfill}%
\pgfsetlinewidth{0.301125pt}%
\definecolor{currentstroke}{rgb}{0.000000,0.000000,0.000000}%
\pgfsetstrokecolor{currentstroke}%
\pgfsetdash{}{0pt}%
\pgfsys@defobject{currentmarker}{\pgfqpoint{0.000000in}{0.000000in}}{\pgfqpoint{0.000000in}{0.027778in}}{%
\pgfpathmoveto{\pgfqpoint{0.000000in}{0.000000in}}%
\pgfpathlineto{\pgfqpoint{0.000000in}{0.027778in}}%
\pgfusepath{stroke,fill}%
}%
\begin{pgfscope}%
\pgfsys@transformshift{0.929383in}{0.352669in}%
\pgfsys@useobject{currentmarker}{}%
\end{pgfscope}%
\end{pgfscope}%
\begin{pgfscope}%
\pgfsetbuttcap%
\pgfsetroundjoin%
\definecolor{currentfill}{rgb}{0.000000,0.000000,0.000000}%
\pgfsetfillcolor{currentfill}%
\pgfsetlinewidth{0.301125pt}%
\definecolor{currentstroke}{rgb}{0.000000,0.000000,0.000000}%
\pgfsetstrokecolor{currentstroke}%
\pgfsetdash{}{0pt}%
\pgfsys@defobject{currentmarker}{\pgfqpoint{0.000000in}{0.000000in}}{\pgfqpoint{0.000000in}{0.027778in}}{%
\pgfpathmoveto{\pgfqpoint{0.000000in}{0.000000in}}%
\pgfpathlineto{\pgfqpoint{0.000000in}{0.027778in}}%
\pgfusepath{stroke,fill}%
}%
\begin{pgfscope}%
\pgfsys@transformshift{1.081743in}{0.352669in}%
\pgfsys@useobject{currentmarker}{}%
\end{pgfscope}%
\end{pgfscope}%
\begin{pgfscope}%
\pgfsetbuttcap%
\pgfsetroundjoin%
\definecolor{currentfill}{rgb}{0.000000,0.000000,0.000000}%
\pgfsetfillcolor{currentfill}%
\pgfsetlinewidth{0.301125pt}%
\definecolor{currentstroke}{rgb}{0.000000,0.000000,0.000000}%
\pgfsetstrokecolor{currentstroke}%
\pgfsetdash{}{0pt}%
\pgfsys@defobject{currentmarker}{\pgfqpoint{0.000000in}{0.000000in}}{\pgfqpoint{0.000000in}{0.027778in}}{%
\pgfpathmoveto{\pgfqpoint{0.000000in}{0.000000in}}%
\pgfpathlineto{\pgfqpoint{0.000000in}{0.027778in}}%
\pgfusepath{stroke,fill}%
}%
\begin{pgfscope}%
\pgfsys@transformshift{1.157924in}{0.352669in}%
\pgfsys@useobject{currentmarker}{}%
\end{pgfscope}%
\end{pgfscope}%
\begin{pgfscope}%
\pgfsetbuttcap%
\pgfsetroundjoin%
\definecolor{currentfill}{rgb}{0.000000,0.000000,0.000000}%
\pgfsetfillcolor{currentfill}%
\pgfsetlinewidth{0.301125pt}%
\definecolor{currentstroke}{rgb}{0.000000,0.000000,0.000000}%
\pgfsetstrokecolor{currentstroke}%
\pgfsetdash{}{0pt}%
\pgfsys@defobject{currentmarker}{\pgfqpoint{0.000000in}{0.000000in}}{\pgfqpoint{0.000000in}{0.027778in}}{%
\pgfpathmoveto{\pgfqpoint{0.000000in}{0.000000in}}%
\pgfpathlineto{\pgfqpoint{0.000000in}{0.027778in}}%
\pgfusepath{stroke,fill}%
}%
\begin{pgfscope}%
\pgfsys@transformshift{1.234104in}{0.352669in}%
\pgfsys@useobject{currentmarker}{}%
\end{pgfscope}%
\end{pgfscope}%
\begin{pgfscope}%
\pgfsetbuttcap%
\pgfsetroundjoin%
\definecolor{currentfill}{rgb}{0.000000,0.000000,0.000000}%
\pgfsetfillcolor{currentfill}%
\pgfsetlinewidth{0.301125pt}%
\definecolor{currentstroke}{rgb}{0.000000,0.000000,0.000000}%
\pgfsetstrokecolor{currentstroke}%
\pgfsetdash{}{0pt}%
\pgfsys@defobject{currentmarker}{\pgfqpoint{0.000000in}{0.000000in}}{\pgfqpoint{0.000000in}{0.027778in}}{%
\pgfpathmoveto{\pgfqpoint{0.000000in}{0.000000in}}%
\pgfpathlineto{\pgfqpoint{0.000000in}{0.027778in}}%
\pgfusepath{stroke,fill}%
}%
\begin{pgfscope}%
\pgfsys@transformshift{1.310284in}{0.352669in}%
\pgfsys@useobject{currentmarker}{}%
\end{pgfscope}%
\end{pgfscope}%
\begin{pgfscope}%
\pgfsetbuttcap%
\pgfsetroundjoin%
\definecolor{currentfill}{rgb}{0.000000,0.000000,0.000000}%
\pgfsetfillcolor{currentfill}%
\pgfsetlinewidth{0.301125pt}%
\definecolor{currentstroke}{rgb}{0.000000,0.000000,0.000000}%
\pgfsetstrokecolor{currentstroke}%
\pgfsetdash{}{0pt}%
\pgfsys@defobject{currentmarker}{\pgfqpoint{0.000000in}{0.000000in}}{\pgfqpoint{0.000000in}{0.027778in}}{%
\pgfpathmoveto{\pgfqpoint{0.000000in}{0.000000in}}%
\pgfpathlineto{\pgfqpoint{0.000000in}{0.027778in}}%
\pgfusepath{stroke,fill}%
}%
\begin{pgfscope}%
\pgfsys@transformshift{1.462645in}{0.352669in}%
\pgfsys@useobject{currentmarker}{}%
\end{pgfscope}%
\end{pgfscope}%
\begin{pgfscope}%
\pgfsetbuttcap%
\pgfsetroundjoin%
\definecolor{currentfill}{rgb}{0.000000,0.000000,0.000000}%
\pgfsetfillcolor{currentfill}%
\pgfsetlinewidth{0.301125pt}%
\definecolor{currentstroke}{rgb}{0.000000,0.000000,0.000000}%
\pgfsetstrokecolor{currentstroke}%
\pgfsetdash{}{0pt}%
\pgfsys@defobject{currentmarker}{\pgfqpoint{0.000000in}{0.000000in}}{\pgfqpoint{0.000000in}{0.027778in}}{%
\pgfpathmoveto{\pgfqpoint{0.000000in}{0.000000in}}%
\pgfpathlineto{\pgfqpoint{0.000000in}{0.027778in}}%
\pgfusepath{stroke,fill}%
}%
\begin{pgfscope}%
\pgfsys@transformshift{1.538825in}{0.352669in}%
\pgfsys@useobject{currentmarker}{}%
\end{pgfscope}%
\end{pgfscope}%
\begin{pgfscope}%
\pgfsetbuttcap%
\pgfsetroundjoin%
\definecolor{currentfill}{rgb}{0.000000,0.000000,0.000000}%
\pgfsetfillcolor{currentfill}%
\pgfsetlinewidth{0.301125pt}%
\definecolor{currentstroke}{rgb}{0.000000,0.000000,0.000000}%
\pgfsetstrokecolor{currentstroke}%
\pgfsetdash{}{0pt}%
\pgfsys@defobject{currentmarker}{\pgfqpoint{0.000000in}{0.000000in}}{\pgfqpoint{0.000000in}{0.027778in}}{%
\pgfpathmoveto{\pgfqpoint{0.000000in}{0.000000in}}%
\pgfpathlineto{\pgfqpoint{0.000000in}{0.027778in}}%
\pgfusepath{stroke,fill}%
}%
\begin{pgfscope}%
\pgfsys@transformshift{1.615005in}{0.352669in}%
\pgfsys@useobject{currentmarker}{}%
\end{pgfscope}%
\end{pgfscope}%
\begin{pgfscope}%
\pgfsetbuttcap%
\pgfsetroundjoin%
\definecolor{currentfill}{rgb}{0.000000,0.000000,0.000000}%
\pgfsetfillcolor{currentfill}%
\pgfsetlinewidth{0.301125pt}%
\definecolor{currentstroke}{rgb}{0.000000,0.000000,0.000000}%
\pgfsetstrokecolor{currentstroke}%
\pgfsetdash{}{0pt}%
\pgfsys@defobject{currentmarker}{\pgfqpoint{0.000000in}{0.000000in}}{\pgfqpoint{0.000000in}{0.027778in}}{%
\pgfpathmoveto{\pgfqpoint{0.000000in}{0.000000in}}%
\pgfpathlineto{\pgfqpoint{0.000000in}{0.027778in}}%
\pgfusepath{stroke,fill}%
}%
\begin{pgfscope}%
\pgfsys@transformshift{1.691186in}{0.352669in}%
\pgfsys@useobject{currentmarker}{}%
\end{pgfscope}%
\end{pgfscope}%
\begin{pgfscope}%
\pgfsetbuttcap%
\pgfsetroundjoin%
\definecolor{currentfill}{rgb}{0.000000,0.000000,0.000000}%
\pgfsetfillcolor{currentfill}%
\pgfsetlinewidth{0.301125pt}%
\definecolor{currentstroke}{rgb}{0.000000,0.000000,0.000000}%
\pgfsetstrokecolor{currentstroke}%
\pgfsetdash{}{0pt}%
\pgfsys@defobject{currentmarker}{\pgfqpoint{0.000000in}{0.000000in}}{\pgfqpoint{0.000000in}{0.027778in}}{%
\pgfpathmoveto{\pgfqpoint{0.000000in}{0.000000in}}%
\pgfpathlineto{\pgfqpoint{0.000000in}{0.027778in}}%
\pgfusepath{stroke,fill}%
}%
\begin{pgfscope}%
\pgfsys@transformshift{1.843546in}{0.352669in}%
\pgfsys@useobject{currentmarker}{}%
\end{pgfscope}%
\end{pgfscope}%
\begin{pgfscope}%
\pgfsetbuttcap%
\pgfsetroundjoin%
\definecolor{currentfill}{rgb}{0.000000,0.000000,0.000000}%
\pgfsetfillcolor{currentfill}%
\pgfsetlinewidth{0.301125pt}%
\definecolor{currentstroke}{rgb}{0.000000,0.000000,0.000000}%
\pgfsetstrokecolor{currentstroke}%
\pgfsetdash{}{0pt}%
\pgfsys@defobject{currentmarker}{\pgfqpoint{0.000000in}{0.000000in}}{\pgfqpoint{0.000000in}{0.027778in}}{%
\pgfpathmoveto{\pgfqpoint{0.000000in}{0.000000in}}%
\pgfpathlineto{\pgfqpoint{0.000000in}{0.027778in}}%
\pgfusepath{stroke,fill}%
}%
\begin{pgfscope}%
\pgfsys@transformshift{1.919727in}{0.352669in}%
\pgfsys@useobject{currentmarker}{}%
\end{pgfscope}%
\end{pgfscope}%
\begin{pgfscope}%
\pgfsetbuttcap%
\pgfsetroundjoin%
\definecolor{currentfill}{rgb}{0.000000,0.000000,0.000000}%
\pgfsetfillcolor{currentfill}%
\pgfsetlinewidth{0.301125pt}%
\definecolor{currentstroke}{rgb}{0.000000,0.000000,0.000000}%
\pgfsetstrokecolor{currentstroke}%
\pgfsetdash{}{0pt}%
\pgfsys@defobject{currentmarker}{\pgfqpoint{0.000000in}{0.000000in}}{\pgfqpoint{0.000000in}{0.027778in}}{%
\pgfpathmoveto{\pgfqpoint{0.000000in}{0.000000in}}%
\pgfpathlineto{\pgfqpoint{0.000000in}{0.027778in}}%
\pgfusepath{stroke,fill}%
}%
\begin{pgfscope}%
\pgfsys@transformshift{1.995907in}{0.352669in}%
\pgfsys@useobject{currentmarker}{}%
\end{pgfscope}%
\end{pgfscope}%
\begin{pgfscope}%
\pgfsetbuttcap%
\pgfsetroundjoin%
\definecolor{currentfill}{rgb}{0.000000,0.000000,0.000000}%
\pgfsetfillcolor{currentfill}%
\pgfsetlinewidth{0.301125pt}%
\definecolor{currentstroke}{rgb}{0.000000,0.000000,0.000000}%
\pgfsetstrokecolor{currentstroke}%
\pgfsetdash{}{0pt}%
\pgfsys@defobject{currentmarker}{\pgfqpoint{0.000000in}{0.000000in}}{\pgfqpoint{0.000000in}{0.027778in}}{%
\pgfpathmoveto{\pgfqpoint{0.000000in}{0.000000in}}%
\pgfpathlineto{\pgfqpoint{0.000000in}{0.027778in}}%
\pgfusepath{stroke,fill}%
}%
\begin{pgfscope}%
\pgfsys@transformshift{2.072087in}{0.352669in}%
\pgfsys@useobject{currentmarker}{}%
\end{pgfscope}%
\end{pgfscope}%
\begin{pgfscope}%
\pgfsetbuttcap%
\pgfsetroundjoin%
\definecolor{currentfill}{rgb}{0.000000,0.000000,0.000000}%
\pgfsetfillcolor{currentfill}%
\pgfsetlinewidth{0.301125pt}%
\definecolor{currentstroke}{rgb}{0.000000,0.000000,0.000000}%
\pgfsetstrokecolor{currentstroke}%
\pgfsetdash{}{0pt}%
\pgfsys@defobject{currentmarker}{\pgfqpoint{0.000000in}{0.000000in}}{\pgfqpoint{0.000000in}{0.027778in}}{%
\pgfpathmoveto{\pgfqpoint{0.000000in}{0.000000in}}%
\pgfpathlineto{\pgfqpoint{0.000000in}{0.027778in}}%
\pgfusepath{stroke,fill}%
}%
\begin{pgfscope}%
\pgfsys@transformshift{2.224448in}{0.352669in}%
\pgfsys@useobject{currentmarker}{}%
\end{pgfscope}%
\end{pgfscope}%
\begin{pgfscope}%
\pgfsetbuttcap%
\pgfsetroundjoin%
\definecolor{currentfill}{rgb}{0.000000,0.000000,0.000000}%
\pgfsetfillcolor{currentfill}%
\pgfsetlinewidth{0.301125pt}%
\definecolor{currentstroke}{rgb}{0.000000,0.000000,0.000000}%
\pgfsetstrokecolor{currentstroke}%
\pgfsetdash{}{0pt}%
\pgfsys@defobject{currentmarker}{\pgfqpoint{0.000000in}{0.000000in}}{\pgfqpoint{0.000000in}{0.027778in}}{%
\pgfpathmoveto{\pgfqpoint{0.000000in}{0.000000in}}%
\pgfpathlineto{\pgfqpoint{0.000000in}{0.027778in}}%
\pgfusepath{stroke,fill}%
}%
\begin{pgfscope}%
\pgfsys@transformshift{2.300628in}{0.352669in}%
\pgfsys@useobject{currentmarker}{}%
\end{pgfscope}%
\end{pgfscope}%
\begin{pgfscope}%
\pgfsetbuttcap%
\pgfsetroundjoin%
\definecolor{currentfill}{rgb}{0.000000,0.000000,0.000000}%
\pgfsetfillcolor{currentfill}%
\pgfsetlinewidth{0.301125pt}%
\definecolor{currentstroke}{rgb}{0.000000,0.000000,0.000000}%
\pgfsetstrokecolor{currentstroke}%
\pgfsetdash{}{0pt}%
\pgfsys@defobject{currentmarker}{\pgfqpoint{0.000000in}{0.000000in}}{\pgfqpoint{0.000000in}{0.027778in}}{%
\pgfpathmoveto{\pgfqpoint{0.000000in}{0.000000in}}%
\pgfpathlineto{\pgfqpoint{0.000000in}{0.027778in}}%
\pgfusepath{stroke,fill}%
}%
\begin{pgfscope}%
\pgfsys@transformshift{2.376809in}{0.352669in}%
\pgfsys@useobject{currentmarker}{}%
\end{pgfscope}%
\end{pgfscope}%
\begin{pgfscope}%
\pgfsetbuttcap%
\pgfsetroundjoin%
\definecolor{currentfill}{rgb}{0.000000,0.000000,0.000000}%
\pgfsetfillcolor{currentfill}%
\pgfsetlinewidth{0.301125pt}%
\definecolor{currentstroke}{rgb}{0.000000,0.000000,0.000000}%
\pgfsetstrokecolor{currentstroke}%
\pgfsetdash{}{0pt}%
\pgfsys@defobject{currentmarker}{\pgfqpoint{0.000000in}{0.000000in}}{\pgfqpoint{0.000000in}{0.027778in}}{%
\pgfpathmoveto{\pgfqpoint{0.000000in}{0.000000in}}%
\pgfpathlineto{\pgfqpoint{0.000000in}{0.027778in}}%
\pgfusepath{stroke,fill}%
}%
\begin{pgfscope}%
\pgfsys@transformshift{2.452989in}{0.352669in}%
\pgfsys@useobject{currentmarker}{}%
\end{pgfscope}%
\end{pgfscope}%
\begin{pgfscope}%
\pgfsetbuttcap%
\pgfsetroundjoin%
\definecolor{currentfill}{rgb}{0.000000,0.000000,0.000000}%
\pgfsetfillcolor{currentfill}%
\pgfsetlinewidth{0.301125pt}%
\definecolor{currentstroke}{rgb}{0.000000,0.000000,0.000000}%
\pgfsetstrokecolor{currentstroke}%
\pgfsetdash{}{0pt}%
\pgfsys@defobject{currentmarker}{\pgfqpoint{0.000000in}{0.000000in}}{\pgfqpoint{0.000000in}{0.027778in}}{%
\pgfpathmoveto{\pgfqpoint{0.000000in}{0.000000in}}%
\pgfpathlineto{\pgfqpoint{0.000000in}{0.027778in}}%
\pgfusepath{stroke,fill}%
}%
\begin{pgfscope}%
\pgfsys@transformshift{2.605350in}{0.352669in}%
\pgfsys@useobject{currentmarker}{}%
\end{pgfscope}%
\end{pgfscope}%
\begin{pgfscope}%
\pgfsetbuttcap%
\pgfsetroundjoin%
\definecolor{currentfill}{rgb}{0.000000,0.000000,0.000000}%
\pgfsetfillcolor{currentfill}%
\pgfsetlinewidth{0.301125pt}%
\definecolor{currentstroke}{rgb}{0.000000,0.000000,0.000000}%
\pgfsetstrokecolor{currentstroke}%
\pgfsetdash{}{0pt}%
\pgfsys@defobject{currentmarker}{\pgfqpoint{0.000000in}{0.000000in}}{\pgfqpoint{0.000000in}{0.027778in}}{%
\pgfpathmoveto{\pgfqpoint{0.000000in}{0.000000in}}%
\pgfpathlineto{\pgfqpoint{0.000000in}{0.027778in}}%
\pgfusepath{stroke,fill}%
}%
\begin{pgfscope}%
\pgfsys@transformshift{2.681530in}{0.352669in}%
\pgfsys@useobject{currentmarker}{}%
\end{pgfscope}%
\end{pgfscope}%
\begin{pgfscope}%
\pgfsetbuttcap%
\pgfsetroundjoin%
\definecolor{currentfill}{rgb}{0.000000,0.000000,0.000000}%
\pgfsetfillcolor{currentfill}%
\pgfsetlinewidth{0.301125pt}%
\definecolor{currentstroke}{rgb}{0.000000,0.000000,0.000000}%
\pgfsetstrokecolor{currentstroke}%
\pgfsetdash{}{0pt}%
\pgfsys@defobject{currentmarker}{\pgfqpoint{0.000000in}{0.000000in}}{\pgfqpoint{0.000000in}{0.027778in}}{%
\pgfpathmoveto{\pgfqpoint{0.000000in}{0.000000in}}%
\pgfpathlineto{\pgfqpoint{0.000000in}{0.027778in}}%
\pgfusepath{stroke,fill}%
}%
\begin{pgfscope}%
\pgfsys@transformshift{2.757710in}{0.352669in}%
\pgfsys@useobject{currentmarker}{}%
\end{pgfscope}%
\end{pgfscope}%
\begin{pgfscope}%
\pgfsetbuttcap%
\pgfsetroundjoin%
\definecolor{currentfill}{rgb}{0.000000,0.000000,0.000000}%
\pgfsetfillcolor{currentfill}%
\pgfsetlinewidth{0.301125pt}%
\definecolor{currentstroke}{rgb}{0.000000,0.000000,0.000000}%
\pgfsetstrokecolor{currentstroke}%
\pgfsetdash{}{0pt}%
\pgfsys@defobject{currentmarker}{\pgfqpoint{0.000000in}{0.000000in}}{\pgfqpoint{0.000000in}{0.027778in}}{%
\pgfpathmoveto{\pgfqpoint{0.000000in}{0.000000in}}%
\pgfpathlineto{\pgfqpoint{0.000000in}{0.027778in}}%
\pgfusepath{stroke,fill}%
}%
\begin{pgfscope}%
\pgfsys@transformshift{2.833891in}{0.352669in}%
\pgfsys@useobject{currentmarker}{}%
\end{pgfscope}%
\end{pgfscope}%
\begin{pgfscope}%
\pgfsetbuttcap%
\pgfsetroundjoin%
\definecolor{currentfill}{rgb}{0.000000,0.000000,0.000000}%
\pgfsetfillcolor{currentfill}%
\pgfsetlinewidth{0.301125pt}%
\definecolor{currentstroke}{rgb}{0.000000,0.000000,0.000000}%
\pgfsetstrokecolor{currentstroke}%
\pgfsetdash{}{0pt}%
\pgfsys@defobject{currentmarker}{\pgfqpoint{0.000000in}{0.000000in}}{\pgfqpoint{0.000000in}{0.027778in}}{%
\pgfpathmoveto{\pgfqpoint{0.000000in}{0.000000in}}%
\pgfpathlineto{\pgfqpoint{0.000000in}{0.027778in}}%
\pgfusepath{stroke,fill}%
}%
\begin{pgfscope}%
\pgfsys@transformshift{2.986251in}{0.352669in}%
\pgfsys@useobject{currentmarker}{}%
\end{pgfscope}%
\end{pgfscope}%
\begin{pgfscope}%
\pgfsetbuttcap%
\pgfsetroundjoin%
\definecolor{currentfill}{rgb}{0.000000,0.000000,0.000000}%
\pgfsetfillcolor{currentfill}%
\pgfsetlinewidth{0.301125pt}%
\definecolor{currentstroke}{rgb}{0.000000,0.000000,0.000000}%
\pgfsetstrokecolor{currentstroke}%
\pgfsetdash{}{0pt}%
\pgfsys@defobject{currentmarker}{\pgfqpoint{0.000000in}{0.000000in}}{\pgfqpoint{0.000000in}{0.027778in}}{%
\pgfpathmoveto{\pgfqpoint{0.000000in}{0.000000in}}%
\pgfpathlineto{\pgfqpoint{0.000000in}{0.027778in}}%
\pgfusepath{stroke,fill}%
}%
\begin{pgfscope}%
\pgfsys@transformshift{3.062432in}{0.352669in}%
\pgfsys@useobject{currentmarker}{}%
\end{pgfscope}%
\end{pgfscope}%
\begin{pgfscope}%
\pgfsetbuttcap%
\pgfsetroundjoin%
\definecolor{currentfill}{rgb}{0.000000,0.000000,0.000000}%
\pgfsetfillcolor{currentfill}%
\pgfsetlinewidth{0.301125pt}%
\definecolor{currentstroke}{rgb}{0.000000,0.000000,0.000000}%
\pgfsetstrokecolor{currentstroke}%
\pgfsetdash{}{0pt}%
\pgfsys@defobject{currentmarker}{\pgfqpoint{0.000000in}{0.000000in}}{\pgfqpoint{0.000000in}{0.027778in}}{%
\pgfpathmoveto{\pgfqpoint{0.000000in}{0.000000in}}%
\pgfpathlineto{\pgfqpoint{0.000000in}{0.027778in}}%
\pgfusepath{stroke,fill}%
}%
\begin{pgfscope}%
\pgfsys@transformshift{3.138612in}{0.352669in}%
\pgfsys@useobject{currentmarker}{}%
\end{pgfscope}%
\end{pgfscope}%
\begin{pgfscope}%
\pgfsetbuttcap%
\pgfsetroundjoin%
\definecolor{currentfill}{rgb}{0.000000,0.000000,0.000000}%
\pgfsetfillcolor{currentfill}%
\pgfsetlinewidth{0.301125pt}%
\definecolor{currentstroke}{rgb}{0.000000,0.000000,0.000000}%
\pgfsetstrokecolor{currentstroke}%
\pgfsetdash{}{0pt}%
\pgfsys@defobject{currentmarker}{\pgfqpoint{0.000000in}{0.000000in}}{\pgfqpoint{0.000000in}{0.027778in}}{%
\pgfpathmoveto{\pgfqpoint{0.000000in}{0.000000in}}%
\pgfpathlineto{\pgfqpoint{0.000000in}{0.027778in}}%
\pgfusepath{stroke,fill}%
}%
\begin{pgfscope}%
\pgfsys@transformshift{3.214792in}{0.352669in}%
\pgfsys@useobject{currentmarker}{}%
\end{pgfscope}%
\end{pgfscope}%
\begin{pgfscope}%
\pgfsetbuttcap%
\pgfsetroundjoin%
\definecolor{currentfill}{rgb}{0.000000,0.000000,0.000000}%
\pgfsetfillcolor{currentfill}%
\pgfsetlinewidth{0.301125pt}%
\definecolor{currentstroke}{rgb}{0.000000,0.000000,0.000000}%
\pgfsetstrokecolor{currentstroke}%
\pgfsetdash{}{0pt}%
\pgfsys@defobject{currentmarker}{\pgfqpoint{0.000000in}{0.000000in}}{\pgfqpoint{0.000000in}{0.027778in}}{%
\pgfpathmoveto{\pgfqpoint{0.000000in}{0.000000in}}%
\pgfpathlineto{\pgfqpoint{0.000000in}{0.027778in}}%
\pgfusepath{stroke,fill}%
}%
\begin{pgfscope}%
\pgfsys@transformshift{3.367153in}{0.352669in}%
\pgfsys@useobject{currentmarker}{}%
\end{pgfscope}%
\end{pgfscope}%
\begin{pgfscope}%
\pgfsetbuttcap%
\pgfsetroundjoin%
\definecolor{currentfill}{rgb}{0.000000,0.000000,0.000000}%
\pgfsetfillcolor{currentfill}%
\pgfsetlinewidth{0.301125pt}%
\definecolor{currentstroke}{rgb}{0.000000,0.000000,0.000000}%
\pgfsetstrokecolor{currentstroke}%
\pgfsetdash{}{0pt}%
\pgfsys@defobject{currentmarker}{\pgfqpoint{0.000000in}{0.000000in}}{\pgfqpoint{0.000000in}{0.027778in}}{%
\pgfpathmoveto{\pgfqpoint{0.000000in}{0.000000in}}%
\pgfpathlineto{\pgfqpoint{0.000000in}{0.027778in}}%
\pgfusepath{stroke,fill}%
}%
\begin{pgfscope}%
\pgfsys@transformshift{3.443333in}{0.352669in}%
\pgfsys@useobject{currentmarker}{}%
\end{pgfscope}%
\end{pgfscope}%
\begin{pgfscope}%
\pgfsetbuttcap%
\pgfsetroundjoin%
\definecolor{currentfill}{rgb}{0.000000,0.000000,0.000000}%
\pgfsetfillcolor{currentfill}%
\pgfsetlinewidth{0.301125pt}%
\definecolor{currentstroke}{rgb}{0.000000,0.000000,0.000000}%
\pgfsetstrokecolor{currentstroke}%
\pgfsetdash{}{0pt}%
\pgfsys@defobject{currentmarker}{\pgfqpoint{0.000000in}{0.000000in}}{\pgfqpoint{0.000000in}{0.027778in}}{%
\pgfpathmoveto{\pgfqpoint{0.000000in}{0.000000in}}%
\pgfpathlineto{\pgfqpoint{0.000000in}{0.027778in}}%
\pgfusepath{stroke,fill}%
}%
\begin{pgfscope}%
\pgfsys@transformshift{3.519514in}{0.352669in}%
\pgfsys@useobject{currentmarker}{}%
\end{pgfscope}%
\end{pgfscope}%
\begin{pgfscope}%
\pgfsetbuttcap%
\pgfsetroundjoin%
\definecolor{currentfill}{rgb}{0.000000,0.000000,0.000000}%
\pgfsetfillcolor{currentfill}%
\pgfsetlinewidth{0.301125pt}%
\definecolor{currentstroke}{rgb}{0.000000,0.000000,0.000000}%
\pgfsetstrokecolor{currentstroke}%
\pgfsetdash{}{0pt}%
\pgfsys@defobject{currentmarker}{\pgfqpoint{0.000000in}{0.000000in}}{\pgfqpoint{0.000000in}{0.027778in}}{%
\pgfpathmoveto{\pgfqpoint{0.000000in}{0.000000in}}%
\pgfpathlineto{\pgfqpoint{0.000000in}{0.027778in}}%
\pgfusepath{stroke,fill}%
}%
\begin{pgfscope}%
\pgfsys@transformshift{3.595694in}{0.352669in}%
\pgfsys@useobject{currentmarker}{}%
\end{pgfscope}%
\end{pgfscope}%
\begin{pgfscope}%
\pgfsetbuttcap%
\pgfsetroundjoin%
\definecolor{currentfill}{rgb}{0.000000,0.000000,0.000000}%
\pgfsetfillcolor{currentfill}%
\pgfsetlinewidth{0.301125pt}%
\definecolor{currentstroke}{rgb}{0.000000,0.000000,0.000000}%
\pgfsetstrokecolor{currentstroke}%
\pgfsetdash{}{0pt}%
\pgfsys@defobject{currentmarker}{\pgfqpoint{0.000000in}{0.000000in}}{\pgfqpoint{0.000000in}{0.027778in}}{%
\pgfpathmoveto{\pgfqpoint{0.000000in}{0.000000in}}%
\pgfpathlineto{\pgfqpoint{0.000000in}{0.027778in}}%
\pgfusepath{stroke,fill}%
}%
\begin{pgfscope}%
\pgfsys@transformshift{3.748055in}{0.352669in}%
\pgfsys@useobject{currentmarker}{}%
\end{pgfscope}%
\end{pgfscope}%
\begin{pgfscope}%
\pgfsetbuttcap%
\pgfsetroundjoin%
\definecolor{currentfill}{rgb}{0.000000,0.000000,0.000000}%
\pgfsetfillcolor{currentfill}%
\pgfsetlinewidth{0.301125pt}%
\definecolor{currentstroke}{rgb}{0.000000,0.000000,0.000000}%
\pgfsetstrokecolor{currentstroke}%
\pgfsetdash{}{0pt}%
\pgfsys@defobject{currentmarker}{\pgfqpoint{0.000000in}{0.000000in}}{\pgfqpoint{0.000000in}{0.027778in}}{%
\pgfpathmoveto{\pgfqpoint{0.000000in}{0.000000in}}%
\pgfpathlineto{\pgfqpoint{0.000000in}{0.027778in}}%
\pgfusepath{stroke,fill}%
}%
\begin{pgfscope}%
\pgfsys@transformshift{3.824235in}{0.352669in}%
\pgfsys@useobject{currentmarker}{}%
\end{pgfscope}%
\end{pgfscope}%
\begin{pgfscope}%
\definecolor{textcolor}{rgb}{0.000000,0.000000,0.000000}%
\pgfsetstrokecolor{textcolor}%
\pgfsetfillcolor{textcolor}%
\pgftext[x=2.148268in,y=0.140972in,,top]{\color{textcolor}{\sffamily\fontsize{9.000000}{10.800000}\selectfont\catcode`\^=\active\def^{\ifmmode\sp\else\^{}\fi}\catcode`\%=\active\def%{\%}$x / M$}}%
\end{pgfscope}%
\begin{pgfscope}%
\pgfsetbuttcap%
\pgfsetroundjoin%
\definecolor{currentfill}{rgb}{0.000000,0.000000,0.000000}%
\pgfsetfillcolor{currentfill}%
\pgfsetlinewidth{0.501875pt}%
\definecolor{currentstroke}{rgb}{0.000000,0.000000,0.000000}%
\pgfsetstrokecolor{currentstroke}%
\pgfsetdash{}{0pt}%
\pgfsys@defobject{currentmarker}{\pgfqpoint{0.000000in}{0.000000in}}{\pgfqpoint{0.041667in}{0.000000in}}{%
\pgfpathmoveto{\pgfqpoint{0.000000in}{0.000000in}}%
\pgfpathlineto{\pgfqpoint{0.041667in}{0.000000in}}%
\pgfusepath{stroke,fill}%
}%
\begin{pgfscope}%
\pgfsys@transformshift{0.472301in}{0.505029in}%
\pgfsys@useobject{currentmarker}{}%
\end{pgfscope}%
\end{pgfscope}%
\begin{pgfscope}%
\definecolor{textcolor}{rgb}{0.000000,0.000000,0.000000}%
\pgfsetstrokecolor{textcolor}%
\pgfsetfillcolor{textcolor}%
\pgftext[x=0.196527in, y=0.462820in, left, base]{\color{textcolor}{\sffamily\fontsize{8.000000}{9.600000}\selectfont\catcode`\^=\active\def^{\ifmmode\sp\else\^{}\fi}\catcode`\%=\active\def%{\%}\ensuremath{-}20}}%
\end{pgfscope}%
\begin{pgfscope}%
\pgfsetbuttcap%
\pgfsetroundjoin%
\definecolor{currentfill}{rgb}{0.000000,0.000000,0.000000}%
\pgfsetfillcolor{currentfill}%
\pgfsetlinewidth{0.501875pt}%
\definecolor{currentstroke}{rgb}{0.000000,0.000000,0.000000}%
\pgfsetstrokecolor{currentstroke}%
\pgfsetdash{}{0pt}%
\pgfsys@defobject{currentmarker}{\pgfqpoint{0.000000in}{0.000000in}}{\pgfqpoint{0.041667in}{0.000000in}}{%
\pgfpathmoveto{\pgfqpoint{0.000000in}{0.000000in}}%
\pgfpathlineto{\pgfqpoint{0.041667in}{0.000000in}}%
\pgfusepath{stroke,fill}%
}%
\begin{pgfscope}%
\pgfsys@transformshift{0.472301in}{0.885931in}%
\pgfsys@useobject{currentmarker}{}%
\end{pgfscope}%
\end{pgfscope}%
\begin{pgfscope}%
\definecolor{textcolor}{rgb}{0.000000,0.000000,0.000000}%
\pgfsetstrokecolor{textcolor}%
\pgfsetfillcolor{textcolor}%
\pgftext[x=0.196527in, y=0.843722in, left, base]{\color{textcolor}{\sffamily\fontsize{8.000000}{9.600000}\selectfont\catcode`\^=\active\def^{\ifmmode\sp\else\^{}\fi}\catcode`\%=\active\def%{\%}\ensuremath{-}15}}%
\end{pgfscope}%
\begin{pgfscope}%
\pgfsetbuttcap%
\pgfsetroundjoin%
\definecolor{currentfill}{rgb}{0.000000,0.000000,0.000000}%
\pgfsetfillcolor{currentfill}%
\pgfsetlinewidth{0.501875pt}%
\definecolor{currentstroke}{rgb}{0.000000,0.000000,0.000000}%
\pgfsetstrokecolor{currentstroke}%
\pgfsetdash{}{0pt}%
\pgfsys@defobject{currentmarker}{\pgfqpoint{0.000000in}{0.000000in}}{\pgfqpoint{0.041667in}{0.000000in}}{%
\pgfpathmoveto{\pgfqpoint{0.000000in}{0.000000in}}%
\pgfpathlineto{\pgfqpoint{0.041667in}{0.000000in}}%
\pgfusepath{stroke,fill}%
}%
\begin{pgfscope}%
\pgfsys@transformshift{0.472301in}{1.266833in}%
\pgfsys@useobject{currentmarker}{}%
\end{pgfscope}%
\end{pgfscope}%
\begin{pgfscope}%
\definecolor{textcolor}{rgb}{0.000000,0.000000,0.000000}%
\pgfsetstrokecolor{textcolor}%
\pgfsetfillcolor{textcolor}%
\pgftext[x=0.196527in, y=1.224623in, left, base]{\color{textcolor}{\sffamily\fontsize{8.000000}{9.600000}\selectfont\catcode`\^=\active\def^{\ifmmode\sp\else\^{}\fi}\catcode`\%=\active\def%{\%}\ensuremath{-}10}}%
\end{pgfscope}%
\begin{pgfscope}%
\pgfsetbuttcap%
\pgfsetroundjoin%
\definecolor{currentfill}{rgb}{0.000000,0.000000,0.000000}%
\pgfsetfillcolor{currentfill}%
\pgfsetlinewidth{0.501875pt}%
\definecolor{currentstroke}{rgb}{0.000000,0.000000,0.000000}%
\pgfsetstrokecolor{currentstroke}%
\pgfsetdash{}{0pt}%
\pgfsys@defobject{currentmarker}{\pgfqpoint{0.000000in}{0.000000in}}{\pgfqpoint{0.041667in}{0.000000in}}{%
\pgfpathmoveto{\pgfqpoint{0.000000in}{0.000000in}}%
\pgfpathlineto{\pgfqpoint{0.041667in}{0.000000in}}%
\pgfusepath{stroke,fill}%
}%
\begin{pgfscope}%
\pgfsys@transformshift{0.472301in}{1.647734in}%
\pgfsys@useobject{currentmarker}{}%
\end{pgfscope}%
\end{pgfscope}%
\begin{pgfscope}%
\definecolor{textcolor}{rgb}{0.000000,0.000000,0.000000}%
\pgfsetstrokecolor{textcolor}%
\pgfsetfillcolor{textcolor}%
\pgftext[x=0.267219in, y=1.605525in, left, base]{\color{textcolor}{\sffamily\fontsize{8.000000}{9.600000}\selectfont\catcode`\^=\active\def^{\ifmmode\sp\else\^{}\fi}\catcode`\%=\active\def%{\%}\ensuremath{-}5}}%
\end{pgfscope}%
\begin{pgfscope}%
\pgfsetbuttcap%
\pgfsetroundjoin%
\definecolor{currentfill}{rgb}{0.000000,0.000000,0.000000}%
\pgfsetfillcolor{currentfill}%
\pgfsetlinewidth{0.501875pt}%
\definecolor{currentstroke}{rgb}{0.000000,0.000000,0.000000}%
\pgfsetstrokecolor{currentstroke}%
\pgfsetdash{}{0pt}%
\pgfsys@defobject{currentmarker}{\pgfqpoint{0.000000in}{0.000000in}}{\pgfqpoint{0.041667in}{0.000000in}}{%
\pgfpathmoveto{\pgfqpoint{0.000000in}{0.000000in}}%
\pgfpathlineto{\pgfqpoint{0.041667in}{0.000000in}}%
\pgfusepath{stroke,fill}%
}%
\begin{pgfscope}%
\pgfsys@transformshift{0.472301in}{2.028636in}%
\pgfsys@useobject{currentmarker}{}%
\end{pgfscope}%
\end{pgfscope}%
\begin{pgfscope}%
\definecolor{textcolor}{rgb}{0.000000,0.000000,0.000000}%
\pgfsetstrokecolor{textcolor}%
\pgfsetfillcolor{textcolor}%
\pgftext[x=0.352997in, y=1.986427in, left, base]{\color{textcolor}{\sffamily\fontsize{8.000000}{9.600000}\selectfont\catcode`\^=\active\def^{\ifmmode\sp\else\^{}\fi}\catcode`\%=\active\def%{\%}0}}%
\end{pgfscope}%
\begin{pgfscope}%
\pgfsetbuttcap%
\pgfsetroundjoin%
\definecolor{currentfill}{rgb}{0.000000,0.000000,0.000000}%
\pgfsetfillcolor{currentfill}%
\pgfsetlinewidth{0.501875pt}%
\definecolor{currentstroke}{rgb}{0.000000,0.000000,0.000000}%
\pgfsetstrokecolor{currentstroke}%
\pgfsetdash{}{0pt}%
\pgfsys@defobject{currentmarker}{\pgfqpoint{0.000000in}{0.000000in}}{\pgfqpoint{0.041667in}{0.000000in}}{%
\pgfpathmoveto{\pgfqpoint{0.000000in}{0.000000in}}%
\pgfpathlineto{\pgfqpoint{0.041667in}{0.000000in}}%
\pgfusepath{stroke,fill}%
}%
\begin{pgfscope}%
\pgfsys@transformshift{0.472301in}{2.409538in}%
\pgfsys@useobject{currentmarker}{}%
\end{pgfscope}%
\end{pgfscope}%
\begin{pgfscope}%
\definecolor{textcolor}{rgb}{0.000000,0.000000,0.000000}%
\pgfsetstrokecolor{textcolor}%
\pgfsetfillcolor{textcolor}%
\pgftext[x=0.352997in, y=2.367328in, left, base]{\color{textcolor}{\sffamily\fontsize{8.000000}{9.600000}\selectfont\catcode`\^=\active\def^{\ifmmode\sp\else\^{}\fi}\catcode`\%=\active\def%{\%}5}}%
\end{pgfscope}%
\begin{pgfscope}%
\pgfsetbuttcap%
\pgfsetroundjoin%
\definecolor{currentfill}{rgb}{0.000000,0.000000,0.000000}%
\pgfsetfillcolor{currentfill}%
\pgfsetlinewidth{0.501875pt}%
\definecolor{currentstroke}{rgb}{0.000000,0.000000,0.000000}%
\pgfsetstrokecolor{currentstroke}%
\pgfsetdash{}{0pt}%
\pgfsys@defobject{currentmarker}{\pgfqpoint{0.000000in}{0.000000in}}{\pgfqpoint{0.041667in}{0.000000in}}{%
\pgfpathmoveto{\pgfqpoint{0.000000in}{0.000000in}}%
\pgfpathlineto{\pgfqpoint{0.041667in}{0.000000in}}%
\pgfusepath{stroke,fill}%
}%
\begin{pgfscope}%
\pgfsys@transformshift{0.472301in}{2.790439in}%
\pgfsys@useobject{currentmarker}{}%
\end{pgfscope}%
\end{pgfscope}%
\begin{pgfscope}%
\definecolor{textcolor}{rgb}{0.000000,0.000000,0.000000}%
\pgfsetstrokecolor{textcolor}%
\pgfsetfillcolor{textcolor}%
\pgftext[x=0.282305in, y=2.748230in, left, base]{\color{textcolor}{\sffamily\fontsize{8.000000}{9.600000}\selectfont\catcode`\^=\active\def^{\ifmmode\sp\else\^{}\fi}\catcode`\%=\active\def%{\%}10}}%
\end{pgfscope}%
\begin{pgfscope}%
\pgfsetbuttcap%
\pgfsetroundjoin%
\definecolor{currentfill}{rgb}{0.000000,0.000000,0.000000}%
\pgfsetfillcolor{currentfill}%
\pgfsetlinewidth{0.501875pt}%
\definecolor{currentstroke}{rgb}{0.000000,0.000000,0.000000}%
\pgfsetstrokecolor{currentstroke}%
\pgfsetdash{}{0pt}%
\pgfsys@defobject{currentmarker}{\pgfqpoint{0.000000in}{0.000000in}}{\pgfqpoint{0.041667in}{0.000000in}}{%
\pgfpathmoveto{\pgfqpoint{0.000000in}{0.000000in}}%
\pgfpathlineto{\pgfqpoint{0.041667in}{0.000000in}}%
\pgfusepath{stroke,fill}%
}%
\begin{pgfscope}%
\pgfsys@transformshift{0.472301in}{3.171341in}%
\pgfsys@useobject{currentmarker}{}%
\end{pgfscope}%
\end{pgfscope}%
\begin{pgfscope}%
\definecolor{textcolor}{rgb}{0.000000,0.000000,0.000000}%
\pgfsetstrokecolor{textcolor}%
\pgfsetfillcolor{textcolor}%
\pgftext[x=0.282305in, y=3.129132in, left, base]{\color{textcolor}{\sffamily\fontsize{8.000000}{9.600000}\selectfont\catcode`\^=\active\def^{\ifmmode\sp\else\^{}\fi}\catcode`\%=\active\def%{\%}15}}%
\end{pgfscope}%
\begin{pgfscope}%
\pgfsetbuttcap%
\pgfsetroundjoin%
\definecolor{currentfill}{rgb}{0.000000,0.000000,0.000000}%
\pgfsetfillcolor{currentfill}%
\pgfsetlinewidth{0.501875pt}%
\definecolor{currentstroke}{rgb}{0.000000,0.000000,0.000000}%
\pgfsetstrokecolor{currentstroke}%
\pgfsetdash{}{0pt}%
\pgfsys@defobject{currentmarker}{\pgfqpoint{0.000000in}{0.000000in}}{\pgfqpoint{0.041667in}{0.000000in}}{%
\pgfpathmoveto{\pgfqpoint{0.000000in}{0.000000in}}%
\pgfpathlineto{\pgfqpoint{0.041667in}{0.000000in}}%
\pgfusepath{stroke,fill}%
}%
\begin{pgfscope}%
\pgfsys@transformshift{0.472301in}{3.552243in}%
\pgfsys@useobject{currentmarker}{}%
\end{pgfscope}%
\end{pgfscope}%
\begin{pgfscope}%
\definecolor{textcolor}{rgb}{0.000000,0.000000,0.000000}%
\pgfsetstrokecolor{textcolor}%
\pgfsetfillcolor{textcolor}%
\pgftext[x=0.282305in, y=3.510033in, left, base]{\color{textcolor}{\sffamily\fontsize{8.000000}{9.600000}\selectfont\catcode`\^=\active\def^{\ifmmode\sp\else\^{}\fi}\catcode`\%=\active\def%{\%}20}}%
\end{pgfscope}%
\begin{pgfscope}%
\pgfsetbuttcap%
\pgfsetroundjoin%
\definecolor{currentfill}{rgb}{0.000000,0.000000,0.000000}%
\pgfsetfillcolor{currentfill}%
\pgfsetlinewidth{0.301125pt}%
\definecolor{currentstroke}{rgb}{0.000000,0.000000,0.000000}%
\pgfsetstrokecolor{currentstroke}%
\pgfsetdash{}{0pt}%
\pgfsys@defobject{currentmarker}{\pgfqpoint{0.000000in}{0.000000in}}{\pgfqpoint{0.027778in}{0.000000in}}{%
\pgfpathmoveto{\pgfqpoint{0.000000in}{0.000000in}}%
\pgfpathlineto{\pgfqpoint{0.027778in}{0.000000in}}%
\pgfusepath{stroke,fill}%
}%
\begin{pgfscope}%
\pgfsys@transformshift{0.472301in}{0.352669in}%
\pgfsys@useobject{currentmarker}{}%
\end{pgfscope}%
\end{pgfscope}%
\begin{pgfscope}%
\pgfsetbuttcap%
\pgfsetroundjoin%
\definecolor{currentfill}{rgb}{0.000000,0.000000,0.000000}%
\pgfsetfillcolor{currentfill}%
\pgfsetlinewidth{0.301125pt}%
\definecolor{currentstroke}{rgb}{0.000000,0.000000,0.000000}%
\pgfsetstrokecolor{currentstroke}%
\pgfsetdash{}{0pt}%
\pgfsys@defobject{currentmarker}{\pgfqpoint{0.000000in}{0.000000in}}{\pgfqpoint{0.027778in}{0.000000in}}{%
\pgfpathmoveto{\pgfqpoint{0.000000in}{0.000000in}}%
\pgfpathlineto{\pgfqpoint{0.027778in}{0.000000in}}%
\pgfusepath{stroke,fill}%
}%
\begin{pgfscope}%
\pgfsys@transformshift{0.472301in}{0.428849in}%
\pgfsys@useobject{currentmarker}{}%
\end{pgfscope}%
\end{pgfscope}%
\begin{pgfscope}%
\pgfsetbuttcap%
\pgfsetroundjoin%
\definecolor{currentfill}{rgb}{0.000000,0.000000,0.000000}%
\pgfsetfillcolor{currentfill}%
\pgfsetlinewidth{0.301125pt}%
\definecolor{currentstroke}{rgb}{0.000000,0.000000,0.000000}%
\pgfsetstrokecolor{currentstroke}%
\pgfsetdash{}{0pt}%
\pgfsys@defobject{currentmarker}{\pgfqpoint{0.000000in}{0.000000in}}{\pgfqpoint{0.027778in}{0.000000in}}{%
\pgfpathmoveto{\pgfqpoint{0.000000in}{0.000000in}}%
\pgfpathlineto{\pgfqpoint{0.027778in}{0.000000in}}%
\pgfusepath{stroke,fill}%
}%
\begin{pgfscope}%
\pgfsys@transformshift{0.472301in}{0.581210in}%
\pgfsys@useobject{currentmarker}{}%
\end{pgfscope}%
\end{pgfscope}%
\begin{pgfscope}%
\pgfsetbuttcap%
\pgfsetroundjoin%
\definecolor{currentfill}{rgb}{0.000000,0.000000,0.000000}%
\pgfsetfillcolor{currentfill}%
\pgfsetlinewidth{0.301125pt}%
\definecolor{currentstroke}{rgb}{0.000000,0.000000,0.000000}%
\pgfsetstrokecolor{currentstroke}%
\pgfsetdash{}{0pt}%
\pgfsys@defobject{currentmarker}{\pgfqpoint{0.000000in}{0.000000in}}{\pgfqpoint{0.027778in}{0.000000in}}{%
\pgfpathmoveto{\pgfqpoint{0.000000in}{0.000000in}}%
\pgfpathlineto{\pgfqpoint{0.027778in}{0.000000in}}%
\pgfusepath{stroke,fill}%
}%
\begin{pgfscope}%
\pgfsys@transformshift{0.472301in}{0.657390in}%
\pgfsys@useobject{currentmarker}{}%
\end{pgfscope}%
\end{pgfscope}%
\begin{pgfscope}%
\pgfsetbuttcap%
\pgfsetroundjoin%
\definecolor{currentfill}{rgb}{0.000000,0.000000,0.000000}%
\pgfsetfillcolor{currentfill}%
\pgfsetlinewidth{0.301125pt}%
\definecolor{currentstroke}{rgb}{0.000000,0.000000,0.000000}%
\pgfsetstrokecolor{currentstroke}%
\pgfsetdash{}{0pt}%
\pgfsys@defobject{currentmarker}{\pgfqpoint{0.000000in}{0.000000in}}{\pgfqpoint{0.027778in}{0.000000in}}{%
\pgfpathmoveto{\pgfqpoint{0.000000in}{0.000000in}}%
\pgfpathlineto{\pgfqpoint{0.027778in}{0.000000in}}%
\pgfusepath{stroke,fill}%
}%
\begin{pgfscope}%
\pgfsys@transformshift{0.472301in}{0.733570in}%
\pgfsys@useobject{currentmarker}{}%
\end{pgfscope}%
\end{pgfscope}%
\begin{pgfscope}%
\pgfsetbuttcap%
\pgfsetroundjoin%
\definecolor{currentfill}{rgb}{0.000000,0.000000,0.000000}%
\pgfsetfillcolor{currentfill}%
\pgfsetlinewidth{0.301125pt}%
\definecolor{currentstroke}{rgb}{0.000000,0.000000,0.000000}%
\pgfsetstrokecolor{currentstroke}%
\pgfsetdash{}{0pt}%
\pgfsys@defobject{currentmarker}{\pgfqpoint{0.000000in}{0.000000in}}{\pgfqpoint{0.027778in}{0.000000in}}{%
\pgfpathmoveto{\pgfqpoint{0.000000in}{0.000000in}}%
\pgfpathlineto{\pgfqpoint{0.027778in}{0.000000in}}%
\pgfusepath{stroke,fill}%
}%
\begin{pgfscope}%
\pgfsys@transformshift{0.472301in}{0.809751in}%
\pgfsys@useobject{currentmarker}{}%
\end{pgfscope}%
\end{pgfscope}%
\begin{pgfscope}%
\pgfsetbuttcap%
\pgfsetroundjoin%
\definecolor{currentfill}{rgb}{0.000000,0.000000,0.000000}%
\pgfsetfillcolor{currentfill}%
\pgfsetlinewidth{0.301125pt}%
\definecolor{currentstroke}{rgb}{0.000000,0.000000,0.000000}%
\pgfsetstrokecolor{currentstroke}%
\pgfsetdash{}{0pt}%
\pgfsys@defobject{currentmarker}{\pgfqpoint{0.000000in}{0.000000in}}{\pgfqpoint{0.027778in}{0.000000in}}{%
\pgfpathmoveto{\pgfqpoint{0.000000in}{0.000000in}}%
\pgfpathlineto{\pgfqpoint{0.027778in}{0.000000in}}%
\pgfusepath{stroke,fill}%
}%
\begin{pgfscope}%
\pgfsys@transformshift{0.472301in}{0.962111in}%
\pgfsys@useobject{currentmarker}{}%
\end{pgfscope}%
\end{pgfscope}%
\begin{pgfscope}%
\pgfsetbuttcap%
\pgfsetroundjoin%
\definecolor{currentfill}{rgb}{0.000000,0.000000,0.000000}%
\pgfsetfillcolor{currentfill}%
\pgfsetlinewidth{0.301125pt}%
\definecolor{currentstroke}{rgb}{0.000000,0.000000,0.000000}%
\pgfsetstrokecolor{currentstroke}%
\pgfsetdash{}{0pt}%
\pgfsys@defobject{currentmarker}{\pgfqpoint{0.000000in}{0.000000in}}{\pgfqpoint{0.027778in}{0.000000in}}{%
\pgfpathmoveto{\pgfqpoint{0.000000in}{0.000000in}}%
\pgfpathlineto{\pgfqpoint{0.027778in}{0.000000in}}%
\pgfusepath{stroke,fill}%
}%
\begin{pgfscope}%
\pgfsys@transformshift{0.472301in}{1.038292in}%
\pgfsys@useobject{currentmarker}{}%
\end{pgfscope}%
\end{pgfscope}%
\begin{pgfscope}%
\pgfsetbuttcap%
\pgfsetroundjoin%
\definecolor{currentfill}{rgb}{0.000000,0.000000,0.000000}%
\pgfsetfillcolor{currentfill}%
\pgfsetlinewidth{0.301125pt}%
\definecolor{currentstroke}{rgb}{0.000000,0.000000,0.000000}%
\pgfsetstrokecolor{currentstroke}%
\pgfsetdash{}{0pt}%
\pgfsys@defobject{currentmarker}{\pgfqpoint{0.000000in}{0.000000in}}{\pgfqpoint{0.027778in}{0.000000in}}{%
\pgfpathmoveto{\pgfqpoint{0.000000in}{0.000000in}}%
\pgfpathlineto{\pgfqpoint{0.027778in}{0.000000in}}%
\pgfusepath{stroke,fill}%
}%
\begin{pgfscope}%
\pgfsys@transformshift{0.472301in}{1.114472in}%
\pgfsys@useobject{currentmarker}{}%
\end{pgfscope}%
\end{pgfscope}%
\begin{pgfscope}%
\pgfsetbuttcap%
\pgfsetroundjoin%
\definecolor{currentfill}{rgb}{0.000000,0.000000,0.000000}%
\pgfsetfillcolor{currentfill}%
\pgfsetlinewidth{0.301125pt}%
\definecolor{currentstroke}{rgb}{0.000000,0.000000,0.000000}%
\pgfsetstrokecolor{currentstroke}%
\pgfsetdash{}{0pt}%
\pgfsys@defobject{currentmarker}{\pgfqpoint{0.000000in}{0.000000in}}{\pgfqpoint{0.027778in}{0.000000in}}{%
\pgfpathmoveto{\pgfqpoint{0.000000in}{0.000000in}}%
\pgfpathlineto{\pgfqpoint{0.027778in}{0.000000in}}%
\pgfusepath{stroke,fill}%
}%
\begin{pgfscope}%
\pgfsys@transformshift{0.472301in}{1.190652in}%
\pgfsys@useobject{currentmarker}{}%
\end{pgfscope}%
\end{pgfscope}%
\begin{pgfscope}%
\pgfsetbuttcap%
\pgfsetroundjoin%
\definecolor{currentfill}{rgb}{0.000000,0.000000,0.000000}%
\pgfsetfillcolor{currentfill}%
\pgfsetlinewidth{0.301125pt}%
\definecolor{currentstroke}{rgb}{0.000000,0.000000,0.000000}%
\pgfsetstrokecolor{currentstroke}%
\pgfsetdash{}{0pt}%
\pgfsys@defobject{currentmarker}{\pgfqpoint{0.000000in}{0.000000in}}{\pgfqpoint{0.027778in}{0.000000in}}{%
\pgfpathmoveto{\pgfqpoint{0.000000in}{0.000000in}}%
\pgfpathlineto{\pgfqpoint{0.027778in}{0.000000in}}%
\pgfusepath{stroke,fill}%
}%
\begin{pgfscope}%
\pgfsys@transformshift{0.472301in}{1.343013in}%
\pgfsys@useobject{currentmarker}{}%
\end{pgfscope}%
\end{pgfscope}%
\begin{pgfscope}%
\pgfsetbuttcap%
\pgfsetroundjoin%
\definecolor{currentfill}{rgb}{0.000000,0.000000,0.000000}%
\pgfsetfillcolor{currentfill}%
\pgfsetlinewidth{0.301125pt}%
\definecolor{currentstroke}{rgb}{0.000000,0.000000,0.000000}%
\pgfsetstrokecolor{currentstroke}%
\pgfsetdash{}{0pt}%
\pgfsys@defobject{currentmarker}{\pgfqpoint{0.000000in}{0.000000in}}{\pgfqpoint{0.027778in}{0.000000in}}{%
\pgfpathmoveto{\pgfqpoint{0.000000in}{0.000000in}}%
\pgfpathlineto{\pgfqpoint{0.027778in}{0.000000in}}%
\pgfusepath{stroke,fill}%
}%
\begin{pgfscope}%
\pgfsys@transformshift{0.472301in}{1.419193in}%
\pgfsys@useobject{currentmarker}{}%
\end{pgfscope}%
\end{pgfscope}%
\begin{pgfscope}%
\pgfsetbuttcap%
\pgfsetroundjoin%
\definecolor{currentfill}{rgb}{0.000000,0.000000,0.000000}%
\pgfsetfillcolor{currentfill}%
\pgfsetlinewidth{0.301125pt}%
\definecolor{currentstroke}{rgb}{0.000000,0.000000,0.000000}%
\pgfsetstrokecolor{currentstroke}%
\pgfsetdash{}{0pt}%
\pgfsys@defobject{currentmarker}{\pgfqpoint{0.000000in}{0.000000in}}{\pgfqpoint{0.027778in}{0.000000in}}{%
\pgfpathmoveto{\pgfqpoint{0.000000in}{0.000000in}}%
\pgfpathlineto{\pgfqpoint{0.027778in}{0.000000in}}%
\pgfusepath{stroke,fill}%
}%
\begin{pgfscope}%
\pgfsys@transformshift{0.472301in}{1.495374in}%
\pgfsys@useobject{currentmarker}{}%
\end{pgfscope}%
\end{pgfscope}%
\begin{pgfscope}%
\pgfsetbuttcap%
\pgfsetroundjoin%
\definecolor{currentfill}{rgb}{0.000000,0.000000,0.000000}%
\pgfsetfillcolor{currentfill}%
\pgfsetlinewidth{0.301125pt}%
\definecolor{currentstroke}{rgb}{0.000000,0.000000,0.000000}%
\pgfsetstrokecolor{currentstroke}%
\pgfsetdash{}{0pt}%
\pgfsys@defobject{currentmarker}{\pgfqpoint{0.000000in}{0.000000in}}{\pgfqpoint{0.027778in}{0.000000in}}{%
\pgfpathmoveto{\pgfqpoint{0.000000in}{0.000000in}}%
\pgfpathlineto{\pgfqpoint{0.027778in}{0.000000in}}%
\pgfusepath{stroke,fill}%
}%
\begin{pgfscope}%
\pgfsys@transformshift{0.472301in}{1.571554in}%
\pgfsys@useobject{currentmarker}{}%
\end{pgfscope}%
\end{pgfscope}%
\begin{pgfscope}%
\pgfsetbuttcap%
\pgfsetroundjoin%
\definecolor{currentfill}{rgb}{0.000000,0.000000,0.000000}%
\pgfsetfillcolor{currentfill}%
\pgfsetlinewidth{0.301125pt}%
\definecolor{currentstroke}{rgb}{0.000000,0.000000,0.000000}%
\pgfsetstrokecolor{currentstroke}%
\pgfsetdash{}{0pt}%
\pgfsys@defobject{currentmarker}{\pgfqpoint{0.000000in}{0.000000in}}{\pgfqpoint{0.027778in}{0.000000in}}{%
\pgfpathmoveto{\pgfqpoint{0.000000in}{0.000000in}}%
\pgfpathlineto{\pgfqpoint{0.027778in}{0.000000in}}%
\pgfusepath{stroke,fill}%
}%
\begin{pgfscope}%
\pgfsys@transformshift{0.472301in}{1.723915in}%
\pgfsys@useobject{currentmarker}{}%
\end{pgfscope}%
\end{pgfscope}%
\begin{pgfscope}%
\pgfsetbuttcap%
\pgfsetroundjoin%
\definecolor{currentfill}{rgb}{0.000000,0.000000,0.000000}%
\pgfsetfillcolor{currentfill}%
\pgfsetlinewidth{0.301125pt}%
\definecolor{currentstroke}{rgb}{0.000000,0.000000,0.000000}%
\pgfsetstrokecolor{currentstroke}%
\pgfsetdash{}{0pt}%
\pgfsys@defobject{currentmarker}{\pgfqpoint{0.000000in}{0.000000in}}{\pgfqpoint{0.027778in}{0.000000in}}{%
\pgfpathmoveto{\pgfqpoint{0.000000in}{0.000000in}}%
\pgfpathlineto{\pgfqpoint{0.027778in}{0.000000in}}%
\pgfusepath{stroke,fill}%
}%
\begin{pgfscope}%
\pgfsys@transformshift{0.472301in}{1.800095in}%
\pgfsys@useobject{currentmarker}{}%
\end{pgfscope}%
\end{pgfscope}%
\begin{pgfscope}%
\pgfsetbuttcap%
\pgfsetroundjoin%
\definecolor{currentfill}{rgb}{0.000000,0.000000,0.000000}%
\pgfsetfillcolor{currentfill}%
\pgfsetlinewidth{0.301125pt}%
\definecolor{currentstroke}{rgb}{0.000000,0.000000,0.000000}%
\pgfsetstrokecolor{currentstroke}%
\pgfsetdash{}{0pt}%
\pgfsys@defobject{currentmarker}{\pgfqpoint{0.000000in}{0.000000in}}{\pgfqpoint{0.027778in}{0.000000in}}{%
\pgfpathmoveto{\pgfqpoint{0.000000in}{0.000000in}}%
\pgfpathlineto{\pgfqpoint{0.027778in}{0.000000in}}%
\pgfusepath{stroke,fill}%
}%
\begin{pgfscope}%
\pgfsys@transformshift{0.472301in}{1.876275in}%
\pgfsys@useobject{currentmarker}{}%
\end{pgfscope}%
\end{pgfscope}%
\begin{pgfscope}%
\pgfsetbuttcap%
\pgfsetroundjoin%
\definecolor{currentfill}{rgb}{0.000000,0.000000,0.000000}%
\pgfsetfillcolor{currentfill}%
\pgfsetlinewidth{0.301125pt}%
\definecolor{currentstroke}{rgb}{0.000000,0.000000,0.000000}%
\pgfsetstrokecolor{currentstroke}%
\pgfsetdash{}{0pt}%
\pgfsys@defobject{currentmarker}{\pgfqpoint{0.000000in}{0.000000in}}{\pgfqpoint{0.027778in}{0.000000in}}{%
\pgfpathmoveto{\pgfqpoint{0.000000in}{0.000000in}}%
\pgfpathlineto{\pgfqpoint{0.027778in}{0.000000in}}%
\pgfusepath{stroke,fill}%
}%
\begin{pgfscope}%
\pgfsys@transformshift{0.472301in}{1.952456in}%
\pgfsys@useobject{currentmarker}{}%
\end{pgfscope}%
\end{pgfscope}%
\begin{pgfscope}%
\pgfsetbuttcap%
\pgfsetroundjoin%
\definecolor{currentfill}{rgb}{0.000000,0.000000,0.000000}%
\pgfsetfillcolor{currentfill}%
\pgfsetlinewidth{0.301125pt}%
\definecolor{currentstroke}{rgb}{0.000000,0.000000,0.000000}%
\pgfsetstrokecolor{currentstroke}%
\pgfsetdash{}{0pt}%
\pgfsys@defobject{currentmarker}{\pgfqpoint{0.000000in}{0.000000in}}{\pgfqpoint{0.027778in}{0.000000in}}{%
\pgfpathmoveto{\pgfqpoint{0.000000in}{0.000000in}}%
\pgfpathlineto{\pgfqpoint{0.027778in}{0.000000in}}%
\pgfusepath{stroke,fill}%
}%
\begin{pgfscope}%
\pgfsys@transformshift{0.472301in}{2.104816in}%
\pgfsys@useobject{currentmarker}{}%
\end{pgfscope}%
\end{pgfscope}%
\begin{pgfscope}%
\pgfsetbuttcap%
\pgfsetroundjoin%
\definecolor{currentfill}{rgb}{0.000000,0.000000,0.000000}%
\pgfsetfillcolor{currentfill}%
\pgfsetlinewidth{0.301125pt}%
\definecolor{currentstroke}{rgb}{0.000000,0.000000,0.000000}%
\pgfsetstrokecolor{currentstroke}%
\pgfsetdash{}{0pt}%
\pgfsys@defobject{currentmarker}{\pgfqpoint{0.000000in}{0.000000in}}{\pgfqpoint{0.027778in}{0.000000in}}{%
\pgfpathmoveto{\pgfqpoint{0.000000in}{0.000000in}}%
\pgfpathlineto{\pgfqpoint{0.027778in}{0.000000in}}%
\pgfusepath{stroke,fill}%
}%
\begin{pgfscope}%
\pgfsys@transformshift{0.472301in}{2.180997in}%
\pgfsys@useobject{currentmarker}{}%
\end{pgfscope}%
\end{pgfscope}%
\begin{pgfscope}%
\pgfsetbuttcap%
\pgfsetroundjoin%
\definecolor{currentfill}{rgb}{0.000000,0.000000,0.000000}%
\pgfsetfillcolor{currentfill}%
\pgfsetlinewidth{0.301125pt}%
\definecolor{currentstroke}{rgb}{0.000000,0.000000,0.000000}%
\pgfsetstrokecolor{currentstroke}%
\pgfsetdash{}{0pt}%
\pgfsys@defobject{currentmarker}{\pgfqpoint{0.000000in}{0.000000in}}{\pgfqpoint{0.027778in}{0.000000in}}{%
\pgfpathmoveto{\pgfqpoint{0.000000in}{0.000000in}}%
\pgfpathlineto{\pgfqpoint{0.027778in}{0.000000in}}%
\pgfusepath{stroke,fill}%
}%
\begin{pgfscope}%
\pgfsys@transformshift{0.472301in}{2.257177in}%
\pgfsys@useobject{currentmarker}{}%
\end{pgfscope}%
\end{pgfscope}%
\begin{pgfscope}%
\pgfsetbuttcap%
\pgfsetroundjoin%
\definecolor{currentfill}{rgb}{0.000000,0.000000,0.000000}%
\pgfsetfillcolor{currentfill}%
\pgfsetlinewidth{0.301125pt}%
\definecolor{currentstroke}{rgb}{0.000000,0.000000,0.000000}%
\pgfsetstrokecolor{currentstroke}%
\pgfsetdash{}{0pt}%
\pgfsys@defobject{currentmarker}{\pgfqpoint{0.000000in}{0.000000in}}{\pgfqpoint{0.027778in}{0.000000in}}{%
\pgfpathmoveto{\pgfqpoint{0.000000in}{0.000000in}}%
\pgfpathlineto{\pgfqpoint{0.027778in}{0.000000in}}%
\pgfusepath{stroke,fill}%
}%
\begin{pgfscope}%
\pgfsys@transformshift{0.472301in}{2.333357in}%
\pgfsys@useobject{currentmarker}{}%
\end{pgfscope}%
\end{pgfscope}%
\begin{pgfscope}%
\pgfsetbuttcap%
\pgfsetroundjoin%
\definecolor{currentfill}{rgb}{0.000000,0.000000,0.000000}%
\pgfsetfillcolor{currentfill}%
\pgfsetlinewidth{0.301125pt}%
\definecolor{currentstroke}{rgb}{0.000000,0.000000,0.000000}%
\pgfsetstrokecolor{currentstroke}%
\pgfsetdash{}{0pt}%
\pgfsys@defobject{currentmarker}{\pgfqpoint{0.000000in}{0.000000in}}{\pgfqpoint{0.027778in}{0.000000in}}{%
\pgfpathmoveto{\pgfqpoint{0.000000in}{0.000000in}}%
\pgfpathlineto{\pgfqpoint{0.027778in}{0.000000in}}%
\pgfusepath{stroke,fill}%
}%
\begin{pgfscope}%
\pgfsys@transformshift{0.472301in}{2.485718in}%
\pgfsys@useobject{currentmarker}{}%
\end{pgfscope}%
\end{pgfscope}%
\begin{pgfscope}%
\pgfsetbuttcap%
\pgfsetroundjoin%
\definecolor{currentfill}{rgb}{0.000000,0.000000,0.000000}%
\pgfsetfillcolor{currentfill}%
\pgfsetlinewidth{0.301125pt}%
\definecolor{currentstroke}{rgb}{0.000000,0.000000,0.000000}%
\pgfsetstrokecolor{currentstroke}%
\pgfsetdash{}{0pt}%
\pgfsys@defobject{currentmarker}{\pgfqpoint{0.000000in}{0.000000in}}{\pgfqpoint{0.027778in}{0.000000in}}{%
\pgfpathmoveto{\pgfqpoint{0.000000in}{0.000000in}}%
\pgfpathlineto{\pgfqpoint{0.027778in}{0.000000in}}%
\pgfusepath{stroke,fill}%
}%
\begin{pgfscope}%
\pgfsys@transformshift{0.472301in}{2.561898in}%
\pgfsys@useobject{currentmarker}{}%
\end{pgfscope}%
\end{pgfscope}%
\begin{pgfscope}%
\pgfsetbuttcap%
\pgfsetroundjoin%
\definecolor{currentfill}{rgb}{0.000000,0.000000,0.000000}%
\pgfsetfillcolor{currentfill}%
\pgfsetlinewidth{0.301125pt}%
\definecolor{currentstroke}{rgb}{0.000000,0.000000,0.000000}%
\pgfsetstrokecolor{currentstroke}%
\pgfsetdash{}{0pt}%
\pgfsys@defobject{currentmarker}{\pgfqpoint{0.000000in}{0.000000in}}{\pgfqpoint{0.027778in}{0.000000in}}{%
\pgfpathmoveto{\pgfqpoint{0.000000in}{0.000000in}}%
\pgfpathlineto{\pgfqpoint{0.027778in}{0.000000in}}%
\pgfusepath{stroke,fill}%
}%
\begin{pgfscope}%
\pgfsys@transformshift{0.472301in}{2.638079in}%
\pgfsys@useobject{currentmarker}{}%
\end{pgfscope}%
\end{pgfscope}%
\begin{pgfscope}%
\pgfsetbuttcap%
\pgfsetroundjoin%
\definecolor{currentfill}{rgb}{0.000000,0.000000,0.000000}%
\pgfsetfillcolor{currentfill}%
\pgfsetlinewidth{0.301125pt}%
\definecolor{currentstroke}{rgb}{0.000000,0.000000,0.000000}%
\pgfsetstrokecolor{currentstroke}%
\pgfsetdash{}{0pt}%
\pgfsys@defobject{currentmarker}{\pgfqpoint{0.000000in}{0.000000in}}{\pgfqpoint{0.027778in}{0.000000in}}{%
\pgfpathmoveto{\pgfqpoint{0.000000in}{0.000000in}}%
\pgfpathlineto{\pgfqpoint{0.027778in}{0.000000in}}%
\pgfusepath{stroke,fill}%
}%
\begin{pgfscope}%
\pgfsys@transformshift{0.472301in}{2.714259in}%
\pgfsys@useobject{currentmarker}{}%
\end{pgfscope}%
\end{pgfscope}%
\begin{pgfscope}%
\pgfsetbuttcap%
\pgfsetroundjoin%
\definecolor{currentfill}{rgb}{0.000000,0.000000,0.000000}%
\pgfsetfillcolor{currentfill}%
\pgfsetlinewidth{0.301125pt}%
\definecolor{currentstroke}{rgb}{0.000000,0.000000,0.000000}%
\pgfsetstrokecolor{currentstroke}%
\pgfsetdash{}{0pt}%
\pgfsys@defobject{currentmarker}{\pgfqpoint{0.000000in}{0.000000in}}{\pgfqpoint{0.027778in}{0.000000in}}{%
\pgfpathmoveto{\pgfqpoint{0.000000in}{0.000000in}}%
\pgfpathlineto{\pgfqpoint{0.027778in}{0.000000in}}%
\pgfusepath{stroke,fill}%
}%
\begin{pgfscope}%
\pgfsys@transformshift{0.472301in}{2.866620in}%
\pgfsys@useobject{currentmarker}{}%
\end{pgfscope}%
\end{pgfscope}%
\begin{pgfscope}%
\pgfsetbuttcap%
\pgfsetroundjoin%
\definecolor{currentfill}{rgb}{0.000000,0.000000,0.000000}%
\pgfsetfillcolor{currentfill}%
\pgfsetlinewidth{0.301125pt}%
\definecolor{currentstroke}{rgb}{0.000000,0.000000,0.000000}%
\pgfsetstrokecolor{currentstroke}%
\pgfsetdash{}{0pt}%
\pgfsys@defobject{currentmarker}{\pgfqpoint{0.000000in}{0.000000in}}{\pgfqpoint{0.027778in}{0.000000in}}{%
\pgfpathmoveto{\pgfqpoint{0.000000in}{0.000000in}}%
\pgfpathlineto{\pgfqpoint{0.027778in}{0.000000in}}%
\pgfusepath{stroke,fill}%
}%
\begin{pgfscope}%
\pgfsys@transformshift{0.472301in}{2.942800in}%
\pgfsys@useobject{currentmarker}{}%
\end{pgfscope}%
\end{pgfscope}%
\begin{pgfscope}%
\pgfsetbuttcap%
\pgfsetroundjoin%
\definecolor{currentfill}{rgb}{0.000000,0.000000,0.000000}%
\pgfsetfillcolor{currentfill}%
\pgfsetlinewidth{0.301125pt}%
\definecolor{currentstroke}{rgb}{0.000000,0.000000,0.000000}%
\pgfsetstrokecolor{currentstroke}%
\pgfsetdash{}{0pt}%
\pgfsys@defobject{currentmarker}{\pgfqpoint{0.000000in}{0.000000in}}{\pgfqpoint{0.027778in}{0.000000in}}{%
\pgfpathmoveto{\pgfqpoint{0.000000in}{0.000000in}}%
\pgfpathlineto{\pgfqpoint{0.027778in}{0.000000in}}%
\pgfusepath{stroke,fill}%
}%
\begin{pgfscope}%
\pgfsys@transformshift{0.472301in}{3.018980in}%
\pgfsys@useobject{currentmarker}{}%
\end{pgfscope}%
\end{pgfscope}%
\begin{pgfscope}%
\pgfsetbuttcap%
\pgfsetroundjoin%
\definecolor{currentfill}{rgb}{0.000000,0.000000,0.000000}%
\pgfsetfillcolor{currentfill}%
\pgfsetlinewidth{0.301125pt}%
\definecolor{currentstroke}{rgb}{0.000000,0.000000,0.000000}%
\pgfsetstrokecolor{currentstroke}%
\pgfsetdash{}{0pt}%
\pgfsys@defobject{currentmarker}{\pgfqpoint{0.000000in}{0.000000in}}{\pgfqpoint{0.027778in}{0.000000in}}{%
\pgfpathmoveto{\pgfqpoint{0.000000in}{0.000000in}}%
\pgfpathlineto{\pgfqpoint{0.027778in}{0.000000in}}%
\pgfusepath{stroke,fill}%
}%
\begin{pgfscope}%
\pgfsys@transformshift{0.472301in}{3.095161in}%
\pgfsys@useobject{currentmarker}{}%
\end{pgfscope}%
\end{pgfscope}%
\begin{pgfscope}%
\pgfsetbuttcap%
\pgfsetroundjoin%
\definecolor{currentfill}{rgb}{0.000000,0.000000,0.000000}%
\pgfsetfillcolor{currentfill}%
\pgfsetlinewidth{0.301125pt}%
\definecolor{currentstroke}{rgb}{0.000000,0.000000,0.000000}%
\pgfsetstrokecolor{currentstroke}%
\pgfsetdash{}{0pt}%
\pgfsys@defobject{currentmarker}{\pgfqpoint{0.000000in}{0.000000in}}{\pgfqpoint{0.027778in}{0.000000in}}{%
\pgfpathmoveto{\pgfqpoint{0.000000in}{0.000000in}}%
\pgfpathlineto{\pgfqpoint{0.027778in}{0.000000in}}%
\pgfusepath{stroke,fill}%
}%
\begin{pgfscope}%
\pgfsys@transformshift{0.472301in}{3.247521in}%
\pgfsys@useobject{currentmarker}{}%
\end{pgfscope}%
\end{pgfscope}%
\begin{pgfscope}%
\pgfsetbuttcap%
\pgfsetroundjoin%
\definecolor{currentfill}{rgb}{0.000000,0.000000,0.000000}%
\pgfsetfillcolor{currentfill}%
\pgfsetlinewidth{0.301125pt}%
\definecolor{currentstroke}{rgb}{0.000000,0.000000,0.000000}%
\pgfsetstrokecolor{currentstroke}%
\pgfsetdash{}{0pt}%
\pgfsys@defobject{currentmarker}{\pgfqpoint{0.000000in}{0.000000in}}{\pgfqpoint{0.027778in}{0.000000in}}{%
\pgfpathmoveto{\pgfqpoint{0.000000in}{0.000000in}}%
\pgfpathlineto{\pgfqpoint{0.027778in}{0.000000in}}%
\pgfusepath{stroke,fill}%
}%
\begin{pgfscope}%
\pgfsys@transformshift{0.472301in}{3.323702in}%
\pgfsys@useobject{currentmarker}{}%
\end{pgfscope}%
\end{pgfscope}%
\begin{pgfscope}%
\pgfsetbuttcap%
\pgfsetroundjoin%
\definecolor{currentfill}{rgb}{0.000000,0.000000,0.000000}%
\pgfsetfillcolor{currentfill}%
\pgfsetlinewidth{0.301125pt}%
\definecolor{currentstroke}{rgb}{0.000000,0.000000,0.000000}%
\pgfsetstrokecolor{currentstroke}%
\pgfsetdash{}{0pt}%
\pgfsys@defobject{currentmarker}{\pgfqpoint{0.000000in}{0.000000in}}{\pgfqpoint{0.027778in}{0.000000in}}{%
\pgfpathmoveto{\pgfqpoint{0.000000in}{0.000000in}}%
\pgfpathlineto{\pgfqpoint{0.027778in}{0.000000in}}%
\pgfusepath{stroke,fill}%
}%
\begin{pgfscope}%
\pgfsys@transformshift{0.472301in}{3.399882in}%
\pgfsys@useobject{currentmarker}{}%
\end{pgfscope}%
\end{pgfscope}%
\begin{pgfscope}%
\pgfsetbuttcap%
\pgfsetroundjoin%
\definecolor{currentfill}{rgb}{0.000000,0.000000,0.000000}%
\pgfsetfillcolor{currentfill}%
\pgfsetlinewidth{0.301125pt}%
\definecolor{currentstroke}{rgb}{0.000000,0.000000,0.000000}%
\pgfsetstrokecolor{currentstroke}%
\pgfsetdash{}{0pt}%
\pgfsys@defobject{currentmarker}{\pgfqpoint{0.000000in}{0.000000in}}{\pgfqpoint{0.027778in}{0.000000in}}{%
\pgfpathmoveto{\pgfqpoint{0.000000in}{0.000000in}}%
\pgfpathlineto{\pgfqpoint{0.027778in}{0.000000in}}%
\pgfusepath{stroke,fill}%
}%
\begin{pgfscope}%
\pgfsys@transformshift{0.472301in}{3.476062in}%
\pgfsys@useobject{currentmarker}{}%
\end{pgfscope}%
\end{pgfscope}%
\begin{pgfscope}%
\pgfsetbuttcap%
\pgfsetroundjoin%
\definecolor{currentfill}{rgb}{0.000000,0.000000,0.000000}%
\pgfsetfillcolor{currentfill}%
\pgfsetlinewidth{0.301125pt}%
\definecolor{currentstroke}{rgb}{0.000000,0.000000,0.000000}%
\pgfsetstrokecolor{currentstroke}%
\pgfsetdash{}{0pt}%
\pgfsys@defobject{currentmarker}{\pgfqpoint{0.000000in}{0.000000in}}{\pgfqpoint{0.027778in}{0.000000in}}{%
\pgfpathmoveto{\pgfqpoint{0.000000in}{0.000000in}}%
\pgfpathlineto{\pgfqpoint{0.027778in}{0.000000in}}%
\pgfusepath{stroke,fill}%
}%
\begin{pgfscope}%
\pgfsys@transformshift{0.472301in}{3.628423in}%
\pgfsys@useobject{currentmarker}{}%
\end{pgfscope}%
\end{pgfscope}%
\begin{pgfscope}%
\pgfsetbuttcap%
\pgfsetroundjoin%
\definecolor{currentfill}{rgb}{0.000000,0.000000,0.000000}%
\pgfsetfillcolor{currentfill}%
\pgfsetlinewidth{0.301125pt}%
\definecolor{currentstroke}{rgb}{0.000000,0.000000,0.000000}%
\pgfsetstrokecolor{currentstroke}%
\pgfsetdash{}{0pt}%
\pgfsys@defobject{currentmarker}{\pgfqpoint{0.000000in}{0.000000in}}{\pgfqpoint{0.027778in}{0.000000in}}{%
\pgfpathmoveto{\pgfqpoint{0.000000in}{0.000000in}}%
\pgfpathlineto{\pgfqpoint{0.027778in}{0.000000in}}%
\pgfusepath{stroke,fill}%
}%
\begin{pgfscope}%
\pgfsys@transformshift{0.472301in}{3.704603in}%
\pgfsys@useobject{currentmarker}{}%
\end{pgfscope}%
\end{pgfscope}%
\begin{pgfscope}%
\definecolor{textcolor}{rgb}{0.000000,0.000000,0.000000}%
\pgfsetstrokecolor{textcolor}%
\pgfsetfillcolor{textcolor}%
\pgftext[x=0.140972in,y=2.028636in,,bottom,rotate=90.000000]{\color{textcolor}{\sffamily\fontsize{9.000000}{10.800000}\selectfont\catcode`\^=\active\def^{\ifmmode\sp\else\^{}\fi}\catcode`\%=\active\def%{\%}$y / M$}}%
\end{pgfscope}%
\begin{pgfscope}%
\pgfpathrectangle{\pgfqpoint{0.472301in}{0.352669in}}{\pgfqpoint{3.351934in}{3.351934in}}%
\pgfusepath{clip}%
\pgfsetbuttcap%
\pgfsetroundjoin%
\pgfsetlinewidth{0.602250pt}%
\definecolor{currentstroke}{rgb}{0.431373,0.423529,0.462745}%
\pgfsetstrokecolor{currentstroke}%
\pgfsetstrokeopacity{0.500000}%
\pgfsetdash{{2.220000pt}{0.960000pt}}{0.000000pt}%
\pgfpathmoveto{\pgfqpoint{2.300628in}{2.028636in}}%
\pgfpathlineto{\pgfqpoint{2.299788in}{2.044615in}}%
\pgfpathlineto{\pgfqpoint{2.297277in}{2.060418in}}%
\pgfpathlineto{\pgfqpoint{2.293122in}{2.075870in}}%
\pgfpathlineto{\pgfqpoint{2.287369in}{2.090802in}}%
\pgfpathlineto{\pgfqpoint{2.280082in}{2.105047in}}%
\pgfpathlineto{\pgfqpoint{2.271342in}{2.118450in}}%
\pgfpathlineto{\pgfqpoint{2.261244in}{2.130862in}}%
\pgfpathlineto{\pgfqpoint{2.249899in}{2.142147in}}%
\pgfpathlineto{\pgfqpoint{2.237434in}{2.152180in}}%
\pgfpathlineto{\pgfqpoint{2.223986in}{2.160850in}}%
\pgfpathlineto{\pgfqpoint{2.209702in}{2.168062in}}%
\pgfpathlineto{\pgfqpoint{2.194740in}{2.173736in}}%
\pgfpathlineto{\pgfqpoint{2.179266in}{2.177810in}}%
\pgfpathlineto{\pgfqpoint{2.163451in}{2.180238in}}%
\pgfpathlineto{\pgfqpoint{2.147467in}{2.180995in}}%
\pgfpathlineto{\pgfqpoint{2.131493in}{2.180070in}}%
\pgfpathlineto{\pgfqpoint{2.115703in}{2.177476in}}%
\pgfpathlineto{\pgfqpoint{2.100273in}{2.173240in}}%
\pgfpathlineto{\pgfqpoint{2.085372in}{2.167409in}}%
\pgfpathlineto{\pgfqpoint{2.071165in}{2.160047in}}%
\pgfpathlineto{\pgfqpoint{2.057808in}{2.151236in}}%
\pgfpathlineto{\pgfqpoint{2.045449in}{2.141073in}}%
\pgfpathlineto{\pgfqpoint{2.034224in}{2.129670in}}%
\pgfpathlineto{\pgfqpoint{2.024257in}{2.117152in}}%
\pgfpathlineto{\pgfqpoint{2.015658in}{2.103658in}}%
\pgfpathlineto{\pgfqpoint{2.008521in}{2.089337in}}%
\pgfpathlineto{\pgfqpoint{2.002925in}{2.074346in}}%
\pgfpathlineto{\pgfqpoint{1.998933in}{2.058851in}}%
\pgfpathlineto{\pgfqpoint{1.996588in}{2.043022in}}%
\pgfpathlineto{\pgfqpoint{1.995916in}{2.027035in}}%
\pgfpathlineto{\pgfqpoint{1.996924in}{2.011066in}}%
\pgfpathlineto{\pgfqpoint{1.999601in}{1.995290in}}%
\pgfpathlineto{\pgfqpoint{2.003918in}{1.979882in}}%
\pgfpathlineto{\pgfqpoint{2.009827in}{1.965012in}}%
\pgfpathlineto{\pgfqpoint{2.017263in}{1.950844in}}%
\pgfpathlineto{\pgfqpoint{2.026144in}{1.937534in}}%
\pgfpathlineto{\pgfqpoint{2.036372in}{1.925228in}}%
\pgfpathlineto{\pgfqpoint{2.047834in}{1.914063in}}%
\pgfpathlineto{\pgfqpoint{2.060404in}{1.904162in}}%
\pgfpathlineto{\pgfqpoint{2.073943in}{1.895634in}}%
\pgfpathlineto{\pgfqpoint{2.088302in}{1.888572in}}%
\pgfpathlineto{\pgfqpoint{2.103322in}{1.883056in}}%
\pgfpathlineto{\pgfqpoint{2.118838in}{1.879145in}}%
\pgfpathlineto{\pgfqpoint{2.134679in}{1.876883in}}%
\pgfpathlineto{\pgfqpoint{2.150669in}{1.876294in}}%
\pgfpathlineto{\pgfqpoint{2.166633in}{1.877386in}}%
\pgfpathlineto{\pgfqpoint{2.182394in}{1.880146in}}%
\pgfpathlineto{\pgfqpoint{2.197779in}{1.884544in}}%
\pgfpathlineto{\pgfqpoint{2.212618in}{1.890532in}}%
\pgfpathlineto{\pgfqpoint{2.226747in}{1.898042in}}%
\pgfpathlineto{\pgfqpoint{2.240010in}{1.906993in}}%
\pgfpathlineto{\pgfqpoint{2.252262in}{1.917285in}}%
\pgfpathlineto{\pgfqpoint{2.263367in}{1.928806in}}%
\pgfpathlineto{\pgfqpoint{2.273202in}{1.941428in}}%
\pgfpathlineto{\pgfqpoint{2.281659in}{1.955011in}}%
\pgfpathlineto{\pgfqpoint{2.288645in}{1.969407in}}%
\pgfpathlineto{\pgfqpoint{2.294082in}{1.984456in}}%
\pgfpathlineto{\pgfqpoint{2.297912in}{1.999992in}}%
\pgfpathlineto{\pgfqpoint{2.300091in}{2.015844in}}%
\pgfpathlineto{\pgfqpoint{2.300628in}{2.028636in}}%
\pgfpathlineto{\pgfqpoint{2.300628in}{2.028636in}}%
\pgfusepath{stroke}%
\end{pgfscope}%
\begin{pgfscope}%
\pgfpathrectangle{\pgfqpoint{0.472301in}{0.352669in}}{\pgfqpoint{3.351934in}{3.351934in}}%
\pgfusepath{clip}%
\pgfsetbuttcap%
\pgfsetroundjoin%
\pgfsetlinewidth{0.602250pt}%
\definecolor{currentstroke}{rgb}{0.431373,0.423529,0.462745}%
\pgfsetstrokecolor{currentstroke}%
\pgfsetstrokeopacity{0.500000}%
\pgfsetdash{{0.600000pt}{0.990000pt}}{0.000000pt}%
\pgfpathmoveto{\pgfqpoint{2.605350in}{2.028636in}}%
\pgfpathlineto{\pgfqpoint{2.604442in}{2.057432in}}%
\pgfpathlineto{\pgfqpoint{2.601721in}{2.086114in}}%
\pgfpathlineto{\pgfqpoint{2.597199in}{2.114568in}}%
\pgfpathlineto{\pgfqpoint{2.590894in}{2.142680in}}%
\pgfpathlineto{\pgfqpoint{2.582830in}{2.170339in}}%
\pgfpathlineto{\pgfqpoint{2.573039in}{2.197435in}}%
\pgfpathlineto{\pgfqpoint{2.561561in}{2.223860in}}%
\pgfpathlineto{\pgfqpoint{2.548441in}{2.249510in}}%
\pgfpathlineto{\pgfqpoint{2.533731in}{2.274282in}}%
\pgfpathlineto{\pgfqpoint{2.517489in}{2.298078in}}%
\pgfpathlineto{\pgfqpoint{2.499781in}{2.320804in}}%
\pgfpathlineto{\pgfqpoint{2.480676in}{2.342369in}}%
\pgfpathlineto{\pgfqpoint{2.460251in}{2.362688in}}%
\pgfpathlineto{\pgfqpoint{2.438585in}{2.381679in}}%
\pgfpathlineto{\pgfqpoint{2.415767in}{2.399268in}}%
\pgfpathlineto{\pgfqpoint{2.391886in}{2.415384in}}%
\pgfpathlineto{\pgfqpoint{2.367036in}{2.429964in}}%
\pgfpathlineto{\pgfqpoint{2.341318in}{2.442949in}}%
\pgfpathlineto{\pgfqpoint{2.314833in}{2.454288in}}%
\pgfpathlineto{\pgfqpoint{2.287686in}{2.463937in}}%
\pgfpathlineto{\pgfqpoint{2.259985in}{2.471855in}}%
\pgfpathlineto{\pgfqpoint{2.231840in}{2.478013in}}%
\pgfpathlineto{\pgfqpoint{2.203363in}{2.482385in}}%
\pgfpathlineto{\pgfqpoint{2.174667in}{2.484955in}}%
\pgfpathlineto{\pgfqpoint{2.145867in}{2.485712in}}%
\pgfpathlineto{\pgfqpoint{2.117075in}{2.484652in}}%
\pgfpathlineto{\pgfqpoint{2.088408in}{2.481781in}}%
\pgfpathlineto{\pgfqpoint{2.059979in}{2.477110in}}%
\pgfpathlineto{\pgfqpoint{2.031900in}{2.470657in}}%
\pgfpathlineto{\pgfqpoint{2.004284in}{2.462448in}}%
\pgfpathlineto{\pgfqpoint{1.977240in}{2.452515in}}%
\pgfpathlineto{\pgfqpoint{1.950875in}{2.440898in}}%
\pgfpathlineto{\pgfqpoint{1.925295in}{2.427643in}}%
\pgfpathlineto{\pgfqpoint{1.900600in}{2.412803in}}%
\pgfpathlineto{\pgfqpoint{1.876889in}{2.396437in}}%
\pgfpathlineto{\pgfqpoint{1.854257in}{2.378610in}}%
\pgfpathlineto{\pgfqpoint{1.832792in}{2.359392in}}%
\pgfpathlineto{\pgfqpoint{1.812581in}{2.338859in}}%
\pgfpathlineto{\pgfqpoint{1.793704in}{2.317095in}}%
\pgfpathlineto{\pgfqpoint{1.776235in}{2.294184in}}%
\pgfpathlineto{\pgfqpoint{1.760245in}{2.270219in}}%
\pgfpathlineto{\pgfqpoint{1.745796in}{2.245293in}}%
\pgfpathlineto{\pgfqpoint{1.732946in}{2.219507in}}%
\pgfpathlineto{\pgfqpoint{1.721746in}{2.192962in}}%
\pgfpathlineto{\pgfqpoint{1.712241in}{2.165765in}}%
\pgfpathlineto{\pgfqpoint{1.704468in}{2.138023in}}%
\pgfpathlineto{\pgfqpoint{1.698458in}{2.109846in}}%
\pgfpathlineto{\pgfqpoint{1.694235in}{2.081347in}}%
\pgfpathlineto{\pgfqpoint{1.691816in}{2.052638in}}%
\pgfpathlineto{\pgfqpoint{1.691211in}{2.023834in}}%
\pgfpathlineto{\pgfqpoint{1.692422in}{1.995048in}}%
\pgfpathlineto{\pgfqpoint{1.695443in}{1.966397in}}%
\pgfpathlineto{\pgfqpoint{1.700264in}{1.937992in}}%
\pgfpathlineto{\pgfqpoint{1.706864in}{1.909948in}}%
\pgfpathlineto{\pgfqpoint{1.715218in}{1.882375in}}%
\pgfpathlineto{\pgfqpoint{1.725293in}{1.855383in}}%
\pgfpathlineto{\pgfqpoint{1.737048in}{1.829080in}}%
\pgfpathlineto{\pgfqpoint{1.750437in}{1.803570in}}%
\pgfpathlineto{\pgfqpoint{1.765407in}{1.778953in}}%
\pgfpathlineto{\pgfqpoint{1.781897in}{1.755329in}}%
\pgfpathlineto{\pgfqpoint{1.799844in}{1.732790in}}%
\pgfpathlineto{\pgfqpoint{1.819174in}{1.711427in}}%
\pgfpathlineto{\pgfqpoint{1.839812in}{1.691325in}}%
\pgfpathlineto{\pgfqpoint{1.861676in}{1.672562in}}%
\pgfpathlineto{\pgfqpoint{1.884678in}{1.655214in}}%
\pgfpathlineto{\pgfqpoint{1.908727in}{1.639349in}}%
\pgfpathlineto{\pgfqpoint{1.933728in}{1.625031in}}%
\pgfpathlineto{\pgfqpoint{1.959581in}{1.612317in}}%
\pgfpathlineto{\pgfqpoint{1.986184in}{1.601257in}}%
\pgfpathlineto{\pgfqpoint{2.013431in}{1.591895in}}%
\pgfpathlineto{\pgfqpoint{2.041214in}{1.584267in}}%
\pgfpathlineto{\pgfqpoint{2.069422in}{1.578406in}}%
\pgfpathlineto{\pgfqpoint{2.097943in}{1.574333in}}%
\pgfpathlineto{\pgfqpoint{2.126664in}{1.572065in}}%
\pgfpathlineto{\pgfqpoint{2.155471in}{1.571611in}}%
\pgfpathlineto{\pgfqpoint{2.184250in}{1.572972in}}%
\pgfpathlineto{\pgfqpoint{2.212885in}{1.576145in}}%
\pgfpathlineto{\pgfqpoint{2.241264in}{1.581114in}}%
\pgfpathlineto{\pgfqpoint{2.269273in}{1.587862in}}%
\pgfpathlineto{\pgfqpoint{2.296802in}{1.596361in}}%
\pgfpathlineto{\pgfqpoint{2.323740in}{1.606577in}}%
\pgfpathlineto{\pgfqpoint{2.349981in}{1.618471in}}%
\pgfpathlineto{\pgfqpoint{2.375421in}{1.631993in}}%
\pgfpathlineto{\pgfqpoint{2.399958in}{1.647092in}}%
\pgfpathlineto{\pgfqpoint{2.423496in}{1.663707in}}%
\pgfpathlineto{\pgfqpoint{2.445940in}{1.681771in}}%
\pgfpathlineto{\pgfqpoint{2.467201in}{1.701213in}}%
\pgfpathlineto{\pgfqpoint{2.487195in}{1.721957in}}%
\pgfpathlineto{\pgfqpoint{2.505843in}{1.743918in}}%
\pgfpathlineto{\pgfqpoint{2.523070in}{1.767011in}}%
\pgfpathlineto{\pgfqpoint{2.538807in}{1.791144in}}%
\pgfpathlineto{\pgfqpoint{2.552994in}{1.816219in}}%
\pgfpathlineto{\pgfqpoint{2.565572in}{1.842139in}}%
\pgfpathlineto{\pgfqpoint{2.576492in}{1.868800in}}%
\pgfpathlineto{\pgfqpoint{2.585711in}{1.896096in}}%
\pgfpathlineto{\pgfqpoint{2.593193in}{1.923918in}}%
\pgfpathlineto{\pgfqpoint{2.598906in}{1.952156in}}%
\pgfpathlineto{\pgfqpoint{2.602829in}{1.980699in}}%
\pgfpathlineto{\pgfqpoint{2.604946in}{2.009431in}}%
\pgfpathlineto{\pgfqpoint{2.605350in}{2.028636in}}%
\pgfpathlineto{\pgfqpoint{2.605350in}{2.028636in}}%
\pgfusepath{stroke}%
\end{pgfscope}%
\begin{pgfscope}%
\pgfpathrectangle{\pgfqpoint{0.472301in}{0.352669in}}{\pgfqpoint{3.351934in}{3.351934in}}%
\pgfusepath{clip}%
\pgfsetrectcap%
\pgfsetroundjoin%
\pgfsetlinewidth{0.803000pt}%
\definecolor{currentstroke}{rgb}{0.768627,0.556863,0.211765}%
\pgfsetstrokecolor{currentstroke}%
\pgfsetdash{}{0pt}%
\pgfpathmoveto{\pgfqpoint{3.474849in}{2.028636in}}%
\pgfpathlineto{\pgfqpoint{3.474212in}{2.065225in}}%
\pgfpathlineto{\pgfqpoint{3.471782in}{2.108917in}}%
\pgfpathlineto{\pgfqpoint{3.467534in}{2.152498in}}%
\pgfpathlineto{\pgfqpoint{3.461471in}{2.195908in}}%
\pgfpathlineto{\pgfqpoint{3.453596in}{2.239088in}}%
\pgfpathlineto{\pgfqpoint{3.443911in}{2.281976in}}%
\pgfpathlineto{\pgfqpoint{3.432422in}{2.324511in}}%
\pgfpathlineto{\pgfqpoint{3.419134in}{2.366631in}}%
\pgfpathlineto{\pgfqpoint{3.407212in}{2.399987in}}%
\pgfpathlineto{\pgfqpoint{3.394146in}{2.433006in}}%
\pgfpathlineto{\pgfqpoint{3.379941in}{2.465655in}}%
\pgfpathlineto{\pgfqpoint{3.364601in}{2.497901in}}%
\pgfpathlineto{\pgfqpoint{3.348132in}{2.529711in}}%
\pgfpathlineto{\pgfqpoint{3.330538in}{2.561051in}}%
\pgfpathlineto{\pgfqpoint{3.311826in}{2.591889in}}%
\pgfpathlineto{\pgfqpoint{3.292003in}{2.622188in}}%
\pgfpathlineto{\pgfqpoint{3.271076in}{2.651914in}}%
\pgfpathlineto{\pgfqpoint{3.249054in}{2.681032in}}%
\pgfpathlineto{\pgfqpoint{3.225944in}{2.709505in}}%
\pgfpathlineto{\pgfqpoint{3.201756in}{2.737297in}}%
\pgfpathlineto{\pgfqpoint{3.176501in}{2.764372in}}%
\pgfpathlineto{\pgfqpoint{3.150190in}{2.790690in}}%
\pgfpathlineto{\pgfqpoint{3.122834in}{2.816214in}}%
\pgfpathlineto{\pgfqpoint{3.094447in}{2.840904in}}%
\pgfpathlineto{\pgfqpoint{3.065042in}{2.864720in}}%
\pgfpathlineto{\pgfqpoint{3.034636in}{2.887623in}}%
\pgfpathlineto{\pgfqpoint{3.003243in}{2.909570in}}%
\pgfpathlineto{\pgfqpoint{2.970882in}{2.930519in}}%
\pgfpathlineto{\pgfqpoint{2.937572in}{2.950427in}}%
\pgfpathlineto{\pgfqpoint{2.903334in}{2.969251in}}%
\pgfpathlineto{\pgfqpoint{2.868189in}{2.986945in}}%
\pgfpathlineto{\pgfqpoint{2.832161in}{3.003464in}}%
\pgfpathlineto{\pgfqpoint{2.795277in}{3.018761in}}%
\pgfpathlineto{\pgfqpoint{2.757563in}{3.032790in}}%
\pgfpathlineto{\pgfqpoint{2.719051in}{3.045501in}}%
\pgfpathlineto{\pgfqpoint{2.679772in}{3.056846in}}%
\pgfpathlineto{\pgfqpoint{2.639761in}{3.066774in}}%
\pgfpathlineto{\pgfqpoint{2.599055in}{3.075235in}}%
\pgfpathlineto{\pgfqpoint{2.557695in}{3.082176in}}%
\pgfpathlineto{\pgfqpoint{2.515725in}{3.087547in}}%
\pgfpathlineto{\pgfqpoint{2.473192in}{3.091292in}}%
\pgfpathlineto{\pgfqpoint{2.430145in}{3.093359in}}%
\pgfpathlineto{\pgfqpoint{2.386641in}{3.093693in}}%
\pgfpathlineto{\pgfqpoint{2.342736in}{3.092239in}}%
\pgfpathlineto{\pgfqpoint{2.298495in}{3.088943in}}%
\pgfpathlineto{\pgfqpoint{2.253985in}{3.083749in}}%
\pgfpathlineto{\pgfqpoint{2.209281in}{3.076603in}}%
\pgfpathlineto{\pgfqpoint{2.175671in}{3.069929in}}%
\pgfpathlineto{\pgfqpoint{2.142032in}{3.062103in}}%
\pgfpathlineto{\pgfqpoint{2.108400in}{3.053103in}}%
\pgfpathlineto{\pgfqpoint{2.074814in}{3.042907in}}%
\pgfpathlineto{\pgfqpoint{2.041317in}{3.031493in}}%
\pgfpathlineto{\pgfqpoint{2.007952in}{3.018841in}}%
\pgfpathlineto{\pgfqpoint{1.974766in}{3.004929in}}%
\pgfpathlineto{\pgfqpoint{1.941807in}{2.989740in}}%
\pgfpathlineto{\pgfqpoint{1.909127in}{2.973253in}}%
\pgfpathlineto{\pgfqpoint{1.876779in}{2.955454in}}%
\pgfpathlineto{\pgfqpoint{1.844822in}{2.936325in}}%
\pgfpathlineto{\pgfqpoint{1.813313in}{2.915854in}}%
\pgfpathlineto{\pgfqpoint{1.782317in}{2.894027in}}%
\pgfpathlineto{\pgfqpoint{1.751898in}{2.870835in}}%
\pgfpathlineto{\pgfqpoint{1.722125in}{2.846271in}}%
\pgfpathlineto{\pgfqpoint{1.693071in}{2.820329in}}%
\pgfpathlineto{\pgfqpoint{1.664810in}{2.793008in}}%
\pgfpathlineto{\pgfqpoint{1.637421in}{2.764310in}}%
\pgfpathlineto{\pgfqpoint{1.610985in}{2.734239in}}%
\pgfpathlineto{\pgfqpoint{1.585586in}{2.702807in}}%
\pgfpathlineto{\pgfqpoint{1.561314in}{2.670027in}}%
\pgfpathlineto{\pgfqpoint{1.538257in}{2.635918in}}%
\pgfpathlineto{\pgfqpoint{1.516509in}{2.600507in}}%
\pgfpathlineto{\pgfqpoint{1.496167in}{2.563825in}}%
\pgfpathlineto{\pgfqpoint{1.477329in}{2.525909in}}%
\pgfpathlineto{\pgfqpoint{1.460094in}{2.486805in}}%
\pgfpathlineto{\pgfqpoint{1.444565in}{2.446567in}}%
\pgfpathlineto{\pgfqpoint{1.430843in}{2.405256in}}%
\pgfpathlineto{\pgfqpoint{1.419032in}{2.362942in}}%
\pgfpathlineto{\pgfqpoint{1.409234in}{2.319707in}}%
\pgfpathlineto{\pgfqpoint{1.403871in}{2.290415in}}%
\pgfpathlineto{\pgfqpoint{1.399477in}{2.260782in}}%
\pgfpathlineto{\pgfqpoint{1.396082in}{2.230840in}}%
\pgfpathlineto{\pgfqpoint{1.393713in}{2.200622in}}%
\pgfpathlineto{\pgfqpoint{1.392399in}{2.170164in}}%
\pgfpathlineto{\pgfqpoint{1.392166in}{2.139504in}}%
\pgfpathlineto{\pgfqpoint{1.393039in}{2.108681in}}%
\pgfpathlineto{\pgfqpoint{1.395042in}{2.077739in}}%
\pgfpathlineto{\pgfqpoint{1.398198in}{2.046722in}}%
\pgfpathlineto{\pgfqpoint{1.402528in}{2.015677in}}%
\pgfpathlineto{\pgfqpoint{1.408050in}{1.984653in}}%
\pgfpathlineto{\pgfqpoint{1.414781in}{1.953701in}}%
\pgfpathlineto{\pgfqpoint{1.422735in}{1.922874in}}%
\pgfpathlineto{\pgfqpoint{1.431924in}{1.892227in}}%
\pgfpathlineto{\pgfqpoint{1.442358in}{1.861818in}}%
\pgfpathlineto{\pgfqpoint{1.454043in}{1.831704in}}%
\pgfpathlineto{\pgfqpoint{1.466981in}{1.801946in}}%
\pgfpathlineto{\pgfqpoint{1.481172in}{1.772605in}}%
\pgfpathlineto{\pgfqpoint{1.496612in}{1.743743in}}%
\pgfpathlineto{\pgfqpoint{1.513293in}{1.715425in}}%
\pgfpathlineto{\pgfqpoint{1.531205in}{1.687714in}}%
\pgfpathlineto{\pgfqpoint{1.550332in}{1.660674in}}%
\pgfpathlineto{\pgfqpoint{1.570653in}{1.634369in}}%
\pgfpathlineto{\pgfqpoint{1.592147in}{1.608865in}}%
\pgfpathlineto{\pgfqpoint{1.614783in}{1.584223in}}%
\pgfpathlineto{\pgfqpoint{1.638531in}{1.560508in}}%
\pgfpathlineto{\pgfqpoint{1.663354in}{1.537780in}}%
\pgfpathlineto{\pgfqpoint{1.689212in}{1.516099in}}%
\pgfpathlineto{\pgfqpoint{1.716058in}{1.495523in}}%
\pgfpathlineto{\pgfqpoint{1.743845in}{1.476107in}}%
\pgfpathlineto{\pgfqpoint{1.772519in}{1.457904in}}%
\pgfpathlineto{\pgfqpoint{1.802025in}{1.440963in}}%
\pgfpathlineto{\pgfqpoint{1.832301in}{1.425331in}}%
\pgfpathlineto{\pgfqpoint{1.863285in}{1.411050in}}%
\pgfpathlineto{\pgfqpoint{1.894910in}{1.398159in}}%
\pgfpathlineto{\pgfqpoint{1.927108in}{1.386692in}}%
\pgfpathlineto{\pgfqpoint{1.959806in}{1.376679in}}%
\pgfpathlineto{\pgfqpoint{1.992932in}{1.368146in}}%
\pgfpathlineto{\pgfqpoint{2.026411in}{1.361113in}}%
\pgfpathlineto{\pgfqpoint{2.060167in}{1.355595in}}%
\pgfpathlineto{\pgfqpoint{2.094122in}{1.351604in}}%
\pgfpathlineto{\pgfqpoint{2.128200in}{1.349145in}}%
\pgfpathlineto{\pgfqpoint{2.162323in}{1.348217in}}%
\pgfpathlineto{\pgfqpoint{2.196415in}{1.348818in}}%
\pgfpathlineto{\pgfqpoint{2.230399in}{1.350938in}}%
\pgfpathlineto{\pgfqpoint{2.264201in}{1.354561in}}%
\pgfpathlineto{\pgfqpoint{2.297748in}{1.359671in}}%
\pgfpathlineto{\pgfqpoint{2.330968in}{1.366242in}}%
\pgfpathlineto{\pgfqpoint{2.363792in}{1.374249in}}%
\pgfpathlineto{\pgfqpoint{2.396154in}{1.383659in}}%
\pgfpathlineto{\pgfqpoint{2.427991in}{1.394437in}}%
\pgfpathlineto{\pgfqpoint{2.459241in}{1.406546in}}%
\pgfpathlineto{\pgfqpoint{2.489847in}{1.419942in}}%
\pgfpathlineto{\pgfqpoint{2.519756in}{1.434582in}}%
\pgfpathlineto{\pgfqpoint{2.548916in}{1.450419in}}%
\pgfpathlineto{\pgfqpoint{2.577282in}{1.467403in}}%
\pgfpathlineto{\pgfqpoint{2.604809in}{1.485484in}}%
\pgfpathlineto{\pgfqpoint{2.631459in}{1.504609in}}%
\pgfpathlineto{\pgfqpoint{2.657196in}{1.524725in}}%
\pgfpathlineto{\pgfqpoint{2.681988in}{1.545776in}}%
\pgfpathlineto{\pgfqpoint{2.705808in}{1.567708in}}%
\pgfpathlineto{\pgfqpoint{2.728629in}{1.590464in}}%
\pgfpathlineto{\pgfqpoint{2.750433in}{1.613990in}}%
\pgfpathlineto{\pgfqpoint{2.771200in}{1.638229in}}%
\pgfpathlineto{\pgfqpoint{2.790918in}{1.663127in}}%
\pgfpathlineto{\pgfqpoint{2.809574in}{1.688628in}}%
\pgfpathlineto{\pgfqpoint{2.827161in}{1.714679in}}%
\pgfpathlineto{\pgfqpoint{2.843674in}{1.741227in}}%
\pgfpathlineto{\pgfqpoint{2.859111in}{1.768221in}}%
\pgfpathlineto{\pgfqpoint{2.873471in}{1.795608in}}%
\pgfpathlineto{\pgfqpoint{2.886758in}{1.823341in}}%
\pgfpathlineto{\pgfqpoint{2.898976in}{1.851371in}}%
\pgfpathlineto{\pgfqpoint{2.910133in}{1.879652in}}%
\pgfpathlineto{\pgfqpoint{2.920237in}{1.908139in}}%
\pgfpathlineto{\pgfqpoint{2.929299in}{1.936789in}}%
\pgfpathlineto{\pgfqpoint{2.937331in}{1.965560in}}%
\pgfpathlineto{\pgfqpoint{2.944347in}{1.994412in}}%
\pgfpathlineto{\pgfqpoint{2.953001in}{2.037760in}}%
\pgfpathlineto{\pgfqpoint{2.959458in}{2.081083in}}%
\pgfpathlineto{\pgfqpoint{2.963780in}{2.124270in}}%
\pgfpathlineto{\pgfqpoint{2.966031in}{2.167214in}}%
\pgfpathlineto{\pgfqpoint{2.966280in}{2.209819in}}%
\pgfpathlineto{\pgfqpoint{2.964596in}{2.251997in}}%
\pgfpathlineto{\pgfqpoint{2.961052in}{2.293667in}}%
\pgfpathlineto{\pgfqpoint{2.955721in}{2.334754in}}%
\pgfpathlineto{\pgfqpoint{2.948677in}{2.375193in}}%
\pgfpathlineto{\pgfqpoint{2.939993in}{2.414922in}}%
\pgfpathlineto{\pgfqpoint{2.929743in}{2.453888in}}%
\pgfpathlineto{\pgfqpoint{2.918002in}{2.492042in}}%
\pgfpathlineto{\pgfqpoint{2.904840in}{2.529341in}}%
\pgfpathlineto{\pgfqpoint{2.890329in}{2.565749in}}%
\pgfpathlineto{\pgfqpoint{2.874539in}{2.601231in}}%
\pgfpathlineto{\pgfqpoint{2.857538in}{2.635758in}}%
\pgfpathlineto{\pgfqpoint{2.839392in}{2.669307in}}%
\pgfpathlineto{\pgfqpoint{2.820168in}{2.701856in}}%
\pgfpathlineto{\pgfqpoint{2.799927in}{2.733387in}}%
\pgfpathlineto{\pgfqpoint{2.778732in}{2.763886in}}%
\pgfpathlineto{\pgfqpoint{2.756642in}{2.793341in}}%
\pgfpathlineto{\pgfqpoint{2.733713in}{2.821743in}}%
\pgfpathlineto{\pgfqpoint{2.710002in}{2.849086in}}%
\pgfpathlineto{\pgfqpoint{2.685563in}{2.875365in}}%
\pgfpathlineto{\pgfqpoint{2.660446in}{2.900577in}}%
\pgfpathlineto{\pgfqpoint{2.634702in}{2.924723in}}%
\pgfpathlineto{\pgfqpoint{2.608380in}{2.947802in}}%
\pgfpathlineto{\pgfqpoint{2.581524in}{2.969817in}}%
\pgfpathlineto{\pgfqpoint{2.554181in}{2.990773in}}%
\pgfpathlineto{\pgfqpoint{2.517038in}{3.017073in}}%
\pgfpathlineto{\pgfqpoint{2.479199in}{3.041512in}}%
\pgfpathlineto{\pgfqpoint{2.440758in}{3.064108in}}%
\pgfpathlineto{\pgfqpoint{2.401802in}{3.084881in}}%
\pgfpathlineto{\pgfqpoint{2.362412in}{3.103852in}}%
\pgfpathlineto{\pgfqpoint{2.322669in}{3.121045in}}%
\pgfpathlineto{\pgfqpoint{2.282645in}{3.136488in}}%
\pgfpathlineto{\pgfqpoint{2.242412in}{3.150205in}}%
\pgfpathlineto{\pgfqpoint{2.202036in}{3.162225in}}%
\pgfpathlineto{\pgfqpoint{2.161581in}{3.172577in}}%
\pgfpathlineto{\pgfqpoint{2.121106in}{3.181291in}}%
\pgfpathlineto{\pgfqpoint{2.080669in}{3.188396in}}%
\pgfpathlineto{\pgfqpoint{2.040324in}{3.193923in}}%
\pgfpathlineto{\pgfqpoint{2.000122in}{3.197902in}}%
\pgfpathlineto{\pgfqpoint{1.960112in}{3.200366in}}%
\pgfpathlineto{\pgfqpoint{1.920340in}{3.201344in}}%
\pgfpathlineto{\pgfqpoint{1.880850in}{3.200869in}}%
\pgfpathlineto{\pgfqpoint{1.841685in}{3.198971in}}%
\pgfpathlineto{\pgfqpoint{1.802884in}{3.195682in}}%
\pgfpathlineto{\pgfqpoint{1.764486in}{3.191033in}}%
\pgfpathlineto{\pgfqpoint{1.726526in}{3.185055in}}%
\pgfpathlineto{\pgfqpoint{1.689039in}{3.177779in}}%
\pgfpathlineto{\pgfqpoint{1.652059in}{3.169235in}}%
\pgfpathlineto{\pgfqpoint{1.615616in}{3.159454in}}%
\pgfpathlineto{\pgfqpoint{1.579742in}{3.148468in}}%
\pgfpathlineto{\pgfqpoint{1.544465in}{3.136304in}}%
\pgfpathlineto{\pgfqpoint{1.509813in}{3.122995in}}%
\pgfpathlineto{\pgfqpoint{1.475812in}{3.108568in}}%
\pgfpathlineto{\pgfqpoint{1.442488in}{3.093055in}}%
\pgfpathlineto{\pgfqpoint{1.409866in}{3.076483in}}%
\pgfpathlineto{\pgfqpoint{1.377968in}{3.058883in}}%
\pgfpathlineto{\pgfqpoint{1.346818in}{3.040282in}}%
\pgfpathlineto{\pgfqpoint{1.316438in}{3.020711in}}%
\pgfpathlineto{\pgfqpoint{1.286847in}{3.000196in}}%
\pgfpathlineto{\pgfqpoint{1.258068in}{2.978767in}}%
\pgfpathlineto{\pgfqpoint{1.230119in}{2.956451in}}%
\pgfpathlineto{\pgfqpoint{1.203020in}{2.933278in}}%
\pgfpathlineto{\pgfqpoint{1.176788in}{2.909273in}}%
\pgfpathlineto{\pgfqpoint{1.151442in}{2.884466in}}%
\pgfpathlineto{\pgfqpoint{1.121030in}{2.852371in}}%
\pgfpathlineto{\pgfqpoint{1.092062in}{2.819118in}}%
\pgfpathlineto{\pgfqpoint{1.064568in}{2.784761in}}%
\pgfpathlineto{\pgfqpoint{1.038578in}{2.749355in}}%
\pgfpathlineto{\pgfqpoint{1.014123in}{2.712952in}}%
\pgfpathlineto{\pgfqpoint{0.991230in}{2.675607in}}%
\pgfpathlineto{\pgfqpoint{0.969928in}{2.637372in}}%
\pgfpathlineto{\pgfqpoint{0.950241in}{2.598302in}}%
\pgfpathlineto{\pgfqpoint{0.932197in}{2.558452in}}%
\pgfpathlineto{\pgfqpoint{0.915821in}{2.517874in}}%
\pgfpathlineto{\pgfqpoint{0.901137in}{2.476626in}}%
\pgfpathlineto{\pgfqpoint{0.888168in}{2.434760in}}%
\pgfpathlineto{\pgfqpoint{0.876938in}{2.392335in}}%
\pgfpathlineto{\pgfqpoint{0.867470in}{2.349406in}}%
\pgfpathlineto{\pgfqpoint{0.859785in}{2.306030in}}%
\pgfpathlineto{\pgfqpoint{0.853904in}{2.262266in}}%
\pgfpathlineto{\pgfqpoint{0.850512in}{2.227014in}}%
\pgfpathlineto{\pgfqpoint{0.848298in}{2.191583in}}%
\pgfpathlineto{\pgfqpoint{0.847273in}{2.156002in}}%
\pgfpathlineto{\pgfqpoint{0.847445in}{2.120305in}}%
\pgfpathlineto{\pgfqpoint{0.848825in}{2.084522in}}%
\pgfpathlineto{\pgfqpoint{0.851421in}{2.048687in}}%
\pgfpathlineto{\pgfqpoint{0.855242in}{2.012834in}}%
\pgfpathlineto{\pgfqpoint{0.860297in}{1.976995in}}%
\pgfpathlineto{\pgfqpoint{0.866593in}{1.941206in}}%
\pgfpathlineto{\pgfqpoint{0.874139in}{1.905503in}}%
\pgfpathlineto{\pgfqpoint{0.882942in}{1.869921in}}%
\pgfpathlineto{\pgfqpoint{0.893009in}{1.834498in}}%
\pgfpathlineto{\pgfqpoint{0.904346in}{1.799271in}}%
\pgfpathlineto{\pgfqpoint{0.916959in}{1.764279in}}%
\pgfpathlineto{\pgfqpoint{0.930853in}{1.729563in}}%
\pgfpathlineto{\pgfqpoint{0.946034in}{1.695163in}}%
\pgfpathlineto{\pgfqpoint{0.962506in}{1.661121in}}%
\pgfpathlineto{\pgfqpoint{0.980272in}{1.627481in}}%
\pgfpathlineto{\pgfqpoint{0.999335in}{1.594286in}}%
\pgfpathlineto{\pgfqpoint{1.019696in}{1.561582in}}%
\pgfpathlineto{\pgfqpoint{1.041358in}{1.529416in}}%
\pgfpathlineto{\pgfqpoint{1.064319in}{1.497836in}}%
\pgfpathlineto{\pgfqpoint{1.088579in}{1.466893in}}%
\pgfpathlineto{\pgfqpoint{1.114135in}{1.436637in}}%
\pgfpathlineto{\pgfqpoint{1.140984in}{1.407121in}}%
\pgfpathlineto{\pgfqpoint{1.169121in}{1.378400in}}%
\pgfpathlineto{\pgfqpoint{1.198539in}{1.350530in}}%
\pgfpathlineto{\pgfqpoint{1.229230in}{1.323570in}}%
\pgfpathlineto{\pgfqpoint{1.261183in}{1.297578in}}%
\pgfpathlineto{\pgfqpoint{1.294387in}{1.272617in}}%
\pgfpathlineto{\pgfqpoint{1.328827in}{1.248751in}}%
\pgfpathlineto{\pgfqpoint{1.364485in}{1.226044in}}%
\pgfpathlineto{\pgfqpoint{1.401343in}{1.204566in}}%
\pgfpathlineto{\pgfqpoint{1.439376in}{1.184386in}}%
\pgfpathlineto{\pgfqpoint{1.478559in}{1.165576in}}%
\pgfpathlineto{\pgfqpoint{1.518862in}{1.148210in}}%
\pgfpathlineto{\pgfqpoint{1.560252in}{1.132364in}}%
\pgfpathlineto{\pgfqpoint{1.602691in}{1.118118in}}%
\pgfpathlineto{\pgfqpoint{1.646135in}{1.105551in}}%
\pgfpathlineto{\pgfqpoint{1.690537in}{1.094746in}}%
\pgfpathlineto{\pgfqpoint{1.724435in}{1.087849in}}%
\pgfpathlineto{\pgfqpoint{1.758817in}{1.082026in}}%
\pgfpathlineto{\pgfqpoint{1.793656in}{1.077316in}}%
\pgfpathlineto{\pgfqpoint{1.828924in}{1.073754in}}%
\pgfpathlineto{\pgfqpoint{1.864587in}{1.071378in}}%
\pgfpathlineto{\pgfqpoint{1.900614in}{1.070226in}}%
\pgfpathlineto{\pgfqpoint{1.936966in}{1.070337in}}%
\pgfpathlineto{\pgfqpoint{1.973605in}{1.071748in}}%
\pgfpathlineto{\pgfqpoint{2.010487in}{1.074498in}}%
\pgfpathlineto{\pgfqpoint{2.047569in}{1.078625in}}%
\pgfpathlineto{\pgfqpoint{2.084800in}{1.084167in}}%
\pgfpathlineto{\pgfqpoint{2.122128in}{1.091160in}}%
\pgfpathlineto{\pgfqpoint{2.159498in}{1.099642in}}%
\pgfpathlineto{\pgfqpoint{2.196848in}{1.109649in}}%
\pgfpathlineto{\pgfqpoint{2.234116in}{1.121213in}}%
\pgfpathlineto{\pgfqpoint{2.271233in}{1.134370in}}%
\pgfpathlineto{\pgfqpoint{2.308127in}{1.149149in}}%
\pgfpathlineto{\pgfqpoint{2.344720in}{1.165580in}}%
\pgfpathlineto{\pgfqpoint{2.380930in}{1.183688in}}%
\pgfpathlineto{\pgfqpoint{2.416672in}{1.203498in}}%
\pgfpathlineto{\pgfqpoint{2.451853in}{1.225029in}}%
\pgfpathlineto{\pgfqpoint{2.486378in}{1.248296in}}%
\pgfpathlineto{\pgfqpoint{2.520145in}{1.273311in}}%
\pgfpathlineto{\pgfqpoint{2.553047in}{1.300080in}}%
\pgfpathlineto{\pgfqpoint{2.584976in}{1.328602in}}%
\pgfpathlineto{\pgfqpoint{2.615813in}{1.358870in}}%
\pgfpathlineto{\pgfqpoint{2.645441in}{1.390870in}}%
\pgfpathlineto{\pgfqpoint{2.673735in}{1.424581in}}%
\pgfpathlineto{\pgfqpoint{2.700567in}{1.459970in}}%
\pgfpathlineto{\pgfqpoint{2.717579in}{1.484476in}}%
\pgfpathlineto{\pgfqpoint{2.733845in}{1.509696in}}%
\pgfpathlineto{\pgfqpoint{2.749323in}{1.535612in}}%
\pgfpathlineto{\pgfqpoint{2.763976in}{1.562206in}}%
\pgfpathlineto{\pgfqpoint{2.777762in}{1.589456in}}%
\pgfpathlineto{\pgfqpoint{2.790642in}{1.617338in}}%
\pgfpathlineto{\pgfqpoint{2.802575in}{1.645826in}}%
\pgfpathlineto{\pgfqpoint{2.813523in}{1.674890in}}%
\pgfpathlineto{\pgfqpoint{2.823445in}{1.704499in}}%
\pgfpathlineto{\pgfqpoint{2.832304in}{1.734619in}}%
\pgfpathlineto{\pgfqpoint{2.840063in}{1.765211in}}%
\pgfpathlineto{\pgfqpoint{2.846684in}{1.796236in}}%
\pgfpathlineto{\pgfqpoint{2.852133in}{1.827652in}}%
\pgfpathlineto{\pgfqpoint{2.856376in}{1.859411in}}%
\pgfpathlineto{\pgfqpoint{2.859382in}{1.891466in}}%
\pgfpathlineto{\pgfqpoint{2.861121in}{1.923765in}}%
\pgfpathlineto{\pgfqpoint{2.861566in}{1.956254in}}%
\pgfpathlineto{\pgfqpoint{2.860691in}{1.988877in}}%
\pgfpathlineto{\pgfqpoint{2.858475in}{2.021575in}}%
\pgfpathlineto{\pgfqpoint{2.854899in}{2.054285in}}%
\pgfpathlineto{\pgfqpoint{2.849947in}{2.086946in}}%
\pgfpathlineto{\pgfqpoint{2.843608in}{2.119492in}}%
\pgfpathlineto{\pgfqpoint{2.835873in}{2.151854in}}%
\pgfpathlineto{\pgfqpoint{2.826737in}{2.183966in}}%
\pgfpathlineto{\pgfqpoint{2.816200in}{2.215757in}}%
\pgfpathlineto{\pgfqpoint{2.804266in}{2.247156in}}%
\pgfpathlineto{\pgfqpoint{2.790945in}{2.278093in}}%
\pgfpathlineto{\pgfqpoint{2.776248in}{2.308496in}}%
\pgfpathlineto{\pgfqpoint{2.760194in}{2.338294in}}%
\pgfpathlineto{\pgfqpoint{2.742804in}{2.367416in}}%
\pgfpathlineto{\pgfqpoint{2.724107in}{2.395792in}}%
\pgfpathlineto{\pgfqpoint{2.704134in}{2.423355in}}%
\pgfpathlineto{\pgfqpoint{2.682920in}{2.450038in}}%
\pgfpathlineto{\pgfqpoint{2.660508in}{2.475774in}}%
\pgfpathlineto{\pgfqpoint{2.636941in}{2.500504in}}%
\pgfpathlineto{\pgfqpoint{2.612269in}{2.524166in}}%
\pgfpathlineto{\pgfqpoint{2.586545in}{2.546705in}}%
\pgfpathlineto{\pgfqpoint{2.559827in}{2.568068in}}%
\pgfpathlineto{\pgfqpoint{2.532174in}{2.588206in}}%
\pgfpathlineto{\pgfqpoint{2.503649in}{2.607074in}}%
\pgfpathlineto{\pgfqpoint{2.474319in}{2.624631in}}%
\pgfpathlineto{\pgfqpoint{2.444252in}{2.640840in}}%
\pgfpathlineto{\pgfqpoint{2.413518in}{2.655671in}}%
\pgfpathlineto{\pgfqpoint{2.382189in}{2.669096in}}%
\pgfpathlineto{\pgfqpoint{2.350338in}{2.681092in}}%
\pgfpathlineto{\pgfqpoint{2.318040in}{2.691643in}}%
\pgfpathlineto{\pgfqpoint{2.285369in}{2.700737in}}%
\pgfpathlineto{\pgfqpoint{2.252398in}{2.708365in}}%
\pgfpathlineto{\pgfqpoint{2.219203in}{2.714525in}}%
\pgfpathlineto{\pgfqpoint{2.185855in}{2.719219in}}%
\pgfpathlineto{\pgfqpoint{2.152428in}{2.722452in}}%
\pgfpathlineto{\pgfqpoint{2.118992in}{2.724236in}}%
\pgfpathlineto{\pgfqpoint{2.085617in}{2.724586in}}%
\pgfpathlineto{\pgfqpoint{2.052368in}{2.723519in}}%
\pgfpathlineto{\pgfqpoint{2.019312in}{2.721059in}}%
\pgfpathlineto{\pgfqpoint{1.986511in}{2.717231in}}%
\pgfpathlineto{\pgfqpoint{1.954024in}{2.712064in}}%
\pgfpathlineto{\pgfqpoint{1.921908in}{2.705590in}}%
\pgfpathlineto{\pgfqpoint{1.890217in}{2.697843in}}%
\pgfpathlineto{\pgfqpoint{1.859003in}{2.688860in}}%
\pgfpathlineto{\pgfqpoint{1.828313in}{2.678680in}}%
\pgfpathlineto{\pgfqpoint{1.798193in}{2.667344in}}%
\pgfpathlineto{\pgfqpoint{1.768682in}{2.654893in}}%
\pgfpathlineto{\pgfqpoint{1.739821in}{2.641371in}}%
\pgfpathlineto{\pgfqpoint{1.711644in}{2.626823in}}%
\pgfpathlineto{\pgfqpoint{1.684184in}{2.611292in}}%
\pgfpathlineto{\pgfqpoint{1.657469in}{2.594826in}}%
\pgfpathlineto{\pgfqpoint{1.631525in}{2.577469in}}%
\pgfpathlineto{\pgfqpoint{1.606376in}{2.559268in}}%
\pgfpathlineto{\pgfqpoint{1.582041in}{2.540269in}}%
\pgfpathlineto{\pgfqpoint{1.558540in}{2.520517in}}%
\pgfpathlineto{\pgfqpoint{1.535885in}{2.500059in}}%
\pgfpathlineto{\pgfqpoint{1.514091in}{2.478938in}}%
\pgfpathlineto{\pgfqpoint{1.493166in}{2.457199in}}%
\pgfpathlineto{\pgfqpoint{1.473120in}{2.434886in}}%
\pgfpathlineto{\pgfqpoint{1.453957in}{2.412040in}}%
\pgfpathlineto{\pgfqpoint{1.426878in}{2.376866in}}%
\pgfpathlineto{\pgfqpoint{1.401800in}{2.340725in}}%
\pgfpathlineto{\pgfqpoint{1.378720in}{2.303747in}}%
\pgfpathlineto{\pgfqpoint{1.357628in}{2.266059in}}%
\pgfpathlineto{\pgfqpoint{1.338504in}{2.227779in}}%
\pgfpathlineto{\pgfqpoint{1.321325in}{2.189019in}}%
\pgfpathlineto{\pgfqpoint{1.306062in}{2.149888in}}%
\pgfpathlineto{\pgfqpoint{1.292682in}{2.110484in}}%
\pgfpathlineto{\pgfqpoint{1.281148in}{2.070904in}}%
\pgfpathlineto{\pgfqpoint{1.271419in}{2.031235in}}%
\pgfpathlineto{\pgfqpoint{1.263453in}{1.991561in}}%
\pgfpathlineto{\pgfqpoint{1.257203in}{1.951958in}}%
\pgfpathlineto{\pgfqpoint{1.252625in}{1.912498in}}%
\pgfpathlineto{\pgfqpoint{1.249668in}{1.873249in}}%
\pgfpathlineto{\pgfqpoint{1.248285in}{1.834271in}}%
\pgfpathlineto{\pgfqpoint{1.248426in}{1.795622in}}%
\pgfpathlineto{\pgfqpoint{1.250041in}{1.757354in}}%
\pgfpathlineto{\pgfqpoint{1.253079in}{1.719515in}}%
\pgfpathlineto{\pgfqpoint{1.257492in}{1.682151in}}%
\pgfpathlineto{\pgfqpoint{1.263228in}{1.645302in}}%
\pgfpathlineto{\pgfqpoint{1.270240in}{1.609005in}}%
\pgfpathlineto{\pgfqpoint{1.278479in}{1.573293in}}%
\pgfpathlineto{\pgfqpoint{1.287896in}{1.538199in}}%
\pgfpathlineto{\pgfqpoint{1.298444in}{1.503749in}}%
\pgfpathlineto{\pgfqpoint{1.310078in}{1.469969in}}%
\pgfpathlineto{\pgfqpoint{1.322752in}{1.436882in}}%
\pgfpathlineto{\pgfqpoint{1.336421in}{1.404508in}}%
\pgfpathlineto{\pgfqpoint{1.351041in}{1.372866in}}%
\pgfpathlineto{\pgfqpoint{1.366572in}{1.341972in}}%
\pgfpathlineto{\pgfqpoint{1.388622in}{1.301967in}}%
\pgfpathlineto{\pgfqpoint{1.412121in}{1.263347in}}%
\pgfpathlineto{\pgfqpoint{1.436976in}{1.226137in}}%
\pgfpathlineto{\pgfqpoint{1.463098in}{1.190355in}}%
\pgfpathlineto{\pgfqpoint{1.490402in}{1.156020in}}%
\pgfpathlineto{\pgfqpoint{1.518805in}{1.123142in}}%
\pgfpathlineto{\pgfqpoint{1.548229in}{1.091731in}}%
\pgfpathlineto{\pgfqpoint{1.578598in}{1.061794in}}%
\pgfpathlineto{\pgfqpoint{1.609837in}{1.033334in}}%
\pgfpathlineto{\pgfqpoint{1.641876in}{1.006353in}}%
\pgfpathlineto{\pgfqpoint{1.674648in}{0.980850in}}%
\pgfpathlineto{\pgfqpoint{1.708087in}{0.956822in}}%
\pgfpathlineto{\pgfqpoint{1.742130in}{0.934266in}}%
\pgfpathlineto{\pgfqpoint{1.776717in}{0.913175in}}%
\pgfpathlineto{\pgfqpoint{1.811791in}{0.893542in}}%
\pgfpathlineto{\pgfqpoint{1.847294in}{0.875359in}}%
\pgfpathlineto{\pgfqpoint{1.883174in}{0.858617in}}%
\pgfpathlineto{\pgfqpoint{1.919377in}{0.843305in}}%
\pgfpathlineto{\pgfqpoint{1.955855in}{0.829412in}}%
\pgfpathlineto{\pgfqpoint{1.992559in}{0.816926in}}%
\pgfpathlineto{\pgfqpoint{2.029442in}{0.805835in}}%
\pgfpathlineto{\pgfqpoint{2.066460in}{0.796125in}}%
\pgfpathlineto{\pgfqpoint{2.103568in}{0.787783in}}%
\pgfpathlineto{\pgfqpoint{2.140726in}{0.780794in}}%
\pgfpathlineto{\pgfqpoint{2.177891in}{0.775144in}}%
\pgfpathlineto{\pgfqpoint{2.215025in}{0.770817in}}%
\pgfpathlineto{\pgfqpoint{2.252090in}{0.767799in}}%
\pgfpathlineto{\pgfqpoint{2.289049in}{0.766074in}}%
\pgfpathlineto{\pgfqpoint{2.325865in}{0.765626in}}%
\pgfpathlineto{\pgfqpoint{2.362504in}{0.766439in}}%
\pgfpathlineto{\pgfqpoint{2.398932in}{0.768496in}}%
\pgfpathlineto{\pgfqpoint{2.435116in}{0.771782in}}%
\pgfpathlineto{\pgfqpoint{2.471023in}{0.776278in}}%
\pgfpathlineto{\pgfqpoint{2.506623in}{0.781969in}}%
\pgfpathlineto{\pgfqpoint{2.541884in}{0.788837in}}%
\pgfpathlineto{\pgfqpoint{2.576777in}{0.796865in}}%
\pgfpathlineto{\pgfqpoint{2.611273in}{0.806037in}}%
\pgfpathlineto{\pgfqpoint{2.645342in}{0.816333in}}%
\pgfpathlineto{\pgfqpoint{2.678957in}{0.827737in}}%
\pgfpathlineto{\pgfqpoint{2.720296in}{0.843523in}}%
\pgfpathlineto{\pgfqpoint{2.760829in}{0.860977in}}%
\pgfpathlineto{\pgfqpoint{2.800506in}{0.880064in}}%
\pgfpathlineto{\pgfqpoint{2.839276in}{0.900748in}}%
\pgfpathlineto{\pgfqpoint{2.877089in}{0.922993in}}%
\pgfpathlineto{\pgfqpoint{2.913898in}{0.946763in}}%
\pgfpathlineto{\pgfqpoint{2.949654in}{0.972020in}}%
\pgfpathlineto{\pgfqpoint{2.984311in}{0.998727in}}%
\pgfpathlineto{\pgfqpoint{3.017820in}{1.026846in}}%
\pgfpathlineto{\pgfqpoint{3.050138in}{1.056339in}}%
\pgfpathlineto{\pgfqpoint{3.081218in}{1.087166in}}%
\pgfpathlineto{\pgfqpoint{3.111016in}{1.119288in}}%
\pgfpathlineto{\pgfqpoint{3.139486in}{1.152663in}}%
\pgfpathlineto{\pgfqpoint{3.166585in}{1.187250in}}%
\pgfpathlineto{\pgfqpoint{3.192268in}{1.223007in}}%
\pgfpathlineto{\pgfqpoint{3.216491in}{1.259890in}}%
\pgfpathlineto{\pgfqpoint{3.234790in}{1.290177in}}%
\pgfpathlineto{\pgfqpoint{3.252104in}{1.321133in}}%
\pgfpathlineto{\pgfqpoint{3.268412in}{1.352734in}}%
\pgfpathlineto{\pgfqpoint{3.283691in}{1.384956in}}%
\pgfpathlineto{\pgfqpoint{3.297920in}{1.417773in}}%
\pgfpathlineto{\pgfqpoint{3.311075in}{1.451161in}}%
\pgfpathlineto{\pgfqpoint{3.323136in}{1.485092in}}%
\pgfpathlineto{\pgfqpoint{3.334080in}{1.519540in}}%
\pgfpathlineto{\pgfqpoint{3.343885in}{1.554477in}}%
\pgfpathlineto{\pgfqpoint{3.352528in}{1.589874in}}%
\pgfpathlineto{\pgfqpoint{3.359989in}{1.625702in}}%
\pgfpathlineto{\pgfqpoint{3.366245in}{1.661931in}}%
\pgfpathlineto{\pgfqpoint{3.371274in}{1.698529in}}%
\pgfpathlineto{\pgfqpoint{3.375055in}{1.735464in}}%
\pgfpathlineto{\pgfqpoint{3.377565in}{1.772702in}}%
\pgfpathlineto{\pgfqpoint{3.378783in}{1.810209in}}%
\pgfpathlineto{\pgfqpoint{3.378687in}{1.847950in}}%
\pgfpathlineto{\pgfqpoint{3.377256in}{1.885886in}}%
\pgfpathlineto{\pgfqpoint{3.374469in}{1.923980in}}%
\pgfpathlineto{\pgfqpoint{3.370305in}{1.962191in}}%
\pgfpathlineto{\pgfqpoint{3.364742in}{2.000478in}}%
\pgfpathlineto{\pgfqpoint{3.357760in}{2.038798in}}%
\pgfpathlineto{\pgfqpoint{3.349340in}{2.077106in}}%
\pgfpathlineto{\pgfqpoint{3.339461in}{2.115354in}}%
\pgfpathlineto{\pgfqpoint{3.328105in}{2.153495in}}%
\pgfpathlineto{\pgfqpoint{3.315251in}{2.191477in}}%
\pgfpathlineto{\pgfqpoint{3.300883in}{2.229248in}}%
\pgfpathlineto{\pgfqpoint{3.284984in}{2.266751in}}%
\pgfpathlineto{\pgfqpoint{3.267536in}{2.303931in}}%
\pgfpathlineto{\pgfqpoint{3.248524in}{2.340725in}}%
\pgfpathlineto{\pgfqpoint{3.227935in}{2.377072in}}%
\pgfpathlineto{\pgfqpoint{3.205756in}{2.412905in}}%
\pgfpathlineto{\pgfqpoint{3.181975in}{2.448156in}}%
\pgfpathlineto{\pgfqpoint{3.156582in}{2.482752in}}%
\pgfpathlineto{\pgfqpoint{3.129571in}{2.516619in}}%
\pgfpathlineto{\pgfqpoint{3.100937in}{2.549679in}}%
\pgfpathlineto{\pgfqpoint{3.070675in}{2.581848in}}%
\pgfpathlineto{\pgfqpoint{3.038787in}{2.613041in}}%
\pgfpathlineto{\pgfqpoint{3.005277in}{2.643169in}}%
\pgfpathlineto{\pgfqpoint{2.970150in}{2.672138in}}%
\pgfpathlineto{\pgfqpoint{2.942752in}{2.693045in}}%
\pgfpathlineto{\pgfqpoint{2.914458in}{2.713203in}}%
\pgfpathlineto{\pgfqpoint{2.885276in}{2.732567in}}%
\pgfpathlineto{\pgfqpoint{2.855215in}{2.751092in}}%
\pgfpathlineto{\pgfqpoint{2.824289in}{2.768732in}}%
\pgfpathlineto{\pgfqpoint{2.792509in}{2.785438in}}%
\pgfpathlineto{\pgfqpoint{2.759893in}{2.801160in}}%
\pgfpathlineto{\pgfqpoint{2.726459in}{2.815849in}}%
\pgfpathlineto{\pgfqpoint{2.692228in}{2.829452in}}%
\pgfpathlineto{\pgfqpoint{2.657223in}{2.841916in}}%
\pgfpathlineto{\pgfqpoint{2.621471in}{2.853186in}}%
\pgfpathlineto{\pgfqpoint{2.585003in}{2.863207in}}%
\pgfpathlineto{\pgfqpoint{2.547853in}{2.871924in}}%
\pgfpathlineto{\pgfqpoint{2.510057in}{2.879277in}}%
\pgfpathlineto{\pgfqpoint{2.471657in}{2.885210in}}%
\pgfpathlineto{\pgfqpoint{2.432700in}{2.889664in}}%
\pgfpathlineto{\pgfqpoint{2.393234in}{2.892580in}}%
\pgfpathlineto{\pgfqpoint{2.353317in}{2.893899in}}%
\pgfpathlineto{\pgfqpoint{2.313008in}{2.893563in}}%
\pgfpathlineto{\pgfqpoint{2.272372in}{2.891514in}}%
\pgfpathlineto{\pgfqpoint{2.231483in}{2.887694in}}%
\pgfpathlineto{\pgfqpoint{2.190417in}{2.882049in}}%
\pgfpathlineto{\pgfqpoint{2.149259in}{2.874523in}}%
\pgfpathlineto{\pgfqpoint{2.108099in}{2.865066in}}%
\pgfpathlineto{\pgfqpoint{2.067035in}{2.853630in}}%
\pgfpathlineto{\pgfqpoint{2.026171in}{2.840170in}}%
\pgfpathlineto{\pgfqpoint{1.985618in}{2.824645in}}%
\pgfpathlineto{\pgfqpoint{1.945494in}{2.807021in}}%
\pgfpathlineto{\pgfqpoint{1.905926in}{2.787270in}}%
\pgfpathlineto{\pgfqpoint{1.879921in}{2.772909in}}%
\pgfpathlineto{\pgfqpoint{1.854261in}{2.757588in}}%
\pgfpathlineto{\pgfqpoint{1.828991in}{2.741305in}}%
\pgfpathlineto{\pgfqpoint{1.804152in}{2.724057in}}%
\pgfpathlineto{\pgfqpoint{1.779790in}{2.705848in}}%
\pgfpathlineto{\pgfqpoint{1.755949in}{2.686678in}}%
\pgfpathlineto{\pgfqpoint{1.732678in}{2.666556in}}%
\pgfpathlineto{\pgfqpoint{1.710023in}{2.645487in}}%
\pgfpathlineto{\pgfqpoint{1.688033in}{2.623482in}}%
\pgfpathlineto{\pgfqpoint{1.666757in}{2.600555in}}%
\pgfpathlineto{\pgfqpoint{1.646244in}{2.576722in}}%
\pgfpathlineto{\pgfqpoint{1.626545in}{2.552002in}}%
\pgfpathlineto{\pgfqpoint{1.607708in}{2.526415in}}%
\pgfpathlineto{\pgfqpoint{1.589785in}{2.499989in}}%
\pgfpathlineto{\pgfqpoint{1.572824in}{2.472750in}}%
\pgfpathlineto{\pgfqpoint{1.556874in}{2.444731in}}%
\pgfpathlineto{\pgfqpoint{1.541984in}{2.415966in}}%
\pgfpathlineto{\pgfqpoint{1.528202in}{2.386494in}}%
\pgfpathlineto{\pgfqpoint{1.515572in}{2.356358in}}%
\pgfpathlineto{\pgfqpoint{1.504140in}{2.325601in}}%
\pgfpathlineto{\pgfqpoint{1.493947in}{2.294272in}}%
\pgfpathlineto{\pgfqpoint{1.485035in}{2.262425in}}%
\pgfpathlineto{\pgfqpoint{1.477440in}{2.230112in}}%
\pgfpathlineto{\pgfqpoint{1.471198in}{2.197394in}}%
\pgfpathlineto{\pgfqpoint{1.466341in}{2.164331in}}%
\pgfpathlineto{\pgfqpoint{1.462898in}{2.130986in}}%
\pgfpathlineto{\pgfqpoint{1.460893in}{2.097427in}}%
\pgfpathlineto{\pgfqpoint{1.460347in}{2.063722in}}%
\pgfpathlineto{\pgfqpoint{1.461278in}{2.029942in}}%
\pgfpathlineto{\pgfqpoint{1.463699in}{1.996159in}}%
\pgfpathlineto{\pgfqpoint{1.467616in}{1.962448in}}%
\pgfpathlineto{\pgfqpoint{1.473035in}{1.928882in}}%
\pgfpathlineto{\pgfqpoint{1.479954in}{1.895538in}}%
\pgfpathlineto{\pgfqpoint{1.488366in}{1.862491in}}%
\pgfpathlineto{\pgfqpoint{1.498260in}{1.829818in}}%
\pgfpathlineto{\pgfqpoint{1.509620in}{1.797593in}}%
\pgfpathlineto{\pgfqpoint{1.522426in}{1.765891in}}%
\pgfpathlineto{\pgfqpoint{1.536652in}{1.734785in}}%
\pgfpathlineto{\pgfqpoint{1.552266in}{1.704347in}}%
\pgfpathlineto{\pgfqpoint{1.569233in}{1.674644in}}%
\pgfpathlineto{\pgfqpoint{1.587515in}{1.645745in}}%
\pgfpathlineto{\pgfqpoint{1.607066in}{1.617714in}}%
\pgfpathlineto{\pgfqpoint{1.627840in}{1.590610in}}%
\pgfpathlineto{\pgfqpoint{1.649785in}{1.564491in}}%
\pgfpathlineto{\pgfqpoint{1.672845in}{1.539412in}}%
\pgfpathlineto{\pgfqpoint{1.696964in}{1.515421in}}%
\pgfpathlineto{\pgfqpoint{1.722080in}{1.492565in}}%
\pgfpathlineto{\pgfqpoint{1.748132in}{1.470885in}}%
\pgfpathlineto{\pgfqpoint{1.775053in}{1.450417in}}%
\pgfpathlineto{\pgfqpoint{1.802779in}{1.431194in}}%
\pgfpathlineto{\pgfqpoint{1.831240in}{1.413245in}}%
\pgfpathlineto{\pgfqpoint{1.860370in}{1.396592in}}%
\pgfpathlineto{\pgfqpoint{1.890099in}{1.381254in}}%
\pgfpathlineto{\pgfqpoint{1.920359in}{1.367247in}}%
\pgfpathlineto{\pgfqpoint{1.951081in}{1.354579in}}%
\pgfpathlineto{\pgfqpoint{1.982196in}{1.343257in}}%
\pgfpathlineto{\pgfqpoint{2.013638in}{1.333283in}}%
\pgfpathlineto{\pgfqpoint{2.045341in}{1.324654in}}%
\pgfpathlineto{\pgfqpoint{2.077240in}{1.317364in}}%
\pgfpathlineto{\pgfqpoint{2.109272in}{1.311405in}}%
\pgfpathlineto{\pgfqpoint{2.141375in}{1.306763in}}%
\pgfpathlineto{\pgfqpoint{2.173491in}{1.303423in}}%
\pgfpathlineto{\pgfqpoint{2.205563in}{1.301365in}}%
\pgfpathlineto{\pgfqpoint{2.237536in}{1.300569in}}%
\pgfpathlineto{\pgfqpoint{2.269357in}{1.301010in}}%
\pgfpathlineto{\pgfqpoint{2.300976in}{1.302663in}}%
\pgfpathlineto{\pgfqpoint{2.332347in}{1.305500in}}%
\pgfpathlineto{\pgfqpoint{2.363424in}{1.309491in}}%
\pgfpathlineto{\pgfqpoint{2.394165in}{1.314606in}}%
\pgfpathlineto{\pgfqpoint{2.424531in}{1.320811in}}%
\pgfpathlineto{\pgfqpoint{2.454484in}{1.328074in}}%
\pgfpathlineto{\pgfqpoint{2.483991in}{1.336360in}}%
\pgfpathlineto{\pgfqpoint{2.513018in}{1.345634in}}%
\pgfpathlineto{\pgfqpoint{2.541536in}{1.355861in}}%
\pgfpathlineto{\pgfqpoint{2.569519in}{1.367005in}}%
\pgfpathlineto{\pgfqpoint{2.596941in}{1.379029in}}%
\pgfpathlineto{\pgfqpoint{2.623780in}{1.391898in}}%
\pgfpathlineto{\pgfqpoint{2.662902in}{1.412706in}}%
\pgfpathlineto{\pgfqpoint{2.700606in}{1.435211in}}%
\pgfpathlineto{\pgfqpoint{2.736842in}{1.459293in}}%
\pgfpathlineto{\pgfqpoint{2.771571in}{1.484836in}}%
\pgfpathlineto{\pgfqpoint{2.804758in}{1.511724in}}%
\pgfpathlineto{\pgfqpoint{2.836380in}{1.539845in}}%
\pgfpathlineto{\pgfqpoint{2.866419in}{1.569092in}}%
\pgfpathlineto{\pgfqpoint{2.894864in}{1.599358in}}%
\pgfpathlineto{\pgfqpoint{2.921710in}{1.630545in}}%
\pgfpathlineto{\pgfqpoint{2.946958in}{1.662556in}}%
\pgfpathlineto{\pgfqpoint{2.970611in}{1.695297in}}%
\pgfpathlineto{\pgfqpoint{2.992680in}{1.728681in}}%
\pgfpathlineto{\pgfqpoint{3.013177in}{1.762624in}}%
\pgfpathlineto{\pgfqpoint{3.032117in}{1.797046in}}%
\pgfpathlineto{\pgfqpoint{3.049521in}{1.831871in}}%
\pgfpathlineto{\pgfqpoint{3.065409in}{1.867028in}}%
\pgfpathlineto{\pgfqpoint{3.079805in}{1.902448in}}%
\pgfpathlineto{\pgfqpoint{3.092734in}{1.938066in}}%
\pgfpathlineto{\pgfqpoint{3.104222in}{1.973823in}}%
\pgfpathlineto{\pgfqpoint{3.114299in}{2.009662in}}%
\pgfpathlineto{\pgfqpoint{3.122992in}{2.045528in}}%
\pgfpathlineto{\pgfqpoint{3.130332in}{2.081370in}}%
\pgfpathlineto{\pgfqpoint{3.136349in}{2.117142in}}%
\pgfpathlineto{\pgfqpoint{3.141075in}{2.152799in}}%
\pgfpathlineto{\pgfqpoint{3.144540in}{2.188298in}}%
\pgfpathlineto{\pgfqpoint{3.146777in}{2.223600in}}%
\pgfpathlineto{\pgfqpoint{3.147816in}{2.258669in}}%
\pgfpathlineto{\pgfqpoint{3.147690in}{2.293469in}}%
\pgfpathlineto{\pgfqpoint{3.146430in}{2.327970in}}%
\pgfpathlineto{\pgfqpoint{3.143041in}{2.373451in}}%
\pgfpathlineto{\pgfqpoint{3.137765in}{2.418280in}}%
\pgfpathlineto{\pgfqpoint{3.130675in}{2.462394in}}%
\pgfpathlineto{\pgfqpoint{3.121844in}{2.505740in}}%
\pgfpathlineto{\pgfqpoint{3.111342in}{2.548268in}}%
\pgfpathlineto{\pgfqpoint{3.099238in}{2.589931in}}%
\pgfpathlineto{\pgfqpoint{3.085599in}{2.630690in}}%
\pgfpathlineto{\pgfqpoint{3.070492in}{2.670505in}}%
\pgfpathlineto{\pgfqpoint{3.053982in}{2.709344in}}%
\pgfpathlineto{\pgfqpoint{3.036131in}{2.747177in}}%
\pgfpathlineto{\pgfqpoint{3.017002in}{2.783974in}}%
\pgfpathlineto{\pgfqpoint{2.996653in}{2.819712in}}%
\pgfpathlineto{\pgfqpoint{2.975144in}{2.854369in}}%
\pgfpathlineto{\pgfqpoint{2.952532in}{2.887924in}}%
\pgfpathlineto{\pgfqpoint{2.928872in}{2.920359in}}%
\pgfpathlineto{\pgfqpoint{2.904218in}{2.951660in}}%
\pgfpathlineto{\pgfqpoint{2.878623in}{2.981811in}}%
\pgfpathlineto{\pgfqpoint{2.852137in}{3.010801in}}%
\pgfpathlineto{\pgfqpoint{2.824812in}{3.038618in}}%
\pgfpathlineto{\pgfqpoint{2.796694in}{3.065255in}}%
\pgfpathlineto{\pgfqpoint{2.767833in}{3.090701in}}%
\pgfpathlineto{\pgfqpoint{2.738273in}{3.114951in}}%
\pgfpathlineto{\pgfqpoint{2.708060in}{3.137999in}}%
\pgfpathlineto{\pgfqpoint{2.677238in}{3.159840in}}%
\pgfpathlineto{\pgfqpoint{2.645850in}{3.180470in}}%
\pgfpathlineto{\pgfqpoint{2.613937in}{3.199887in}}%
\pgfpathlineto{\pgfqpoint{2.581541in}{3.218088in}}%
\pgfpathlineto{\pgfqpoint{2.548702in}{3.235072in}}%
\pgfpathlineto{\pgfqpoint{2.515458in}{3.250838in}}%
\pgfpathlineto{\pgfqpoint{2.481848in}{3.265387in}}%
\pgfpathlineto{\pgfqpoint{2.447910in}{3.278718in}}%
\pgfpathlineto{\pgfqpoint{2.413680in}{3.290833in}}%
\pgfpathlineto{\pgfqpoint{2.379195in}{3.301734in}}%
\pgfpathlineto{\pgfqpoint{2.344490in}{3.311422in}}%
\pgfpathlineto{\pgfqpoint{2.309600in}{3.319901in}}%
\pgfpathlineto{\pgfqpoint{2.274559in}{3.327173in}}%
\pgfpathlineto{\pgfqpoint{2.239401in}{3.333242in}}%
\pgfpathlineto{\pgfqpoint{2.204160in}{3.338111in}}%
\pgfpathlineto{\pgfqpoint{2.168867in}{3.341785in}}%
\pgfpathlineto{\pgfqpoint{2.124728in}{3.344704in}}%
\pgfpathlineto{\pgfqpoint{2.080622in}{3.345772in}}%
\pgfpathlineto{\pgfqpoint{2.036611in}{3.344999in}}%
\pgfpathlineto{\pgfqpoint{1.992755in}{3.342397in}}%
\pgfpathlineto{\pgfqpoint{1.949114in}{3.337978in}}%
\pgfpathlineto{\pgfqpoint{1.905748in}{3.331754in}}%
\pgfpathlineto{\pgfqpoint{1.862715in}{3.323739in}}%
\pgfpathlineto{\pgfqpoint{1.820073in}{3.313947in}}%
\pgfpathlineto{\pgfqpoint{1.777882in}{3.302392in}}%
\pgfpathlineto{\pgfqpoint{1.736198in}{3.289091in}}%
\pgfpathlineto{\pgfqpoint{1.695080in}{3.274060in}}%
\pgfpathlineto{\pgfqpoint{1.654586in}{3.257316in}}%
\pgfpathlineto{\pgfqpoint{1.614772in}{3.238878in}}%
\pgfpathlineto{\pgfqpoint{1.575697in}{3.218764in}}%
\pgfpathlineto{\pgfqpoint{1.537418in}{3.196995in}}%
\pgfpathlineto{\pgfqpoint{1.499994in}{3.173592in}}%
\pgfpathlineto{\pgfqpoint{1.470709in}{3.153707in}}%
\pgfpathlineto{\pgfqpoint{1.442039in}{3.132802in}}%
\pgfpathlineto{\pgfqpoint{1.414013in}{3.110889in}}%
\pgfpathlineto{\pgfqpoint{1.386662in}{3.087981in}}%
\pgfpathlineto{\pgfqpoint{1.360017in}{3.064092in}}%
\pgfpathlineto{\pgfqpoint{1.334109in}{3.039236in}}%
\pgfpathlineto{\pgfqpoint{1.308968in}{3.013426in}}%
\pgfpathlineto{\pgfqpoint{1.284627in}{2.986679in}}%
\pgfpathlineto{\pgfqpoint{1.261116in}{2.959010in}}%
\pgfpathlineto{\pgfqpoint{1.238468in}{2.930437in}}%
\pgfpathlineto{\pgfqpoint{1.216716in}{2.900976in}}%
\pgfpathlineto{\pgfqpoint{1.195892in}{2.870646in}}%
\pgfpathlineto{\pgfqpoint{1.176029in}{2.839466in}}%
\pgfpathlineto{\pgfqpoint{1.157160in}{2.807457in}}%
\pgfpathlineto{\pgfqpoint{1.139321in}{2.774640in}}%
\pgfpathlineto{\pgfqpoint{1.122544in}{2.741036in}}%
\pgfpathlineto{\pgfqpoint{1.106865in}{2.706671in}}%
\pgfpathlineto{\pgfqpoint{1.092319in}{2.671568in}}%
\pgfpathlineto{\pgfqpoint{1.078941in}{2.635754in}}%
\pgfpathlineto{\pgfqpoint{1.066767in}{2.599256in}}%
\pgfpathlineto{\pgfqpoint{1.055835in}{2.562104in}}%
\pgfpathlineto{\pgfqpoint{1.046180in}{2.524328in}}%
\pgfpathlineto{\pgfqpoint{1.037841in}{2.485960in}}%
\pgfpathlineto{\pgfqpoint{1.030854in}{2.447036in}}%
\pgfpathlineto{\pgfqpoint{1.025259in}{2.407590in}}%
\pgfpathlineto{\pgfqpoint{1.021093in}{2.367663in}}%
\pgfpathlineto{\pgfqpoint{1.018394in}{2.327294in}}%
\pgfpathlineto{\pgfqpoint{1.017203in}{2.286527in}}%
\pgfpathlineto{\pgfqpoint{1.017559in}{2.245408in}}%
\pgfpathlineto{\pgfqpoint{1.019500in}{2.203985in}}%
\pgfpathlineto{\pgfqpoint{1.023066in}{2.162310in}}%
\pgfpathlineto{\pgfqpoint{1.028296in}{2.120438in}}%
\pgfpathlineto{\pgfqpoint{1.035230in}{2.078426in}}%
\pgfpathlineto{\pgfqpoint{1.043907in}{2.036336in}}%
\pgfpathlineto{\pgfqpoint{1.054365in}{1.994235in}}%
\pgfpathlineto{\pgfqpoint{1.066641in}{1.952190in}}%
\pgfpathlineto{\pgfqpoint{1.080773in}{1.910278in}}%
\pgfpathlineto{\pgfqpoint{1.096797in}{1.868576in}}%
\pgfpathlineto{\pgfqpoint{1.114747in}{1.827167in}}%
\pgfpathlineto{\pgfqpoint{1.134655in}{1.786140in}}%
\pgfpathlineto{\pgfqpoint{1.150890in}{1.755678in}}%
\pgfpathlineto{\pgfqpoint{1.168256in}{1.725525in}}%
\pgfpathlineto{\pgfqpoint{1.186764in}{1.695725in}}%
\pgfpathlineto{\pgfqpoint{1.206425in}{1.666325in}}%
\pgfpathlineto{\pgfqpoint{1.227246in}{1.637372in}}%
\pgfpathlineto{\pgfqpoint{1.249236in}{1.608915in}}%
\pgfpathlineto{\pgfqpoint{1.272400in}{1.581007in}}%
\pgfpathlineto{\pgfqpoint{1.296743in}{1.553701in}}%
\pgfpathlineto{\pgfqpoint{1.322268in}{1.527055in}}%
\pgfpathlineto{\pgfqpoint{1.348975in}{1.501127in}}%
\pgfpathlineto{\pgfqpoint{1.376861in}{1.475978in}}%
\pgfpathlineto{\pgfqpoint{1.405923in}{1.451671in}}%
\pgfpathlineto{\pgfqpoint{1.436154in}{1.428273in}}%
\pgfpathlineto{\pgfqpoint{1.467542in}{1.405851in}}%
\pgfpathlineto{\pgfqpoint{1.500074in}{1.384477in}}%
\pgfpathlineto{\pgfqpoint{1.533733in}{1.364222in}}%
\pgfpathlineto{\pgfqpoint{1.568496in}{1.345163in}}%
\pgfpathlineto{\pgfqpoint{1.604338in}{1.327375in}}%
\pgfpathlineto{\pgfqpoint{1.641227in}{1.310939in}}%
\pgfpathlineto{\pgfqpoint{1.679128in}{1.295934in}}%
\pgfpathlineto{\pgfqpoint{1.717997in}{1.282443in}}%
\pgfpathlineto{\pgfqpoint{1.757786in}{1.270550in}}%
\pgfpathlineto{\pgfqpoint{1.798442in}{1.260340in}}%
\pgfpathlineto{\pgfqpoint{1.839901in}{1.251896in}}%
\pgfpathlineto{\pgfqpoint{1.882095in}{1.245303in}}%
\pgfpathlineto{\pgfqpoint{1.924946in}{1.240647in}}%
\pgfpathlineto{\pgfqpoint{1.968369in}{1.238010in}}%
\pgfpathlineto{\pgfqpoint{2.012269in}{1.237473in}}%
\pgfpathlineto{\pgfqpoint{2.056543in}{1.239114in}}%
\pgfpathlineto{\pgfqpoint{2.086211in}{1.241456in}}%
\pgfpathlineto{\pgfqpoint{2.115959in}{1.244820in}}%
\pgfpathlineto{\pgfqpoint{2.145751in}{1.249227in}}%
\pgfpathlineto{\pgfqpoint{2.175545in}{1.254695in}}%
\pgfpathlineto{\pgfqpoint{2.205302in}{1.261243in}}%
\pgfpathlineto{\pgfqpoint{2.234977in}{1.268884in}}%
\pgfpathlineto{\pgfqpoint{2.264525in}{1.277635in}}%
\pgfpathlineto{\pgfqpoint{2.293899in}{1.287507in}}%
\pgfpathlineto{\pgfqpoint{2.323051in}{1.298511in}}%
\pgfpathlineto{\pgfqpoint{2.351930in}{1.310654in}}%
\pgfpathlineto{\pgfqpoint{2.380484in}{1.323942in}}%
\pgfpathlineto{\pgfqpoint{2.408659in}{1.338379in}}%
\pgfpathlineto{\pgfqpoint{2.436401in}{1.353965in}}%
\pgfpathlineto{\pgfqpoint{2.463653in}{1.370697in}}%
\pgfpathlineto{\pgfqpoint{2.490358in}{1.388569in}}%
\pgfpathlineto{\pgfqpoint{2.516458in}{1.407573in}}%
\pgfpathlineto{\pgfqpoint{2.541894in}{1.427697in}}%
\pgfpathlineto{\pgfqpoint{2.566608in}{1.448923in}}%
\pgfpathlineto{\pgfqpoint{2.590539in}{1.471232in}}%
\pgfpathlineto{\pgfqpoint{2.613630in}{1.494602in}}%
\pgfpathlineto{\pgfqpoint{2.635820in}{1.519005in}}%
\pgfpathlineto{\pgfqpoint{2.657051in}{1.544409in}}%
\pgfpathlineto{\pgfqpoint{2.677267in}{1.570779in}}%
\pgfpathlineto{\pgfqpoint{2.696412in}{1.598077in}}%
\pgfpathlineto{\pgfqpoint{2.714430in}{1.626259in}}%
\pgfpathlineto{\pgfqpoint{2.731271in}{1.655278in}}%
\pgfpathlineto{\pgfqpoint{2.746884in}{1.685083in}}%
\pgfpathlineto{\pgfqpoint{2.761221in}{1.715620in}}%
\pgfpathlineto{\pgfqpoint{2.774239in}{1.746830in}}%
\pgfpathlineto{\pgfqpoint{2.785895in}{1.778654in}}%
\pgfpathlineto{\pgfqpoint{2.796154in}{1.811025in}}%
\pgfpathlineto{\pgfqpoint{2.804980in}{1.843877in}}%
\pgfpathlineto{\pgfqpoint{2.812345in}{1.877140in}}%
\pgfpathlineto{\pgfqpoint{2.818223in}{1.910743in}}%
\pgfpathlineto{\pgfqpoint{2.822595in}{1.944610in}}%
\pgfpathlineto{\pgfqpoint{2.825444in}{1.978667in}}%
\pgfpathlineto{\pgfqpoint{2.826759in}{2.012837in}}%
\pgfpathlineto{\pgfqpoint{2.826534in}{2.047044in}}%
\pgfpathlineto{\pgfqpoint{2.824770in}{2.081208in}}%
\pgfpathlineto{\pgfqpoint{2.821469in}{2.115253in}}%
\pgfpathlineto{\pgfqpoint{2.816641in}{2.149101in}}%
\pgfpathlineto{\pgfqpoint{2.810300in}{2.182677in}}%
\pgfpathlineto{\pgfqpoint{2.802466in}{2.215905in}}%
\pgfpathlineto{\pgfqpoint{2.793163in}{2.248713in}}%
\pgfpathlineto{\pgfqpoint{2.782417in}{2.281029in}}%
\pgfpathlineto{\pgfqpoint{2.770264in}{2.312785in}}%
\pgfpathlineto{\pgfqpoint{2.756738in}{2.343913in}}%
\pgfpathlineto{\pgfqpoint{2.741881in}{2.374352in}}%
\pgfpathlineto{\pgfqpoint{2.725736in}{2.404042in}}%
\pgfpathlineto{\pgfqpoint{2.708352in}{2.432925in}}%
\pgfpathlineto{\pgfqpoint{2.689778in}{2.460950in}}%
\pgfpathlineto{\pgfqpoint{2.670067in}{2.488068in}}%
\pgfpathlineto{\pgfqpoint{2.649274in}{2.514233in}}%
\pgfpathlineto{\pgfqpoint{2.627456in}{2.539404in}}%
\pgfpathlineto{\pgfqpoint{2.604671in}{2.563546in}}%
\pgfpathlineto{\pgfqpoint{2.580979in}{2.586624in}}%
\pgfpathlineto{\pgfqpoint{2.556440in}{2.608610in}}%
\pgfpathlineto{\pgfqpoint{2.531115in}{2.629480in}}%
\pgfpathlineto{\pgfqpoint{2.505065in}{2.649212in}}%
\pgfpathlineto{\pgfqpoint{2.478351in}{2.667790in}}%
\pgfpathlineto{\pgfqpoint{2.451034in}{2.685201in}}%
\pgfpathlineto{\pgfqpoint{2.423174in}{2.701435in}}%
\pgfpathlineto{\pgfqpoint{2.394830in}{2.716485in}}%
\pgfpathlineto{\pgfqpoint{2.366061in}{2.730350in}}%
\pgfpathlineto{\pgfqpoint{2.336924in}{2.743029in}}%
\pgfpathlineto{\pgfqpoint{2.307475in}{2.754526in}}%
\pgfpathlineto{\pgfqpoint{2.277767in}{2.764847in}}%
\pgfpathlineto{\pgfqpoint{2.247854in}{2.774000in}}%
\pgfpathlineto{\pgfqpoint{2.217787in}{2.781997in}}%
\pgfpathlineto{\pgfqpoint{2.187613in}{2.788850in}}%
\pgfpathlineto{\pgfqpoint{2.157381in}{2.794576in}}%
\pgfpathlineto{\pgfqpoint{2.127136in}{2.799190in}}%
\pgfpathlineto{\pgfqpoint{2.096920in}{2.802712in}}%
\pgfpathlineto{\pgfqpoint{2.066775in}{2.805162in}}%
\pgfpathlineto{\pgfqpoint{2.036741in}{2.806562in}}%
\pgfpathlineto{\pgfqpoint{2.006853in}{2.806933in}}%
\pgfpathlineto{\pgfqpoint{1.977149in}{2.806299in}}%
\pgfpathlineto{\pgfqpoint{1.933005in}{2.803519in}}%
\pgfpathlineto{\pgfqpoint{1.889450in}{2.798620in}}%
\pgfpathlineto{\pgfqpoint{1.846577in}{2.791688in}}%
\pgfpathlineto{\pgfqpoint{1.804471in}{2.782813in}}%
\pgfpathlineto{\pgfqpoint{1.763209in}{2.772087in}}%
\pgfpathlineto{\pgfqpoint{1.722858in}{2.759600in}}%
\pgfpathlineto{\pgfqpoint{1.683478in}{2.745443in}}%
\pgfpathlineto{\pgfqpoint{1.645123in}{2.729706in}}%
\pgfpathlineto{\pgfqpoint{1.607837in}{2.712477in}}%
\pgfpathlineto{\pgfqpoint{1.571660in}{2.693844in}}%
\pgfpathlineto{\pgfqpoint{1.536626in}{2.673893in}}%
\pgfpathlineto{\pgfqpoint{1.502762in}{2.652706in}}%
\pgfpathlineto{\pgfqpoint{1.470091in}{2.630365in}}%
\pgfpathlineto{\pgfqpoint{1.438632in}{2.606948in}}%
\pgfpathlineto{\pgfqpoint{1.408399in}{2.582532in}}%
\pgfpathlineto{\pgfqpoint{1.379401in}{2.557190in}}%
\pgfpathlineto{\pgfqpoint{1.351647in}{2.530994in}}%
\pgfpathlineto{\pgfqpoint{1.325138in}{2.504011in}}%
\pgfpathlineto{\pgfqpoint{1.299877in}{2.476308in}}%
\pgfpathlineto{\pgfqpoint{1.275861in}{2.447947in}}%
\pgfpathlineto{\pgfqpoint{1.253087in}{2.418989in}}%
\pgfpathlineto{\pgfqpoint{1.231547in}{2.389493in}}%
\pgfpathlineto{\pgfqpoint{1.211235in}{2.359513in}}%
\pgfpathlineto{\pgfqpoint{1.192140in}{2.329103in}}%
\pgfpathlineto{\pgfqpoint{1.174253in}{2.298313in}}%
\pgfpathlineto{\pgfqpoint{1.157559in}{2.267193in}}%
\pgfpathlineto{\pgfqpoint{1.142047in}{2.235788in}}%
\pgfpathlineto{\pgfqpoint{1.127701in}{2.204142in}}%
\pgfpathlineto{\pgfqpoint{1.114507in}{2.172298in}}%
\pgfpathlineto{\pgfqpoint{1.098679in}{2.129601in}}%
\pgfpathlineto{\pgfqpoint{1.084831in}{2.086713in}}%
\pgfpathlineto{\pgfqpoint{1.072922in}{2.043718in}}%
\pgfpathlineto{\pgfqpoint{1.062910in}{2.000697in}}%
\pgfpathlineto{\pgfqpoint{1.054751in}{1.957722in}}%
\pgfpathlineto{\pgfqpoint{1.048404in}{1.914865in}}%
\pgfpathlineto{\pgfqpoint{1.043823in}{1.872190in}}%
\pgfpathlineto{\pgfqpoint{1.040964in}{1.829759in}}%
\pgfpathlineto{\pgfqpoint{1.039783in}{1.787629in}}%
\pgfpathlineto{\pgfqpoint{1.040235in}{1.745854in}}%
\pgfpathlineto{\pgfqpoint{1.042277in}{1.704484in}}%
\pgfpathlineto{\pgfqpoint{1.045865in}{1.663567in}}%
\pgfpathlineto{\pgfqpoint{1.050955in}{1.623147in}}%
\pgfpathlineto{\pgfqpoint{1.057503in}{1.583266in}}%
\pgfpathlineto{\pgfqpoint{1.065468in}{1.543963in}}%
\pgfpathlineto{\pgfqpoint{1.074807in}{1.505275in}}%
\pgfpathlineto{\pgfqpoint{1.085478in}{1.467236in}}%
\pgfpathlineto{\pgfqpoint{1.097440in}{1.429879in}}%
\pgfpathlineto{\pgfqpoint{1.110652in}{1.393234in}}%
\pgfpathlineto{\pgfqpoint{1.125075in}{1.357330in}}%
\pgfpathlineto{\pgfqpoint{1.140669in}{1.322195in}}%
\pgfpathlineto{\pgfqpoint{1.157396in}{1.287853in}}%
\pgfpathlineto{\pgfqpoint{1.175216in}{1.254329in}}%
\pgfpathlineto{\pgfqpoint{1.194092in}{1.221645in}}%
\pgfpathlineto{\pgfqpoint{1.213987in}{1.189822in}}%
\pgfpathlineto{\pgfqpoint{1.234864in}{1.158882in}}%
\pgfpathlineto{\pgfqpoint{1.256688in}{1.128842in}}%
\pgfpathlineto{\pgfqpoint{1.279422in}{1.099721in}}%
\pgfpathlineto{\pgfqpoint{1.303031in}{1.071535in}}%
\pgfpathlineto{\pgfqpoint{1.327481in}{1.044302in}}%
\pgfpathlineto{\pgfqpoint{1.352737in}{1.018035in}}%
\pgfpathlineto{\pgfqpoint{1.378766in}{0.992750in}}%
\pgfpathlineto{\pgfqpoint{1.405534in}{0.968459in}}%
\pgfpathlineto{\pgfqpoint{1.433009in}{0.945177in}}%
\pgfpathlineto{\pgfqpoint{1.461157in}{0.922916in}}%
\pgfpathlineto{\pgfqpoint{1.489946in}{0.901686in}}%
\pgfpathlineto{\pgfqpoint{1.519345in}{0.881499in}}%
\pgfpathlineto{\pgfqpoint{1.549321in}{0.862365in}}%
\pgfpathlineto{\pgfqpoint{1.579844in}{0.844295in}}%
\pgfpathlineto{\pgfqpoint{1.618718in}{0.823217in}}%
\pgfpathlineto{\pgfqpoint{1.658336in}{0.803834in}}%
\pgfpathlineto{\pgfqpoint{1.698638in}{0.786160in}}%
\pgfpathlineto{\pgfqpoint{1.739564in}{0.770212in}}%
\pgfpathlineto{\pgfqpoint{1.781056in}{0.756004in}}%
\pgfpathlineto{\pgfqpoint{1.823054in}{0.743549in}}%
\pgfpathlineto{\pgfqpoint{1.865498in}{0.732861in}}%
\pgfpathlineto{\pgfqpoint{1.908331in}{0.723950in}}%
\pgfpathlineto{\pgfqpoint{1.951492in}{0.716827in}}%
\pgfpathlineto{\pgfqpoint{1.994924in}{0.711504in}}%
\pgfpathlineto{\pgfqpoint{2.038566in}{0.707989in}}%
\pgfpathlineto{\pgfqpoint{2.082358in}{0.706290in}}%
\pgfpathlineto{\pgfqpoint{2.126242in}{0.706416in}}%
\pgfpathlineto{\pgfqpoint{2.170156in}{0.708373in}}%
\pgfpathlineto{\pgfqpoint{2.214040in}{0.712168in}}%
\pgfpathlineto{\pgfqpoint{2.257831in}{0.717806in}}%
\pgfpathlineto{\pgfqpoint{2.292756in}{0.723647in}}%
\pgfpathlineto{\pgfqpoint{2.327550in}{0.730671in}}%
\pgfpathlineto{\pgfqpoint{2.362180in}{0.738882in}}%
\pgfpathlineto{\pgfqpoint{2.396613in}{0.748280in}}%
\pgfpathlineto{\pgfqpoint{2.430815in}{0.758865in}}%
\pgfpathlineto{\pgfqpoint{2.464753in}{0.770639in}}%
\pgfpathlineto{\pgfqpoint{2.498393in}{0.783600in}}%
\pgfpathlineto{\pgfqpoint{2.531698in}{0.797748in}}%
\pgfpathlineto{\pgfqpoint{2.564634in}{0.813083in}}%
\pgfpathlineto{\pgfqpoint{2.597164in}{0.829602in}}%
\pgfpathlineto{\pgfqpoint{2.629252in}{0.847303in}}%
\pgfpathlineto{\pgfqpoint{2.660859in}{0.866183in}}%
\pgfpathlineto{\pgfqpoint{2.691948in}{0.886238in}}%
\pgfpathlineto{\pgfqpoint{2.722479in}{0.907465in}}%
\pgfpathlineto{\pgfqpoint{2.752414in}{0.929857in}}%
\pgfpathlineto{\pgfqpoint{2.781710in}{0.953410in}}%
\pgfpathlineto{\pgfqpoint{2.809036in}{0.976963in}}%
\pgfpathlineto{\pgfqpoint{2.809036in}{0.976963in}}%
\pgfusepath{stroke}%
\end{pgfscope}%
\begin{pgfscope}%
\pgfpathrectangle{\pgfqpoint{0.472301in}{0.352669in}}{\pgfqpoint{3.351934in}{3.351934in}}%
\pgfusepath{clip}%
\pgfsetrectcap%
\pgfsetroundjoin%
\pgfsetlinewidth{0.803000pt}%
\definecolor{currentstroke}{rgb}{0.250980,0.400000,0.600000}%
\pgfsetstrokecolor{currentstroke}%
\pgfsetdash{}{0pt}%
\pgfpathmoveto{\pgfqpoint{2.910071in}{2.028636in}}%
\pgfpathlineto{\pgfqpoint{2.884443in}{2.065554in}}%
\pgfpathlineto{\pgfqpoint{2.858459in}{2.100500in}}%
\pgfpathlineto{\pgfqpoint{2.831953in}{2.133759in}}%
\pgfpathlineto{\pgfqpoint{2.805022in}{2.165255in}}%
\pgfpathlineto{\pgfqpoint{2.777778in}{2.194919in}}%
\pgfpathlineto{\pgfqpoint{2.750350in}{2.222691in}}%
\pgfpathlineto{\pgfqpoint{2.721313in}{2.249943in}}%
\pgfpathlineto{\pgfqpoint{2.692224in}{2.275150in}}%
\pgfpathlineto{\pgfqpoint{2.663256in}{2.298280in}}%
\pgfpathlineto{\pgfqpoint{2.634599in}{2.319320in}}%
\pgfpathlineto{\pgfqpoint{2.604529in}{2.339513in}}%
\pgfpathlineto{\pgfqpoint{2.575007in}{2.357538in}}%
\pgfpathlineto{\pgfqpoint{2.546254in}{2.373438in}}%
\pgfpathlineto{\pgfqpoint{2.516319in}{2.388311in}}%
\pgfpathlineto{\pgfqpoint{2.487448in}{2.401074in}}%
\pgfpathlineto{\pgfqpoint{2.457537in}{2.412694in}}%
\pgfpathlineto{\pgfqpoint{2.429011in}{2.422283in}}%
\pgfpathlineto{\pgfqpoint{2.399619in}{2.430657in}}%
\pgfpathlineto{\pgfqpoint{2.369388in}{2.437685in}}%
\pgfpathlineto{\pgfqpoint{2.340973in}{2.442829in}}%
\pgfpathlineto{\pgfqpoint{2.311927in}{2.446617in}}%
\pgfpathlineto{\pgfqpoint{2.282306in}{2.448931in}}%
\pgfpathlineto{\pgfqpoint{2.254936in}{2.449657in}}%
\pgfpathlineto{\pgfqpoint{2.227211in}{2.448976in}}%
\pgfpathlineto{\pgfqpoint{2.199206in}{2.446797in}}%
\pgfpathlineto{\pgfqpoint{2.171008in}{2.443029in}}%
\pgfpathlineto{\pgfqpoint{2.145549in}{2.438207in}}%
\pgfpathlineto{\pgfqpoint{2.120101in}{2.431965in}}%
\pgfpathlineto{\pgfqpoint{2.094756in}{2.424246in}}%
\pgfpathlineto{\pgfqpoint{2.069619in}{2.414998in}}%
\pgfpathlineto{\pgfqpoint{2.044804in}{2.404176in}}%
\pgfpathlineto{\pgfqpoint{2.023115in}{2.393203in}}%
\pgfpathlineto{\pgfqpoint{2.001872in}{2.380936in}}%
\pgfpathlineto{\pgfqpoint{1.981172in}{2.367362in}}%
\pgfpathlineto{\pgfqpoint{1.961121in}{2.352476in}}%
\pgfpathlineto{\pgfqpoint{1.941826in}{2.336281in}}%
\pgfpathlineto{\pgfqpoint{1.923402in}{2.318791in}}%
\pgfpathlineto{\pgfqpoint{1.905964in}{2.300029in}}%
\pgfpathlineto{\pgfqpoint{1.889632in}{2.280030in}}%
\pgfpathlineto{\pgfqpoint{1.876341in}{2.261553in}}%
\pgfpathlineto{\pgfqpoint{1.864069in}{2.242206in}}%
\pgfpathlineto{\pgfqpoint{1.852894in}{2.222033in}}%
\pgfpathlineto{\pgfqpoint{1.842895in}{2.201092in}}%
\pgfpathlineto{\pgfqpoint{1.834147in}{2.179444in}}%
\pgfpathlineto{\pgfqpoint{1.826724in}{2.157161in}}%
\pgfpathlineto{\pgfqpoint{1.820694in}{2.134322in}}%
\pgfpathlineto{\pgfqpoint{1.816122in}{2.111013in}}%
\pgfpathlineto{\pgfqpoint{1.813065in}{2.087328in}}%
\pgfpathlineto{\pgfqpoint{1.811577in}{2.063369in}}%
\pgfpathlineto{\pgfqpoint{1.811702in}{2.039244in}}%
\pgfpathlineto{\pgfqpoint{1.813476in}{2.015068in}}%
\pgfpathlineto{\pgfqpoint{1.816927in}{1.990960in}}%
\pgfpathlineto{\pgfqpoint{1.822074in}{1.967045in}}%
\pgfpathlineto{\pgfqpoint{1.827840in}{1.946798in}}%
\pgfpathlineto{\pgfqpoint{1.834857in}{1.926871in}}%
\pgfpathlineto{\pgfqpoint{1.843115in}{1.907347in}}%
\pgfpathlineto{\pgfqpoint{1.852602in}{1.888311in}}%
\pgfpathlineto{\pgfqpoint{1.863296in}{1.869849in}}%
\pgfpathlineto{\pgfqpoint{1.875172in}{1.852045in}}%
\pgfpathlineto{\pgfqpoint{1.888195in}{1.834983in}}%
\pgfpathlineto{\pgfqpoint{1.902325in}{1.818746in}}%
\pgfpathlineto{\pgfqpoint{1.917516in}{1.803414in}}%
\pgfpathlineto{\pgfqpoint{1.933714in}{1.789066in}}%
\pgfpathlineto{\pgfqpoint{1.950859in}{1.775776in}}%
\pgfpathlineto{\pgfqpoint{1.968885in}{1.763615in}}%
\pgfpathlineto{\pgfqpoint{1.987722in}{1.752650in}}%
\pgfpathlineto{\pgfqpoint{2.007290in}{1.742944in}}%
\pgfpathlineto{\pgfqpoint{2.027508in}{1.734555in}}%
\pgfpathlineto{\pgfqpoint{2.048287in}{1.727533in}}%
\pgfpathlineto{\pgfqpoint{2.069535in}{1.721923in}}%
\pgfpathlineto{\pgfqpoint{2.091155in}{1.717766in}}%
\pgfpathlineto{\pgfqpoint{2.113048in}{1.715091in}}%
\pgfpathlineto{\pgfqpoint{2.135110in}{1.713925in}}%
\pgfpathlineto{\pgfqpoint{2.157236in}{1.714284in}}%
\pgfpathlineto{\pgfqpoint{2.179318in}{1.716177in}}%
\pgfpathlineto{\pgfqpoint{2.201247in}{1.719605in}}%
\pgfpathlineto{\pgfqpoint{2.222914in}{1.724562in}}%
\pgfpathlineto{\pgfqpoint{2.244209in}{1.731033in}}%
\pgfpathlineto{\pgfqpoint{2.265024in}{1.738994in}}%
\pgfpathlineto{\pgfqpoint{2.285250in}{1.748413in}}%
\pgfpathlineto{\pgfqpoint{2.304783in}{1.759252in}}%
\pgfpathlineto{\pgfqpoint{2.323519in}{1.771461in}}%
\pgfpathlineto{\pgfqpoint{2.341359in}{1.784986in}}%
\pgfpathlineto{\pgfqpoint{2.358207in}{1.799763in}}%
\pgfpathlineto{\pgfqpoint{2.373970in}{1.815721in}}%
\pgfpathlineto{\pgfqpoint{2.388564in}{1.832782in}}%
\pgfpathlineto{\pgfqpoint{2.401907in}{1.850863in}}%
\pgfpathlineto{\pgfqpoint{2.413925in}{1.869873in}}%
\pgfpathlineto{\pgfqpoint{2.424550in}{1.889716in}}%
\pgfpathlineto{\pgfqpoint{2.433720in}{1.910289in}}%
\pgfpathlineto{\pgfqpoint{2.441384in}{1.931487in}}%
\pgfpathlineto{\pgfqpoint{2.447495in}{1.953199in}}%
\pgfpathlineto{\pgfqpoint{2.452016in}{1.975313in}}%
\pgfpathlineto{\pgfqpoint{2.454919in}{1.997710in}}%
\pgfpathlineto{\pgfqpoint{2.456184in}{2.020272in}}%
\pgfpathlineto{\pgfqpoint{2.455799in}{2.042880in}}%
\pgfpathlineto{\pgfqpoint{2.453764in}{2.065410in}}%
\pgfpathlineto{\pgfqpoint{2.450084in}{2.087743in}}%
\pgfpathlineto{\pgfqpoint{2.444777in}{2.109757in}}%
\pgfpathlineto{\pgfqpoint{2.437867in}{2.131332in}}%
\pgfpathlineto{\pgfqpoint{2.429388in}{2.152351in}}%
\pgfpathlineto{\pgfqpoint{2.419384in}{2.172699in}}%
\pgfpathlineto{\pgfqpoint{2.407906in}{2.192264in}}%
\pgfpathlineto{\pgfqpoint{2.395015in}{2.210938in}}%
\pgfpathlineto{\pgfqpoint{2.380778in}{2.228617in}}%
\pgfpathlineto{\pgfqpoint{2.365270in}{2.245204in}}%
\pgfpathlineto{\pgfqpoint{2.348576in}{2.260606in}}%
\pgfpathlineto{\pgfqpoint{2.330784in}{2.274737in}}%
\pgfpathlineto{\pgfqpoint{2.311991in}{2.287518in}}%
\pgfpathlineto{\pgfqpoint{2.292299in}{2.298878in}}%
\pgfpathlineto{\pgfqpoint{2.271815in}{2.308751in}}%
\pgfpathlineto{\pgfqpoint{2.250650in}{2.317082in}}%
\pgfpathlineto{\pgfqpoint{2.228922in}{2.323823in}}%
\pgfpathlineto{\pgfqpoint{2.206748in}{2.328936in}}%
\pgfpathlineto{\pgfqpoint{2.184252in}{2.332391in}}%
\pgfpathlineto{\pgfqpoint{2.161556in}{2.334167in}}%
\pgfpathlineto{\pgfqpoint{2.138787in}{2.334252in}}%
\pgfpathlineto{\pgfqpoint{2.116071in}{2.332645in}}%
\pgfpathlineto{\pgfqpoint{2.093533in}{2.329354in}}%
\pgfpathlineto{\pgfqpoint{2.071299in}{2.324395in}}%
\pgfpathlineto{\pgfqpoint{2.049492in}{2.317795in}}%
\pgfpathlineto{\pgfqpoint{2.028233in}{2.309590in}}%
\pgfpathlineto{\pgfqpoint{2.007641in}{2.299823in}}%
\pgfpathlineto{\pgfqpoint{1.987831in}{2.288548in}}%
\pgfpathlineto{\pgfqpoint{1.968914in}{2.275828in}}%
\pgfpathlineto{\pgfqpoint{1.950995in}{2.261732in}}%
\pgfpathlineto{\pgfqpoint{1.934174in}{2.246337in}}%
\pgfpathlineto{\pgfqpoint{1.918547in}{2.229730in}}%
\pgfpathlineto{\pgfqpoint{1.904200in}{2.212002in}}%
\pgfpathlineto{\pgfqpoint{1.891214in}{2.193251in}}%
\pgfpathlineto{\pgfqpoint{1.879662in}{2.173582in}}%
\pgfpathlineto{\pgfqpoint{1.869609in}{2.153103in}}%
\pgfpathlineto{\pgfqpoint{1.861113in}{2.131930in}}%
\pgfpathlineto{\pgfqpoint{1.854220in}{2.110180in}}%
\pgfpathlineto{\pgfqpoint{1.848970in}{2.087975in}}%
\pgfpathlineto{\pgfqpoint{1.845393in}{2.065437in}}%
\pgfpathlineto{\pgfqpoint{1.843509in}{2.042694in}}%
\pgfpathlineto{\pgfqpoint{1.843329in}{2.019872in}}%
\pgfpathlineto{\pgfqpoint{1.844855in}{1.997100in}}%
\pgfpathlineto{\pgfqpoint{1.848079in}{1.974504in}}%
\pgfpathlineto{\pgfqpoint{1.852984in}{1.952210in}}%
\pgfpathlineto{\pgfqpoint{1.859541in}{1.930345in}}%
\pgfpathlineto{\pgfqpoint{1.867715in}{1.909030in}}%
\pgfpathlineto{\pgfqpoint{1.877460in}{1.888385in}}%
\pgfpathlineto{\pgfqpoint{1.888722in}{1.868525in}}%
\pgfpathlineto{\pgfqpoint{1.901438in}{1.849563in}}%
\pgfpathlineto{\pgfqpoint{1.915537in}{1.831604in}}%
\pgfpathlineto{\pgfqpoint{1.930941in}{1.814749in}}%
\pgfpathlineto{\pgfqpoint{1.947563in}{1.799094in}}%
\pgfpathlineto{\pgfqpoint{1.965311in}{1.784725in}}%
\pgfpathlineto{\pgfqpoint{1.984084in}{1.771724in}}%
\pgfpathlineto{\pgfqpoint{2.003778in}{1.760164in}}%
\pgfpathlineto{\pgfqpoint{2.024282in}{1.750109in}}%
\pgfpathlineto{\pgfqpoint{2.045482in}{1.741617in}}%
\pgfpathlineto{\pgfqpoint{2.067258in}{1.734735in}}%
\pgfpathlineto{\pgfqpoint{2.089489in}{1.729503in}}%
\pgfpathlineto{\pgfqpoint{2.112050in}{1.725948in}}%
\pgfpathlineto{\pgfqpoint{2.134814in}{1.724093in}}%
\pgfpathlineto{\pgfqpoint{2.157654in}{1.723946in}}%
\pgfpathlineto{\pgfqpoint{2.180441in}{1.725510in}}%
\pgfpathlineto{\pgfqpoint{2.203047in}{1.728776in}}%
\pgfpathlineto{\pgfqpoint{2.225346in}{1.733725in}}%
\pgfpathlineto{\pgfqpoint{2.247212in}{1.740329in}}%
\pgfpathlineto{\pgfqpoint{2.268522in}{1.748553in}}%
\pgfpathlineto{\pgfqpoint{2.289158in}{1.758349in}}%
\pgfpathlineto{\pgfqpoint{2.309002in}{1.769663in}}%
\pgfpathlineto{\pgfqpoint{2.327943in}{1.782431in}}%
\pgfpathlineto{\pgfqpoint{2.345875in}{1.796582in}}%
\pgfpathlineto{\pgfqpoint{2.362698in}{1.812036in}}%
\pgfpathlineto{\pgfqpoint{2.378316in}{1.828708in}}%
\pgfpathlineto{\pgfqpoint{2.392642in}{1.846502in}}%
\pgfpathlineto{\pgfqpoint{2.405595in}{1.865319in}}%
\pgfpathlineto{\pgfqpoint{2.417102in}{1.885053in}}%
\pgfpathlineto{\pgfqpoint{2.427099in}{1.905595in}}%
\pgfpathlineto{\pgfqpoint{2.435530in}{1.926827in}}%
\pgfpathlineto{\pgfqpoint{2.442347in}{1.948632in}}%
\pgfpathlineto{\pgfqpoint{2.447511in}{1.970886in}}%
\pgfpathlineto{\pgfqpoint{2.450994in}{1.993464in}}%
\pgfpathlineto{\pgfqpoint{2.452776in}{2.016240in}}%
\pgfpathlineto{\pgfqpoint{2.452847in}{2.039086in}}%
\pgfpathlineto{\pgfqpoint{2.451206in}{2.061873in}}%
\pgfpathlineto{\pgfqpoint{2.447863in}{2.084473in}}%
\pgfpathlineto{\pgfqpoint{2.442837in}{2.106760in}}%
\pgfpathlineto{\pgfqpoint{2.436155in}{2.128607in}}%
\pgfpathlineto{\pgfqpoint{2.427855in}{2.149893in}}%
\pgfpathlineto{\pgfqpoint{2.417984in}{2.170497in}}%
\pgfpathlineto{\pgfqpoint{2.406597in}{2.190304in}}%
\pgfpathlineto{\pgfqpoint{2.393758in}{2.209202in}}%
\pgfpathlineto{\pgfqpoint{2.379540in}{2.227086in}}%
\pgfpathlineto{\pgfqpoint{2.364021in}{2.243854in}}%
\pgfpathlineto{\pgfqpoint{2.347290in}{2.259412in}}%
\pgfpathlineto{\pgfqpoint{2.329441in}{2.273673in}}%
\pgfpathlineto{\pgfqpoint{2.310573in}{2.286558in}}%
\pgfpathlineto{\pgfqpoint{2.290792in}{2.297992in}}%
\pgfpathlineto{\pgfqpoint{2.270211in}{2.307912in}}%
\pgfpathlineto{\pgfqpoint{2.248944in}{2.316262in}}%
\pgfpathlineto{\pgfqpoint{2.227111in}{2.322996in}}%
\pgfpathlineto{\pgfqpoint{2.204835in}{2.328075in}}%
\pgfpathlineto{\pgfqpoint{2.182242in}{2.331471in}}%
\pgfpathlineto{\pgfqpoint{2.159457in}{2.333165in}}%
\pgfpathlineto{\pgfqpoint{2.136609in}{2.333146in}}%
\pgfpathlineto{\pgfqpoint{2.113827in}{2.331416in}}%
\pgfpathlineto{\pgfqpoint{2.091238in}{2.327984in}}%
\pgfpathlineto{\pgfqpoint{2.068970in}{2.322869in}}%
\pgfpathlineto{\pgfqpoint{2.047148in}{2.316100in}}%
\pgfpathlineto{\pgfqpoint{2.025894in}{2.307715in}}%
\pgfpathlineto{\pgfqpoint{2.005328in}{2.297762in}}%
\pgfpathlineto{\pgfqpoint{1.985565in}{2.286295in}}%
\pgfpathlineto{\pgfqpoint{1.966718in}{2.273380in}}%
\pgfpathlineto{\pgfqpoint{1.948891in}{2.259089in}}%
\pgfpathlineto{\pgfqpoint{1.932185in}{2.243502in}}%
\pgfpathlineto{\pgfqpoint{1.916693in}{2.226708in}}%
\pgfpathlineto{\pgfqpoint{1.902503in}{2.208800in}}%
\pgfpathlineto{\pgfqpoint{1.889696in}{2.189879in}}%
\pgfpathlineto{\pgfqpoint{1.878341in}{2.170051in}}%
\pgfpathlineto{\pgfqpoint{1.868504in}{2.149429in}}%
\pgfpathlineto{\pgfqpoint{1.860241in}{2.128127in}}%
\pgfpathlineto{\pgfqpoint{1.853596in}{2.106267in}}%
\pgfpathlineto{\pgfqpoint{1.848608in}{2.083969in}}%
\pgfpathlineto{\pgfqpoint{1.845304in}{2.061361in}}%
\pgfpathlineto{\pgfqpoint{1.843704in}{2.038569in}}%
\pgfpathlineto{\pgfqpoint{1.843816in}{2.015720in}}%
\pgfpathlineto{\pgfqpoint{1.845640in}{1.992945in}}%
\pgfpathlineto{\pgfqpoint{1.849165in}{1.970370in}}%
\pgfpathlineto{\pgfqpoint{1.854372in}{1.948123in}}%
\pgfpathlineto{\pgfqpoint{1.861231in}{1.926328in}}%
\pgfpathlineto{\pgfqpoint{1.869704in}{1.905109in}}%
\pgfpathlineto{\pgfqpoint{1.879743in}{1.884584in}}%
\pgfpathlineto{\pgfqpoint{1.891291in}{1.864868in}}%
\pgfpathlineto{\pgfqpoint{1.904285in}{1.846074in}}%
\pgfpathlineto{\pgfqpoint{1.918650in}{1.828306in}}%
\pgfpathlineto{\pgfqpoint{1.934306in}{1.811664in}}%
\pgfpathlineto{\pgfqpoint{1.951164in}{1.796243in}}%
\pgfpathlineto{\pgfqpoint{1.969132in}{1.782127in}}%
\pgfpathlineto{\pgfqpoint{1.988106in}{1.769398in}}%
\pgfpathlineto{\pgfqpoint{2.007980in}{1.758126in}}%
\pgfpathlineto{\pgfqpoint{2.028644in}{1.748375in}}%
\pgfpathlineto{\pgfqpoint{2.049980in}{1.740200in}}%
\pgfpathlineto{\pgfqpoint{2.071868in}{1.733646in}}%
\pgfpathlineto{\pgfqpoint{2.094186in}{1.728751in}}%
\pgfpathlineto{\pgfqpoint{2.116809in}{1.725541in}}%
\pgfpathlineto{\pgfqpoint{2.139608in}{1.724036in}}%
\pgfpathlineto{\pgfqpoint{2.162455in}{1.724244in}}%
\pgfpathlineto{\pgfqpoint{2.185223in}{1.726163in}}%
\pgfpathlineto{\pgfqpoint{2.207783in}{1.729782in}}%
\pgfpathlineto{\pgfqpoint{2.230009in}{1.735082in}}%
\pgfpathlineto{\pgfqpoint{2.251775in}{1.742032in}}%
\pgfpathlineto{\pgfqpoint{2.272959in}{1.750593in}}%
\pgfpathlineto{\pgfqpoint{2.293442in}{1.760718in}}%
\pgfpathlineto{\pgfqpoint{2.313109in}{1.772349in}}%
\pgfpathlineto{\pgfqpoint{2.331849in}{1.785421in}}%
\pgfpathlineto{\pgfqpoint{2.349556in}{1.799860in}}%
\pgfpathlineto{\pgfqpoint{2.366133in}{1.815586in}}%
\pgfpathlineto{\pgfqpoint{2.381484in}{1.832509in}}%
\pgfpathlineto{\pgfqpoint{2.395524in}{1.850535in}}%
\pgfpathlineto{\pgfqpoint{2.408174in}{1.869563in}}%
\pgfpathlineto{\pgfqpoint{2.419362in}{1.889484in}}%
\pgfpathlineto{\pgfqpoint{2.429027in}{1.910189in}}%
\pgfpathlineto{\pgfqpoint{2.437113in}{1.931559in}}%
\pgfpathlineto{\pgfqpoint{2.443575in}{1.953475in}}%
\pgfpathlineto{\pgfqpoint{2.448376in}{1.975813in}}%
\pgfpathlineto{\pgfqpoint{2.451491in}{1.998448in}}%
\pgfpathlineto{\pgfqpoint{2.452900in}{2.021254in}}%
\pgfpathlineto{\pgfqpoint{2.452597in}{2.044100in}}%
\pgfpathlineto{\pgfqpoint{2.450583in}{2.066860in}}%
\pgfpathlineto{\pgfqpoint{2.446869in}{2.089405in}}%
\pgfpathlineto{\pgfqpoint{2.441476in}{2.111608in}}%
\pgfpathlineto{\pgfqpoint{2.434434in}{2.133345in}}%
\pgfpathlineto{\pgfqpoint{2.425784in}{2.154493in}}%
\pgfpathlineto{\pgfqpoint{2.415574in}{2.174933in}}%
\pgfpathlineto{\pgfqpoint{2.403860in}{2.194551in}}%
\pgfpathlineto{\pgfqpoint{2.390710in}{2.213236in}}%
\pgfpathlineto{\pgfqpoint{2.376196in}{2.230883in}}%
\pgfpathlineto{\pgfqpoint{2.360401in}{2.247393in}}%
\pgfpathlineto{\pgfqpoint{2.343413in}{2.262673in}}%
\pgfpathlineto{\pgfqpoint{2.325328in}{2.276637in}}%
\pgfpathlineto{\pgfqpoint{2.306248in}{2.289207in}}%
\pgfpathlineto{\pgfqpoint{2.286279in}{2.300312in}}%
\pgfpathlineto{\pgfqpoint{2.265535in}{2.309890in}}%
\pgfpathlineto{\pgfqpoint{2.244131in}{2.317886in}}%
\pgfpathlineto{\pgfqpoint{2.222188in}{2.324256in}}%
\pgfpathlineto{\pgfqpoint{2.199830in}{2.328963in}}%
\pgfpathlineto{\pgfqpoint{2.177181in}{2.331983in}}%
\pgfpathlineto{\pgfqpoint{2.154370in}{2.333296in}}%
\pgfpathlineto{\pgfqpoint{2.131525in}{2.332897in}}%
\pgfpathlineto{\pgfqpoint{2.108774in}{2.330787in}}%
\pgfpathlineto{\pgfqpoint{2.086245in}{2.326979in}}%
\pgfpathlineto{\pgfqpoint{2.064065in}{2.321493in}}%
\pgfpathlineto{\pgfqpoint{2.042358in}{2.314360in}}%
\pgfpathlineto{\pgfqpoint{2.021246in}{2.305621in}}%
\pgfpathlineto{\pgfqpoint{2.000849in}{2.295325in}}%
\pgfpathlineto{\pgfqpoint{1.981280in}{2.283529in}}%
\pgfpathlineto{\pgfqpoint{1.962651in}{2.270300in}}%
\pgfpathlineto{\pgfqpoint{1.945065in}{2.255712in}}%
\pgfpathlineto{\pgfqpoint{1.928621in}{2.239848in}}%
\pgfpathlineto{\pgfqpoint{1.913413in}{2.222796in}}%
\pgfpathlineto{\pgfqpoint{1.899525in}{2.204653in}}%
\pgfpathlineto{\pgfqpoint{1.887035in}{2.185520in}}%
\pgfpathlineto{\pgfqpoint{1.876014in}{2.165505in}}%
\pgfpathlineto{\pgfqpoint{1.866524in}{2.144720in}}%
\pgfpathlineto{\pgfqpoint{1.858618in}{2.123283in}}%
\pgfpathlineto{\pgfqpoint{1.852340in}{2.101313in}}%
\pgfpathlineto{\pgfqpoint{1.847726in}{2.078935in}}%
\pgfpathlineto{\pgfqpoint{1.844802in}{2.056274in}}%
\pgfpathlineto{\pgfqpoint{1.843585in}{2.033458in}}%
\pgfpathlineto{\pgfqpoint{1.844080in}{2.010615in}}%
\pgfpathlineto{\pgfqpoint{1.846285in}{1.987873in}}%
\pgfpathlineto{\pgfqpoint{1.850189in}{1.965360in}}%
\pgfpathlineto{\pgfqpoint{1.855768in}{1.943203in}}%
\pgfpathlineto{\pgfqpoint{1.862992in}{1.921526in}}%
\pgfpathlineto{\pgfqpoint{1.871819in}{1.900451in}}%
\pgfpathlineto{\pgfqpoint{1.882201in}{1.880098in}}%
\pgfpathlineto{\pgfqpoint{1.894079in}{1.860579in}}%
\pgfpathlineto{\pgfqpoint{1.907386in}{1.842005in}}%
\pgfpathlineto{\pgfqpoint{1.922048in}{1.824480in}}%
\pgfpathlineto{\pgfqpoint{1.937981in}{1.808104in}}%
\pgfpathlineto{\pgfqpoint{1.955096in}{1.792967in}}%
\pgfpathlineto{\pgfqpoint{1.973298in}{1.779155in}}%
\pgfpathlineto{\pgfqpoint{1.992483in}{1.766746in}}%
\pgfpathlineto{\pgfqpoint{2.012545in}{1.755809in}}%
\pgfpathlineto{\pgfqpoint{2.033369in}{1.746407in}}%
\pgfpathlineto{\pgfqpoint{2.054839in}{1.738591in}}%
\pgfpathlineto{\pgfqpoint{2.076835in}{1.732406in}}%
\pgfpathlineto{\pgfqpoint{2.099232in}{1.727886in}}%
\pgfpathlineto{\pgfqpoint{2.121905in}{1.725057in}}%
\pgfpathlineto{\pgfqpoint{2.144726in}{1.723935in}}%
\pgfpathlineto{\pgfqpoint{2.167568in}{1.724526in}}%
\pgfpathlineto{\pgfqpoint{2.190300in}{1.726827in}}%
\pgfpathlineto{\pgfqpoint{2.212796in}{1.730825in}}%
\pgfpathlineto{\pgfqpoint{2.234930in}{1.736498in}}%
\pgfpathlineto{\pgfqpoint{2.256576in}{1.743812in}}%
\pgfpathlineto{\pgfqpoint{2.260128in}{1.745189in}}%
\pgfpathlineto{\pgfqpoint{2.260128in}{1.745189in}}%
\pgfusepath{stroke}%
\end{pgfscope}%
\begin{pgfscope}%
\pgfpathrectangle{\pgfqpoint{0.472301in}{0.352669in}}{\pgfqpoint{3.351934in}{3.351934in}}%
\pgfusepath{clip}%
\pgfsetrectcap%
\pgfsetroundjoin%
\pgfsetlinewidth{0.803000pt}%
\definecolor{currentstroke}{rgb}{0.619608,0.188235,0.172549}%
\pgfsetstrokecolor{currentstroke}%
\pgfsetdash{}{0pt}%
\pgfpathmoveto{\pgfqpoint{3.671874in}{2.028636in}}%
\pgfpathlineto{\pgfqpoint{3.649489in}{2.068293in}}%
\pgfpathlineto{\pgfqpoint{3.625959in}{2.106999in}}%
\pgfpathlineto{\pgfqpoint{3.600738in}{2.145615in}}%
\pgfpathlineto{\pgfqpoint{3.574277in}{2.183391in}}%
\pgfpathlineto{\pgfqpoint{3.546616in}{2.220285in}}%
\pgfpathlineto{\pgfqpoint{3.517797in}{2.256252in}}%
\pgfpathlineto{\pgfqpoint{3.487863in}{2.291250in}}%
\pgfpathlineto{\pgfqpoint{3.456859in}{2.325237in}}%
\pgfpathlineto{\pgfqpoint{3.424830in}{2.358171in}}%
\pgfpathlineto{\pgfqpoint{3.391825in}{2.390012in}}%
\pgfpathlineto{\pgfqpoint{3.357894in}{2.420721in}}%
\pgfpathlineto{\pgfqpoint{3.323090in}{2.450259in}}%
\pgfpathlineto{\pgfqpoint{3.287465in}{2.478589in}}%
\pgfpathlineto{\pgfqpoint{3.251078in}{2.505675in}}%
\pgfpathlineto{\pgfqpoint{3.213106in}{2.532075in}}%
\pgfpathlineto{\pgfqpoint{3.174419in}{2.557128in}}%
\pgfpathlineto{\pgfqpoint{3.135079in}{2.580800in}}%
\pgfpathlineto{\pgfqpoint{3.095154in}{2.603055in}}%
\pgfpathlineto{\pgfqpoint{3.054712in}{2.623864in}}%
\pgfpathlineto{\pgfqpoint{3.013826in}{2.643197in}}%
\pgfpathlineto{\pgfqpoint{2.972571in}{2.661027in}}%
\pgfpathlineto{\pgfqpoint{2.931026in}{2.677332in}}%
\pgfpathlineto{\pgfqpoint{2.889272in}{2.692093in}}%
\pgfpathlineto{\pgfqpoint{2.847396in}{2.705296in}}%
\pgfpathlineto{\pgfqpoint{2.805484in}{2.716929in}}%
\pgfpathlineto{\pgfqpoint{2.763629in}{2.726987in}}%
\pgfpathlineto{\pgfqpoint{2.721924in}{2.735471in}}%
\pgfpathlineto{\pgfqpoint{2.680466in}{2.742389in}}%
\pgfpathlineto{\pgfqpoint{2.639354in}{2.747753in}}%
\pgfpathlineto{\pgfqpoint{2.598688in}{2.751586in}}%
\pgfpathlineto{\pgfqpoint{2.558571in}{2.753916in}}%
\pgfpathlineto{\pgfqpoint{2.517629in}{2.754786in}}%
\pgfpathlineto{\pgfqpoint{2.477388in}{2.754128in}}%
\pgfpathlineto{\pgfqpoint{2.437956in}{2.751992in}}%
\pgfpathlineto{\pgfqpoint{2.399440in}{2.748439in}}%
\pgfpathlineto{\pgfqpoint{2.361945in}{2.743539in}}%
\pgfpathlineto{\pgfqpoint{2.323984in}{2.737075in}}%
\pgfpathlineto{\pgfqpoint{2.287213in}{2.729305in}}%
\pgfpathlineto{\pgfqpoint{2.251736in}{2.720332in}}%
\pgfpathlineto{\pgfqpoint{2.216019in}{2.709749in}}%
\pgfpathlineto{\pgfqpoint{2.181766in}{2.698049in}}%
\pgfpathlineto{\pgfqpoint{2.149066in}{2.685365in}}%
\pgfpathlineto{\pgfqpoint{2.116366in}{2.671090in}}%
\pgfpathlineto{\pgfqpoint{2.085378in}{2.655971in}}%
\pgfpathlineto{\pgfqpoint{2.054551in}{2.639255in}}%
\pgfpathlineto{\pgfqpoint{2.025584in}{2.621872in}}%
\pgfpathlineto{\pgfqpoint{1.996943in}{2.602912in}}%
\pgfpathlineto{\pgfqpoint{1.970291in}{2.583496in}}%
\pgfpathlineto{\pgfqpoint{1.944129in}{2.562555in}}%
\pgfpathlineto{\pgfqpoint{1.920057in}{2.541405in}}%
\pgfpathlineto{\pgfqpoint{1.896626in}{2.518814in}}%
\pgfpathlineto{\pgfqpoint{1.875349in}{2.496292in}}%
\pgfpathlineto{\pgfqpoint{1.854848in}{2.472442in}}%
\pgfpathlineto{\pgfqpoint{1.835240in}{2.447242in}}%
\pgfpathlineto{\pgfqpoint{1.817858in}{2.422492in}}%
\pgfpathlineto{\pgfqpoint{1.801479in}{2.396549in}}%
\pgfpathlineto{\pgfqpoint{1.787274in}{2.371392in}}%
\pgfpathlineto{\pgfqpoint{1.774141in}{2.345217in}}%
\pgfpathlineto{\pgfqpoint{1.762192in}{2.318042in}}%
\pgfpathlineto{\pgfqpoint{1.752313in}{2.292093in}}%
\pgfpathlineto{\pgfqpoint{1.743637in}{2.265345in}}%
\pgfpathlineto{\pgfqpoint{1.736265in}{2.237835in}}%
\pgfpathlineto{\pgfqpoint{1.730299in}{2.209611in}}%
\pgfpathlineto{\pgfqpoint{1.726158in}{2.183163in}}%
\pgfpathlineto{\pgfqpoint{1.723371in}{2.156221in}}%
\pgfpathlineto{\pgfqpoint{1.722023in}{2.128846in}}%
\pgfpathlineto{\pgfqpoint{1.722198in}{2.101115in}}%
\pgfpathlineto{\pgfqpoint{1.723979in}{2.073111in}}%
\pgfpathlineto{\pgfqpoint{1.727061in}{2.047498in}}%
\pgfpathlineto{\pgfqpoint{1.731597in}{2.021821in}}%
\pgfpathlineto{\pgfqpoint{1.737642in}{1.996168in}}%
\pgfpathlineto{\pgfqpoint{1.745248in}{1.970640in}}%
\pgfpathlineto{\pgfqpoint{1.754463in}{1.945345in}}%
\pgfpathlineto{\pgfqpoint{1.764164in}{1.922874in}}%
\pgfpathlineto{\pgfqpoint{1.775226in}{1.900778in}}%
\pgfpathlineto{\pgfqpoint{1.787666in}{1.879156in}}%
\pgfpathlineto{\pgfqpoint{1.801497in}{1.858111in}}%
\pgfpathlineto{\pgfqpoint{1.816725in}{1.837753in}}%
\pgfpathlineto{\pgfqpoint{1.833345in}{1.818198in}}%
\pgfpathlineto{\pgfqpoint{1.849277in}{1.801587in}}%
\pgfpathlineto{\pgfqpoint{1.866285in}{1.785793in}}%
\pgfpathlineto{\pgfqpoint{1.884345in}{1.770906in}}%
\pgfpathlineto{\pgfqpoint{1.903428in}{1.757018in}}%
\pgfpathlineto{\pgfqpoint{1.923495in}{1.744224in}}%
\pgfpathlineto{\pgfqpoint{1.944501in}{1.732618in}}%
\pgfpathlineto{\pgfqpoint{1.966390in}{1.722296in}}%
\pgfpathlineto{\pgfqpoint{1.989097in}{1.713353in}}%
\pgfpathlineto{\pgfqpoint{2.012547in}{1.705884in}}%
\pgfpathlineto{\pgfqpoint{2.033608in}{1.700629in}}%
\pgfpathlineto{\pgfqpoint{2.055109in}{1.696634in}}%
\pgfpathlineto{\pgfqpoint{2.076979in}{1.693956in}}%
\pgfpathlineto{\pgfqpoint{2.099141in}{1.692651in}}%
\pgfpathlineto{\pgfqpoint{2.121511in}{1.692771in}}%
\pgfpathlineto{\pgfqpoint{2.143998in}{1.694365in}}%
\pgfpathlineto{\pgfqpoint{2.166503in}{1.697477in}}%
\pgfpathlineto{\pgfqpoint{2.188923in}{1.702145in}}%
\pgfpathlineto{\pgfqpoint{2.211143in}{1.708401in}}%
\pgfpathlineto{\pgfqpoint{2.233046in}{1.716273in}}%
\pgfpathlineto{\pgfqpoint{2.254505in}{1.725779in}}%
\pgfpathlineto{\pgfqpoint{2.272444in}{1.735234in}}%
\pgfpathlineto{\pgfqpoint{2.289871in}{1.745898in}}%
\pgfpathlineto{\pgfqpoint{2.306695in}{1.757765in}}%
\pgfpathlineto{\pgfqpoint{2.322821in}{1.770824in}}%
\pgfpathlineto{\pgfqpoint{2.338151in}{1.785059in}}%
\pgfpathlineto{\pgfqpoint{2.352583in}{1.800444in}}%
\pgfpathlineto{\pgfqpoint{2.366013in}{1.816947in}}%
\pgfpathlineto{\pgfqpoint{2.378334in}{1.834527in}}%
\pgfpathlineto{\pgfqpoint{2.389436in}{1.853134in}}%
\pgfpathlineto{\pgfqpoint{2.399208in}{1.872709in}}%
\pgfpathlineto{\pgfqpoint{2.407535in}{1.893180in}}%
\pgfpathlineto{\pgfqpoint{2.414302in}{1.914466in}}%
\pgfpathlineto{\pgfqpoint{2.418665in}{1.932757in}}%
\pgfpathlineto{\pgfqpoint{2.421795in}{1.951485in}}%
\pgfpathlineto{\pgfqpoint{2.423624in}{1.970579in}}%
\pgfpathlineto{\pgfqpoint{2.424084in}{1.989960in}}%
\pgfpathlineto{\pgfqpoint{2.423108in}{2.009543in}}%
\pgfpathlineto{\pgfqpoint{2.420628in}{2.029229in}}%
\pgfpathlineto{\pgfqpoint{2.416579in}{2.048911in}}%
\pgfpathlineto{\pgfqpoint{2.410897in}{2.068468in}}%
\pgfpathlineto{\pgfqpoint{2.403520in}{2.087766in}}%
\pgfpathlineto{\pgfqpoint{2.394390in}{2.106657in}}%
\pgfpathlineto{\pgfqpoint{2.385787in}{2.121365in}}%
\pgfpathlineto{\pgfqpoint{2.376002in}{2.135610in}}%
\pgfpathlineto{\pgfqpoint{2.365012in}{2.149289in}}%
\pgfpathlineto{\pgfqpoint{2.352799in}{2.162285in}}%
\pgfpathlineto{\pgfqpoint{2.339349in}{2.174471in}}%
\pgfpathlineto{\pgfqpoint{2.324652in}{2.185703in}}%
\pgfpathlineto{\pgfqpoint{2.308707in}{2.195823in}}%
\pgfpathlineto{\pgfqpoint{2.291523in}{2.204647in}}%
\pgfpathlineto{\pgfqpoint{2.273121in}{2.211971in}}%
\pgfpathlineto{\pgfqpoint{2.258541in}{2.216336in}}%
\pgfpathlineto{\pgfqpoint{2.243322in}{2.219609in}}%
\pgfpathlineto{\pgfqpoint{2.227495in}{2.221659in}}%
\pgfpathlineto{\pgfqpoint{2.211109in}{2.222339in}}%
\pgfpathlineto{\pgfqpoint{2.194224in}{2.221477in}}%
\pgfpathlineto{\pgfqpoint{2.176928in}{2.218873in}}%
\pgfpathlineto{\pgfqpoint{2.159338in}{2.214287in}}%
\pgfpathlineto{\pgfqpoint{2.147527in}{2.209988in}}%
\pgfpathlineto{\pgfqpoint{2.135721in}{2.204580in}}%
\pgfpathlineto{\pgfqpoint{2.124005in}{2.197942in}}%
\pgfpathlineto{\pgfqpoint{2.112489in}{2.189933in}}%
\pgfpathlineto{\pgfqpoint{2.101318in}{2.180382in}}%
\pgfpathlineto{\pgfqpoint{2.091505in}{2.170036in}}%
\pgfpathlineto{\pgfqpoint{2.091505in}{2.170036in}}%
\pgfusepath{stroke}%
\end{pgfscope}%
\begin{pgfscope}%
\pgfsetrectcap%
\pgfsetmiterjoin%
\pgfsetlinewidth{0.602250pt}%
\definecolor{currentstroke}{rgb}{0.000000,0.000000,0.000000}%
\pgfsetstrokecolor{currentstroke}%
\pgfsetdash{}{0pt}%
\pgfpathmoveto{\pgfqpoint{0.472301in}{0.352669in}}%
\pgfpathlineto{\pgfqpoint{0.472301in}{3.704603in}}%
\pgfusepath{stroke}%
\end{pgfscope}%
\begin{pgfscope}%
\pgfsetrectcap%
\pgfsetmiterjoin%
\pgfsetlinewidth{0.602250pt}%
\definecolor{currentstroke}{rgb}{0.000000,0.000000,0.000000}%
\pgfsetstrokecolor{currentstroke}%
\pgfsetdash{}{0pt}%
\pgfpathmoveto{\pgfqpoint{0.472301in}{0.352669in}}%
\pgfpathlineto{\pgfqpoint{3.824235in}{0.352669in}}%
\pgfusepath{stroke}%
\end{pgfscope}%
\begin{pgfscope}%
\pgfpathrectangle{\pgfqpoint{0.472301in}{0.352669in}}{\pgfqpoint{3.351934in}{3.351934in}}%
\pgfusepath{clip}%
\pgfsetbuttcap%
\pgfsetroundjoin%
\definecolor{currentfill}{rgb}{0.768627,0.556863,0.211765}%
\pgfsetfillcolor{currentfill}%
\pgfsetlinewidth{1.003750pt}%
\definecolor{currentstroke}{rgb}{0.768627,0.556863,0.211765}%
\pgfsetstrokecolor{currentstroke}%
\pgfsetdash{}{0pt}%
\pgfsys@defobject{currentmarker}{\pgfqpoint{-0.017361in}{-0.017361in}}{\pgfqpoint{0.017361in}{0.017361in}}{%
\pgfpathmoveto{\pgfqpoint{0.000000in}{-0.017361in}}%
\pgfpathcurveto{\pgfqpoint{0.004604in}{-0.017361in}}{\pgfqpoint{0.009020in}{-0.015532in}}{\pgfqpoint{0.012276in}{-0.012276in}}%
\pgfpathcurveto{\pgfqpoint{0.015532in}{-0.009020in}}{\pgfqpoint{0.017361in}{-0.004604in}}{\pgfqpoint{0.017361in}{0.000000in}}%
\pgfpathcurveto{\pgfqpoint{0.017361in}{0.004604in}}{\pgfqpoint{0.015532in}{0.009020in}}{\pgfqpoint{0.012276in}{0.012276in}}%
\pgfpathcurveto{\pgfqpoint{0.009020in}{0.015532in}}{\pgfqpoint{0.004604in}{0.017361in}}{\pgfqpoint{0.000000in}{0.017361in}}%
\pgfpathcurveto{\pgfqpoint{-0.004604in}{0.017361in}}{\pgfqpoint{-0.009020in}{0.015532in}}{\pgfqpoint{-0.012276in}{0.012276in}}%
\pgfpathcurveto{\pgfqpoint{-0.015532in}{0.009020in}}{\pgfqpoint{-0.017361in}{0.004604in}}{\pgfqpoint{-0.017361in}{0.000000in}}%
\pgfpathcurveto{\pgfqpoint{-0.017361in}{-0.004604in}}{\pgfqpoint{-0.015532in}{-0.009020in}}{\pgfqpoint{-0.012276in}{-0.012276in}}%
\pgfpathcurveto{\pgfqpoint{-0.009020in}{-0.015532in}}{\pgfqpoint{-0.004604in}{-0.017361in}}{\pgfqpoint{0.000000in}{-0.017361in}}%
\pgfpathlineto{\pgfqpoint{0.000000in}{-0.017361in}}%
\pgfpathclose%
\pgfusepath{stroke,fill}%
}%
\begin{pgfscope}%
\pgfsys@transformshift{3.474849in}{2.028636in}%
\pgfsys@useobject{currentmarker}{}%
\end{pgfscope}%
\end{pgfscope}%
\begin{pgfscope}%
\pgfpathrectangle{\pgfqpoint{0.472301in}{0.352669in}}{\pgfqpoint{3.351934in}{3.351934in}}%
\pgfusepath{clip}%
\pgfsetbuttcap%
\pgfsetroundjoin%
\definecolor{currentfill}{rgb}{0.250980,0.400000,0.600000}%
\pgfsetfillcolor{currentfill}%
\pgfsetlinewidth{1.003750pt}%
\definecolor{currentstroke}{rgb}{0.250980,0.400000,0.600000}%
\pgfsetstrokecolor{currentstroke}%
\pgfsetdash{}{0pt}%
\pgfsys@defobject{currentmarker}{\pgfqpoint{-0.017361in}{-0.017361in}}{\pgfqpoint{0.017361in}{0.017361in}}{%
\pgfpathmoveto{\pgfqpoint{0.000000in}{-0.017361in}}%
\pgfpathcurveto{\pgfqpoint{0.004604in}{-0.017361in}}{\pgfqpoint{0.009020in}{-0.015532in}}{\pgfqpoint{0.012276in}{-0.012276in}}%
\pgfpathcurveto{\pgfqpoint{0.015532in}{-0.009020in}}{\pgfqpoint{0.017361in}{-0.004604in}}{\pgfqpoint{0.017361in}{0.000000in}}%
\pgfpathcurveto{\pgfqpoint{0.017361in}{0.004604in}}{\pgfqpoint{0.015532in}{0.009020in}}{\pgfqpoint{0.012276in}{0.012276in}}%
\pgfpathcurveto{\pgfqpoint{0.009020in}{0.015532in}}{\pgfqpoint{0.004604in}{0.017361in}}{\pgfqpoint{0.000000in}{0.017361in}}%
\pgfpathcurveto{\pgfqpoint{-0.004604in}{0.017361in}}{\pgfqpoint{-0.009020in}{0.015532in}}{\pgfqpoint{-0.012276in}{0.012276in}}%
\pgfpathcurveto{\pgfqpoint{-0.015532in}{0.009020in}}{\pgfqpoint{-0.017361in}{0.004604in}}{\pgfqpoint{-0.017361in}{0.000000in}}%
\pgfpathcurveto{\pgfqpoint{-0.017361in}{-0.004604in}}{\pgfqpoint{-0.015532in}{-0.009020in}}{\pgfqpoint{-0.012276in}{-0.012276in}}%
\pgfpathcurveto{\pgfqpoint{-0.009020in}{-0.015532in}}{\pgfqpoint{-0.004604in}{-0.017361in}}{\pgfqpoint{0.000000in}{-0.017361in}}%
\pgfpathlineto{\pgfqpoint{0.000000in}{-0.017361in}}%
\pgfpathclose%
\pgfusepath{stroke,fill}%
}%
\begin{pgfscope}%
\pgfsys@transformshift{2.910071in}{2.028636in}%
\pgfsys@useobject{currentmarker}{}%
\end{pgfscope}%
\end{pgfscope}%
\begin{pgfscope}%
\pgfpathrectangle{\pgfqpoint{0.472301in}{0.352669in}}{\pgfqpoint{3.351934in}{3.351934in}}%
\pgfusepath{clip}%
\pgfsetbuttcap%
\pgfsetroundjoin%
\definecolor{currentfill}{rgb}{0.619608,0.188235,0.172549}%
\pgfsetfillcolor{currentfill}%
\pgfsetlinewidth{1.003750pt}%
\definecolor{currentstroke}{rgb}{0.619608,0.188235,0.172549}%
\pgfsetstrokecolor{currentstroke}%
\pgfsetdash{}{0pt}%
\pgfsys@defobject{currentmarker}{\pgfqpoint{-0.017361in}{-0.017361in}}{\pgfqpoint{0.017361in}{0.017361in}}{%
\pgfpathmoveto{\pgfqpoint{0.000000in}{-0.017361in}}%
\pgfpathcurveto{\pgfqpoint{0.004604in}{-0.017361in}}{\pgfqpoint{0.009020in}{-0.015532in}}{\pgfqpoint{0.012276in}{-0.012276in}}%
\pgfpathcurveto{\pgfqpoint{0.015532in}{-0.009020in}}{\pgfqpoint{0.017361in}{-0.004604in}}{\pgfqpoint{0.017361in}{0.000000in}}%
\pgfpathcurveto{\pgfqpoint{0.017361in}{0.004604in}}{\pgfqpoint{0.015532in}{0.009020in}}{\pgfqpoint{0.012276in}{0.012276in}}%
\pgfpathcurveto{\pgfqpoint{0.009020in}{0.015532in}}{\pgfqpoint{0.004604in}{0.017361in}}{\pgfqpoint{0.000000in}{0.017361in}}%
\pgfpathcurveto{\pgfqpoint{-0.004604in}{0.017361in}}{\pgfqpoint{-0.009020in}{0.015532in}}{\pgfqpoint{-0.012276in}{0.012276in}}%
\pgfpathcurveto{\pgfqpoint{-0.015532in}{0.009020in}}{\pgfqpoint{-0.017361in}{0.004604in}}{\pgfqpoint{-0.017361in}{0.000000in}}%
\pgfpathcurveto{\pgfqpoint{-0.017361in}{-0.004604in}}{\pgfqpoint{-0.015532in}{-0.009020in}}{\pgfqpoint{-0.012276in}{-0.012276in}}%
\pgfpathcurveto{\pgfqpoint{-0.009020in}{-0.015532in}}{\pgfqpoint{-0.004604in}{-0.017361in}}{\pgfqpoint{0.000000in}{-0.017361in}}%
\pgfpathlineto{\pgfqpoint{0.000000in}{-0.017361in}}%
\pgfpathclose%
\pgfusepath{stroke,fill}%
}%
\begin{pgfscope}%
\pgfsys@transformshift{3.671874in}{2.028636in}%
\pgfsys@useobject{currentmarker}{}%
\end{pgfscope}%
\end{pgfscope}%
\begin{pgfscope}%
\pgfsetfillopacity{0.700000}%
\pgfsetstrokeopacity{0.700000}%
\definecolor{textcolor}{rgb}{0.431373,0.423529,0.462745}%
\pgfsetstrokecolor{textcolor}%
\pgfsetfillcolor{textcolor}%
\pgftext[x=2.224448in,y=1.815331in,,base]{\color{textcolor}{\sffamily\fontsize{6.000000}{7.200000}\selectfont\catcode`\^=\active\def^{\ifmmode\sp\else\^{}\fi}\catcode`\%=\active\def%{\%}$2M$}}%
\end{pgfscope}%
\begin{pgfscope}%
\pgfsetfillopacity{0.700000}%
\pgfsetstrokeopacity{0.700000}%
\definecolor{textcolor}{rgb}{0.431373,0.423529,0.462745}%
\pgfsetstrokecolor{textcolor}%
\pgfsetfillcolor{textcolor}%
\pgftext[x=2.452989in,y=1.510610in,,base]{\color{textcolor}{\sffamily\fontsize{6.000000}{7.200000}\selectfont\catcode`\^=\active\def^{\ifmmode\sp\else\^{}\fi}\catcode`\%=\active\def%{\%}$6M$}}%
\end{pgfscope}%
\begin{pgfscope}%
\pgfsetrectcap%
\pgfsetroundjoin%
\pgfsetlinewidth{0.803000pt}%
\definecolor{currentstroke}{rgb}{0.768627,0.556863,0.211765}%
\pgfsetstrokecolor{currentstroke}%
\pgfsetdash{}{0pt}%
\pgfpathmoveto{\pgfqpoint{2.624360in}{3.595388in}}%
\pgfpathlineto{\pgfqpoint{2.714638in}{3.595388in}}%
\pgfpathlineto{\pgfqpoint{2.804916in}{3.595388in}}%
\pgfusepath{stroke}%
\end{pgfscope}%
\begin{pgfscope}%
\definecolor{textcolor}{rgb}{0.000000,0.000000,0.000000}%
\pgfsetstrokecolor{textcolor}%
\pgfsetfillcolor{textcolor}%
\pgftext[x=2.877138in,y=3.563791in,left,base]{\color{textcolor}{\sffamily\fontsize{6.500000}{7.800000}\selectfont\catcode`\^=\active\def^{\ifmmode\sp\else\^{}\fi}\catcode`\%=\active\def%{\%}Bound (precessing)}}%
\end{pgfscope}%
\begin{pgfscope}%
\pgfsetrectcap%
\pgfsetroundjoin%
\pgfsetlinewidth{0.803000pt}%
\definecolor{currentstroke}{rgb}{0.250980,0.400000,0.600000}%
\pgfsetstrokecolor{currentstroke}%
\pgfsetdash{}{0pt}%
\pgfpathmoveto{\pgfqpoint{2.624360in}{3.462881in}}%
\pgfpathlineto{\pgfqpoint{2.714638in}{3.462881in}}%
\pgfpathlineto{\pgfqpoint{2.804916in}{3.462881in}}%
\pgfusepath{stroke}%
\end{pgfscope}%
\begin{pgfscope}%
\definecolor{textcolor}{rgb}{0.000000,0.000000,0.000000}%
\pgfsetstrokecolor{textcolor}%
\pgfsetfillcolor{textcolor}%
\pgftext[x=2.877138in,y=3.431284in,left,base]{\color{textcolor}{\sffamily\fontsize{6.500000}{7.800000}\selectfont\catcode`\^=\active\def^{\ifmmode\sp\else\^{}\fi}\catcode`\%=\active\def%{\%}Marginal}}%
\end{pgfscope}%
\begin{pgfscope}%
\pgfsetrectcap%
\pgfsetroundjoin%
\pgfsetlinewidth{0.803000pt}%
\definecolor{currentstroke}{rgb}{0.619608,0.188235,0.172549}%
\pgfsetstrokecolor{currentstroke}%
\pgfsetdash{}{0pt}%
\pgfpathmoveto{\pgfqpoint{2.624360in}{3.330373in}}%
\pgfpathlineto{\pgfqpoint{2.714638in}{3.330373in}}%
\pgfpathlineto{\pgfqpoint{2.804916in}{3.330373in}}%
\pgfusepath{stroke}%
\end{pgfscope}%
\begin{pgfscope}%
\definecolor{textcolor}{rgb}{0.000000,0.000000,0.000000}%
\pgfsetstrokecolor{textcolor}%
\pgfsetfillcolor{textcolor}%
\pgftext[x=2.877138in,y=3.298776in,left,base]{\color{textcolor}{\sffamily\fontsize{6.500000}{7.800000}\selectfont\catcode`\^=\active\def^{\ifmmode\sp\else\^{}\fi}\catcode`\%=\active\def%{\%}Plunge}}%
\end{pgfscope}%
\end{pgfpicture}%
\makeatother%
\endgroup%

  \caption[Representative timelike orbits in Schwarzschild]{%
    Representative timelike orbits in the Schwarzschild geometry
    ($M = 1$), plotted as Cartesian projections of the equatorial
    plane.  Gold: a bound orbit with $L = 4M$, showing the
    characteristic relativistic perihelion advance --- the rosette
    pattern arises from the $-2ML^2/r^3$ correction in
    $V_{\text{eff}}$.  Blue: a marginally bound orbit
    ($E^2 = V_{\text{eff}}^{\text{max}}$) that asymptotically
    spirals toward the unstable circular orbit at $r = 4M$.
    Crimson: a plunge orbit ($L = 3.7M$, $E = 0.97$) whose energy
    exceeds the potential barrier, falling through the horizon
    without a turning point.  Dashed and dotted circles mark the
    event horizon ($r = 2M$) and ISCO ($r = 6M$).}
  \label{fig:timelike-orbits}
\end{figure}

\paragraph{Accretion disc physics.}

Circular timelike geodesics form the backbone of thin accretion disc
models.  In the Novikov--Thorne picture
\cite{Novikov1973,Page1974}, gas follows quasi-circular orbits,
slowly spiralling inward through viscous torques while radiating
binding energy at each radius.  At the ISCO, the gas reaches the
boundary of stability and plunges toward the horizon on a nearly
geodesic trajectory, radiating negligibly during the plunge.  The
radiative efficiency $\eta = 1 - E_{\text{ISCO}}$
(\cref{eq:kerr-efficiency}) connects the geodesic theory to
observable luminosities: $5.7\%$ for Schwarzschild
(\cref{eq:isco-efficiency}), up to $42.3\%$ for extremal Kerr
(\cref{eq:kerr-efficiency-max}).

The ray tracer does not integrate timelike geodesics directly ---
it traces photons, not particles.  But timelike orbit results enter
indirectly through the disc model.  The ISCO sets the inner edge of
the accretion disc in the termination configuration
(\cref{sec:schw-orbits}), and the circular orbit energy and angular
momentum at each radius determine the local four-velocity of the
emitting gas.  This four-velocity produces a Doppler shift that,
combined with the gravitational redshift, creates the characteristic
asymmetric brightness pattern visible in ray-traced images of
accretion discs: the approaching side is blue-shifted and
brightened, the receding side is red-shifted and dimmed.
\xrefnote{The frequency-shift calculation combining Doppler and
gravitational effects is developed in the camera and rendering
chapters.}
\warnnote{The Novikov--Thorne model assumes zero stress at the ISCO
--- matter in the plunge region exerts no torque on the disc above.
Magnetohydrodynamic (MHD) simulations --- which model the coupled
evolution of plasma and magnetic fields (Part~V) --- show that
magnetic stresses can produce non-negligible torques across the
ISCO, modestly increasing the efficiency and altering the emission
profile near the inner edge \cite{Noble2010}.}
